\documentclass[a4paper]{scrartcl}

\usepackage[
	fancytheorems,
	fancyproofs,
	noindent
]{adam}

\usepackage{tikz}

\newcommand{\A}{\mathbb{A}}
\newcommand{\V}{\mathbb{V}}
\newcommand{\mm}{\mathfrak{m}}
\newcommand{\pp}{\mathfrak{p}}
\newcommand{\Oh}{\mathcal{O}}
\newcommand{\dee}{\mathrm{d}}
\newcommand{\dto}{\dashrightarrow}
\newcommand{\Frac}{\operatorname{Frac}}
\DeclareMathOperator{\Mor}{Mor}
\DeclareMathOperator{\Div}{Div}
\DeclareMathOperator{\Divv}{div}
\DeclareMathOperator{\Prin}{Prin}
\DeclareMathOperator{\Pic}{Pic}
\DeclareMathOperator{\Cl}{Cl}
\DeclareMathOperator{\rk}{rank}
\DeclareMathOperator{\trdeg}{tr.deg}
\DeclareMathOperator{\Bir}{Bir}
\DeclareMathOperator{\dom}{dom}

\title{Algebraic Geometry}
\author{Adam Kelly (\texttt{ak2316@cam.ac.uk})}
\date{\today}

\allowdisplaybreaks

\begin{document}

\maketitle

Algebraic geometry is the study of the shapes cut out by polynomial equations. Write down a system of polynomials, look at the set of points where they all vanish, and you have a geometric object; write down the ring of polynomial functions on that object, and you have an algebraic one. The whole subject is an extended meditation on how much each of these two things knows about the other. It turns out that they know almost everything about each other, and the translations back and forth are what give algebraic geometry its distinctive flavour -- a topological question about a solution set becomes a question about prime ideals, and a question about ideals becomes a question about geometry.

We will spend the first half of the course setting up this dictionary carefully, first for affine varieties and then for projective ones, where the geometry is much better behaved. In the second half we specialise to curves, where the theory becomes remarkably complete: we will be able to attach an integer -- the genus -- to every smooth projective curve, and the Riemann--Roch theorem will let us compute with it.

This article constitutes my notes for the `Algebraic Geometry' course, held in Lent 2021 at Cambridge. These notes are \emph{not a transcription of the lectures}, and differ significantly in quite a few areas. Still, all lectured material should be covered.

\tableofcontents

\section{Introduction}

Let us begin with an example, since the whole subject grows out of examples of this kind. Consider the polynomial
$$f(X,Y) = Y^2 - X^3 + X \in \Z[X,Y],$$
and ask for its solution set. The answer depends rather dramatically on where we look for solutions.

If we look for real solutions, we get a subset of $\R^2$ that we can draw. Since $Y^2 = X(X-1)(X+1)$, a real point requires $X(X^2-1) \geq 0$, which happens exactly when $-1 \leq X \leq 0$ or $X \geq 1$. So the real locus has two pieces: a compact oval sitting over $[-1,0]$, and an unbounded branch over $[1,\infty)$.

\begin{center}
\begin{tikzpicture}[x=1.5cm,y=1.1cm]
	\draw[ultra thin, color=lightgray] (-1.5,-2.4) grid[step=0.5] (2.2,2.4);
	\draw[<->] (-1.6,0) -- (2.3,0);
	\draw[<->] (0,-2.5) -- (0,2.5);
	\draw[thick, color=blue, smooth] plot coordinates {(0.00,0.00) (-0.00,0.06) (-0.02,0.13) (-0.04,0.19) (-0.06,0.25) (-0.10,0.31) (-0.14,0.36) (-0.18,0.42) (-0.23,0.47) (-0.29,0.51) (-0.35,0.55) (-0.41,0.58) (-0.47,0.60) (-0.53,0.62) (-0.59,0.62) (-0.65,0.61) (-0.71,0.59) (-0.77,0.56) (-0.82,0.52) (-0.86,0.47) (-0.90,0.41) (-0.94,0.34) (-0.96,0.26) (-0.98,0.18) (-1.00,0.09) (-1.00,0.00) (-1.00,-0.09) (-0.98,-0.18) (-0.96,-0.26) (-0.94,-0.34) (-0.90,-0.41) (-0.86,-0.47) (-0.82,-0.52) (-0.77,-0.56) (-0.71,-0.59) (-0.65,-0.61) (-0.59,-0.62) (-0.53,-0.62) (-0.47,-0.60) (-0.41,-0.58) (-0.35,-0.55) (-0.29,-0.51) (-0.23,-0.47) (-0.18,-0.42) (-0.14,-0.36) (-0.10,-0.31) (-0.06,-0.25) (-0.04,-0.19) (-0.02,-0.13) (-0.00,-0.06) (0.00,0.00)};
	\draw[thick, color=blue, smooth] plot coordinates {(1.90,2.23) (1.82,2.06) (1.75,1.90) (1.68,1.75) (1.61,1.61) (1.55,1.48) (1.49,1.35) (1.44,1.23) (1.38,1.12) (1.33,1.02) (1.29,0.92) (1.24,0.83) (1.21,0.74) (1.17,0.66) (1.14,0.58) (1.11,0.50) (1.08,0.43) (1.06,0.37) (1.04,0.30) (1.03,0.24) (1.02,0.18) (1.01,0.12) (1.00,0.06) (1.00,0.00) (1.00,-0.06) (1.01,-0.12) (1.02,-0.18) (1.03,-0.24) (1.04,-0.30) (1.06,-0.37) (1.08,-0.43) (1.11,-0.50) (1.14,-0.58) (1.17,-0.66) (1.21,-0.74) (1.24,-0.83) (1.29,-0.92) (1.33,-1.02) (1.38,-1.12) (1.44,-1.23) (1.49,-1.35) (1.55,-1.48) (1.61,-1.61) (1.68,-1.75) (1.75,-1.90) (1.82,-2.06) (1.90,-2.23)};
	\node[below left] at (-1,0) {$-1$};
	\node[below right] at (1,0) {$1$};
\end{tikzpicture}
\end{center}

This picture is already suggestive. A `generic' line meets the curve in three points, which is no accident: substituting $Y = mX+c$ into $f$ produces a cubic in $X$, so there should be three roots. That there are sometimes fewer real roots than three is precisely the sort of irritation that pushes us towards algebraically closed fields.

If instead we look for solutions in $\C^2$, we can no longer draw the picture -- the solution set is a real surface sitting inside a real four-dimensional space -- but the algebra becomes uniform. Every line really does meet the curve three times. In fact, once we compactify appropriately, the solution set turns out to be a torus, and this is the first hint of the deep relationship between algebraic geometry and topology that the second half of these notes will explore.

Finally, if we look for solutions over $\Q$ or over $\Z$, we are doing number theory, and the questions become extremely hard. We will not do that here, but it is worth knowing that this is the same picture viewed through a different lens.

\subsection{Conventions}

We will work throughout over the complex numbers. There are two properties of $\C$ that we care about:
\begin{enumerate}[label=(\roman*)]
	\item $\C$ is algebraically closed -- every non-constant polynomial in one variable has a root. This is what makes the dictionary between algebra and geometry work at all; without it, $\V(X^2+1) \subseteq \R^1$ is empty even though $X^2+1$ is a perfectly good non-unit.
	\item $\C$ has characteristic zero, which means we can differentiate polynomials without anything going wrong, and $\C$ is uncountable, which we will exploit once (in the proof of the Nullstellensatz).
\end{enumerate}
Almost everything we do works verbatim over an arbitrary algebraically closed field, and the reader who prefers to write $k$ everywhere is encouraged to do so. We will write $\C[\mathbf{X}] = \C[X_1,\dots,X_n]$ when the number of variables is clear from context.

Two facts from IB Groups, Rings and Modules will be used constantly, and it is worth recalling them now.

\begin{theorem}[Hilbert Basis Theorem]
	If $R$ is a Noetherian ring then so is $R[X]$. In particular $\C[X_1,\dots,X_n]$ is Noetherian: every ideal of $\C[\mathbf{X}]$ is finitely generated.
\end{theorem}

\begin{theorem}
	$\C[X_1,\dots,X_n]$ is a unique factorisation domain.
\end{theorem}

The first of these is what makes all of our geometric objects `finite' in nature, and the second is what lets us talk about the irreducible polynomial cutting out a hypersurface.

\section{Affine Varieties}

\subsection{Affine Space and Polynomial Functions}

We begin with the ambient space.

\begin{definition}[Affine Space]
	\vocab{Affine $n$-space} over $\C$ is the set $\A^n = \C^n$. Its elements are called \vocab{points}, and we typically write a point as $P = (a_1,\dots,a_n)$.
\end{definition}

Why the new notation for something we already have a name for? The point is one of emphasis. When we write $\C^n$ we are usually thinking of a vector space, with a distinguished origin and a linear structure. When we write $\A^n$ we are forgetting all of that and remembering only that $\A^n$ is the natural home of the polynomial ring: every $f \in \C[X_1,\dots,X_n]$ determines a function
$$f \colon \A^n \to \A^1, \qquad P \mapsto f(P),$$
and it is this ring of functions, rather than any linear structure, that we regard as the essential data.

\begin{definition}[Affine Subspace]
	An \vocab{affine subspace} of $\A^n$ is a subset of the form $v + U$, where $U \leq \C^n$ is a linear subspace and $v \in \C^n$.
\end{definition}

Note that a polynomial is determined by the function it induces on $\A^n$: if $f$ and $g$ agree at every point of $\C^n$ then $f = g$ as polynomials. (This uses that $\C$ is infinite; over $\F_2$ the polynomials $X$ and $X^2$ induce the same function on $\A^1$.) So we may freely conflate polynomials with the functions they define.

\subsection{Vanishing Loci}

\begin{definition}[Vanishing Locus]
	Let $S \subseteq \C[X_1,\dots,X_n]$ be any subset. The \vocab{vanishing locus} of $S$ is
	$$\V(S) = \{P \in \A^n : f(P) = 0 \text{ for all } f \in S\} \subseteq \A^n.$$
	A subset of $\A^n$ of the form $\V(S)$ is called an \vocab{affine (algebraic) variety}.
\end{definition}

A warning about terminology: some authors (and some textbooks you may consult) reserve the word `variety' for what we will call an \emph{irreducible} variety, and use `algebraic set' for the general notion. We will not do that, but you should be alert to it.

Let us see some examples straight away.

\begin{example}
	\begin{enumerate}[label=(\roman*)]
		\item Take $n = 1$. A single polynomial $f \in \C[X]$ has finitely many roots, so $\V(f)$ is a finite set; and conversely, any finite set $\{a_1,\dots,a_m\} \subseteq \A^1$ is $\V(f)$ for $f = \prod_i (X - a_i)$. Also $\V(0) = \A^1$. So the affine varieties in $\A^1$ are exactly the finite sets together with $\A^1$ itself.
		\item A \vocab{hypersurface} in $\A^n$ is a variety of the form $\V(f)$ for a single non-constant $f$. For $n = 2$ these are the \vocab{plane curves}.
		\item Not all varieties are naturally presented by equations. The \vocab{twisted cubic} is
		$$C = \{(t,t^2,t^3) : t \in \C\} \subseteq \A^3,$$
		presented parametrically. It is nonetheless a variety: $C = \V(X_2 - X_1^2,\ X_3 - X_1^3)$. Indeed if $x_2 = x_1^2$ and $x_3 = x_1^3$ then the point is $(t,t^2,t^3)$ for $t = x_1$.
		\item Every affine subspace is a variety, being the vanishing locus of a collection of degree one polynomials. In particular a point $(a_1,\dots,a_n)$ is the variety $\V(X_1-a_1,\dots,X_n-a_n)$.
	\end{enumerate}
\end{example}

Passing from a set $S$ of polynomials to $\V(S)$ loses a lot of information, and the first thing to do is to understand exactly what it remembers. The following is the basic finiteness statement of the subject.

\begin{theorem}\label{thm:Vfinite}
	Let $S \subseteq \C[\mathbf{X}]$ and let $I = (S)$ be the ideal it generates. Then
	\begin{enumerate}[label=(\roman*)]
		\item $\V(S) = \V(I)$;
		\item there is a \emph{finite} subset $\{f_1,\dots,f_r\} \subseteq S$ with $\V(S) = \V(f_1,\dots,f_r)$.
	\end{enumerate}
\end{theorem}

\begin{proof}
	For (i), certainly $\V(I) \subseteq \V(S)$ since $S \subseteq I$. Conversely if $P \in \V(S)$ and $g = \sum_j h_j f_j \in I$ with $f_j \in S$, then $g(P) = \sum_j h_j(P) f_j(P) = 0$, so $P \in \V(I)$.

	For (ii), by the Hilbert Basis Theorem $I$ is finitely generated, say $I = (g_1,\dots,g_m)$. Each $g_i$ lies in the ideal generated by $S$, so is a finite combination $g_i = \sum_j h_{ij} f_{ij}$ with $f_{ij} \in S$. Let $T$ be the (finite) set of all the $f_{ij}$ occurring. Then $I = (g_1,\dots,g_m) \subseteq (T) \subseteq I$, so $(T) = I$ and by (i) we get $\V(S) = \V(I) = \V(T)$.
\end{proof}

In other words, \emph{every} system of polynomial equations, however large, is equivalent to a finite subsystem. This is the geometric content of the Hilbert Basis Theorem, and it is not at all obvious from the geometric side.

The operation $\V$ reverses inclusions and interacts with the lattice operations on ideals in a way we now record.

\begin{proposition}\label{prop:Vprops}
	Let $S, T \subseteq \C[\mathbf{X}]$ be subsets and let $(I_j)_{j \in J}$ be a family of ideals of $\C[\mathbf{X}]$. Then
	\begin{enumerate}[label=(\roman*)]
		\item if $S \subseteq T$ then $\V(T) \subseteq \V(S)$;
		\item $\V(0) = \A^n$ and $\V(\C[\mathbf{X}]) = \V(1) = \varnothing$;
		\item $\bigcap_{j \in J} \V(I_j) = \V\big(\sum_{j} I_j\big)$;
		\item $\V(I) \cup \V(J) = \V(I \cap J) = \V(IJ)$.
	\end{enumerate}
\end{proposition}

\begin{proof}
	Parts (i) and (ii) are immediate from the definitions.

	For (iii), a point lies in every $\V(I_j)$ if and only if it is killed by every element of every $I_j$, if and only if it is killed by $\bigcup_j I_j$; and $\V$ of a set equals $\V$ of the ideal it generates, which is $\sum_j I_j$.

	For (iv), we have $IJ \subseteq I \cap J \subseteq I$ and likewise for $J$, so by (i),
	$$\V(I) \cup \V(J) \subseteq \V(I\cap J) \subseteq \V(IJ).$$
	It remains to show $\V(IJ) \subseteq \V(I) \cup \V(J)$. Suppose $P \in \V(IJ)$ but $P \notin \V(I)$, so there is $g \in I$ with $g(P) \neq 0$. For any $f \in J$ we have $fg \in IJ$, so $f(P)g(P) = 0$, whence $f(P) = 0$. Thus $P \in \V(J)$.
\end{proof}

\subsection{The Zariski Topology}

Proposition~\ref{prop:Vprops} says exactly that the affine varieties are the closed sets of a topology: they contain $\varnothing$ and the whole space, and are closed under arbitrary intersections and finite unions.

\begin{definition}[Zariski Topology]
	The \vocab{Zariski topology} on $\A^n$ is the topology whose closed sets are precisely the affine varieties. If $V \subseteq \A^n$ is a variety, the Zariski topology on $V$ is the subspace topology.
\end{definition}

\begin{definition}[Euclidean Topology]
	The \vocab{Euclidean} or \vocab{analytic} topology on $\A^n = \C^n$ is the usual metric topology. Again, a variety inherits a subspace topology.
\end{definition}

Every Zariski closed set is Euclidean closed, since polynomials are continuous; so the Zariski topology is \emph{coarser} than the Euclidean one. It is in fact very much coarser, and the reader should be warned that it does not behave like the topologies met in IB Analysis and Topology.

\begin{proposition}
	The Zariski topology on $\A^1$ is the cofinite topology. It is compact and irreducible, but not Hausdorff. The Euclidean topology on $\A^1$ is Hausdorff but not compact.
\end{proposition}

\begin{proof}
	We saw that the varieties in $\A^1$ are the finite sets and $\A^1$ itself, which is the definition of the cofinite topology. Any two nonempty open sets are cofinite, hence meet (as $\C$ is infinite), so the space is not Hausdorff. Compactness is immediate: given an open cover, one member $U$ is nonempty, and its complement is finite, so finitely many further members finish the job.
\end{proof}

\begin{remark}
	The Zariski topology on $\A^2$ is \emph{not} the product topology of the Zariski topologies on $\A^1 \times \A^1$. In the product topology the closed sets are finite unions of sets of the form $A \times B$ with $A, B$ closed; the diagonal $\V(X_1 - X_2)$ is Zariski closed in $\A^2$ but is not of this form.
\end{remark}

Because the Zariski topology is so coarse, nonempty open sets are enormous. We will use this constantly: a Zariski open set is `almost all' of the space, and a statement holding on a nonempty open set is a statement holding `generically'.

\subsection{Irreducibility}

A cofinite topology has no interesting disconnections, but varieties can still fall apart in an algebraic sense, as the following notion makes precise.

\begin{definition}[Irreducible]
	A nonempty variety $V$ is \vocab{irreducible} if whenever $V = V_1 \cup V_2$ with $V_1, V_2$ closed subsets of $V$, we have $V = V_1$ or $V = V_2$. Otherwise $V$ is \vocab{reducible}.
\end{definition}

\begin{example}
	The variety $V = \V(XY) \subseteq \A^2$ is reducible: it is the union of the two coordinate axes $\V(X)$ and $\V(Y)$, neither of which is all of $V$.
	\begin{center}
	\begin{tikzpicture}[x=1cm,y=1cm]
		\draw[ultra thin, color=lightgray] (-1.8,-1.8) grid[step=0.6] (1.8,1.8);
		\draw[<->] (-1.9,0) -- (1.9,0);
		\draw[<->] (0,-1.9) -- (0,1.9);
		\draw[very thick, color=blue] (-1.7,0) -- (1.7,0);
		\draw[very thick, color=blue] (0,-1.7) -- (0,1.7);
		\node[above right, color=blue] at (1.1,0) {$\V(Y)$};
		\node[right, color=blue] at (0.05,1.5) {$\V(X)$};
	\end{tikzpicture}
	\end{center}
	By contrast, we shall see shortly that each axis is irreducible.
\end{example}

Note that irreducibility is a genuinely topological condition -- it says that $V$ is not the union of two proper closed subsets -- but it is a condition no Hausdorff space with more than one point satisfies. It is a condition suited to the Zariski topology.

\begin{proposition}\label{prop:decomp}
	Every affine variety is a finite union of irreducible varieties.
\end{proposition}

\begin{proof}
	Suppose not, and let $V$ be a variety that is not a finite union of irreducibles. In particular $V$ is not itself irreducible, so $V = V_1 \cup W_1$ with $V_1, W_1$ proper closed subsets. If both $V_1$ and $W_1$ were finite unions of irreducibles then so would $V$ be; so (relabelling) $V_1$ is not a finite union of irreducibles. Repeating, we obtain a strictly decreasing chain
	$$V = V_0 \supsetneq V_1 \supsetneq V_2 \supsetneq \cdots$$
	of varieties. Let $I_j = I(V_j)$ (defined below in Section~\ref{subsec:ideals}); the containments reverse, giving a strictly increasing chain of ideals $I_0 \subsetneq I_1 \subsetneq \cdots$. But $\C[\mathbf{X}]$ is Noetherian, so no such chain exists, and we have our contradiction.
\end{proof}

The reader will notice that we have used the ideal $I(V)$ of a variety before defining it; the reader who objects may instead run the argument with $\bigcup_j I_j$, which is an ideal and hence finitely generated, so equals $\sum_{j \leq N} I_j$ for some $N$, forcing the chain to stabilise. Either way, the essential point is the Noetherian property.

\begin{definition}[Minimal Decomposition]
	A \vocab{minimal decomposition} of a variety $V$ is an expression $V = V_1 \cup \cdots \cup V_r$ with each $V_i$ irreducible and $V_i \not\subseteq V_j$ for $i \neq j$.
\end{definition}

Given any decomposition into irreducibles, we can throw away any piece contained in another and obtain a minimal one. What is not obvious is that the result does not depend on the choices made.

\begin{proposition}
	Minimal decompositions are unique up to reordering.
\end{proposition}

\begin{proof}
	Suppose $V = \bigcup_{i=1}^r V_i = \bigcup_{j=1}^s W_j$ are two minimal decompositions. Fix $i$. Then
	$$V_i = V_i \cap V = \bigcup_{j=1}^s (V_i \cap W_j),$$
	a finite union of closed subsets of $V_i$. By irreducibility, $V_i = V_i \cap W_j$ for some $j$, i.e.\ $V_i \subseteq W_j$. Symmetrically $W_j \subseteq V_k$ for some $k$, so $V_i \subseteq V_k$, and minimality forces $i = k$ and hence $V_i = W_j$. So each $V_i$ appears among the $W_j$ and vice versa.
\end{proof}

\begin{definition}[Irreducible Components]
	The irreducible varieties appearing in the (unique) minimal decomposition of $V$ are the \vocab{irreducible components} of $V$.
\end{definition}

\subsection{The Ideal of a Variety}\label{subsec:ideals}

We now go the other way, from geometry back to algebra.

\begin{definition}
	Let $X \subseteq \A^n$ be any subset. The \vocab{ideal of $X$} is
	$$I(X) = \{f \in \C[\mathbf{X}] : f(P) = 0 \text{ for all } P \in X\}.$$
\end{definition}

This really is an ideal: it is closed under addition, and if $f$ vanishes on $X$ then so does $gf$ for any $g$. It is the largest set of polynomials with vanishing locus containing $X$.

\begin{proposition}\label{prop:IVbasic}
	Let $V, W \subseteq \A^n$ be affine varieties and $S \subseteq \C[\mathbf{X}]$.
	\begin{enumerate}[label=(\roman*)]
		\item If $V = \V(S)$ then $S \subseteq I(V)$.
		\item $V = \V(I(V))$.
		\item $V \subseteq W$ if and only if $I(W) \subseteq I(V)$. In particular $V = W$ if and only if $I(V) = I(W)$.
		\item $I(V \cup W) = I(V) \cap I(W)$.
	\end{enumerate}
\end{proposition}

\begin{proof}
	(i) is immediate. For (ii), $S \subseteq I(V)$ gives $\V(I(V)) \subseteq \V(S) = V$, and $V \subseteq \V(I(V))$ holds by definition of $I(V)$.

	For (iii), if $V \subseteq W$ then any polynomial vanishing on $W$ vanishes on $V$. Conversely suppose $V \not\subseteq W$ and pick $P \in V \setminus W$. As $W = \V(I(W))$ and $P \notin W$, there is $f \in I(W)$ with $f(P) \neq 0$; then $f \notin I(V)$.

	(iv) is immediate from the definitions.
\end{proof}

So $I$ is an injection from varieties in $\A^n$ to ideals of $\C[\mathbf{X}]$, with $\V$ as a left inverse. It is certainly not surjective: the ideals $(X)$ and $(X^2)$ in $\C[X]$ have the same vanishing locus $\{0\}$, and $I(\{0\}) = (X)$. So the composite $I \circ \V$ is not the identity on ideals, and understanding exactly what it does is the content of the Nullstellensatz.

First, however, let us record the algebraic characterisation of irreducibility.

\begin{proposition}\label{prop:irrprime}
	A variety $V$ is irreducible if and only if $I(V)$ is a prime ideal.
\end{proposition}

Recall that $I(V)$ prime means: whenever $f_1 f_2 \in I(V)$, either $f_1 \in I(V)$ or $f_2 \in I(V)$. Geometrically, the failure of primality is the existence of two functions, neither vanishing identically on $V$, whose product does; and such a pair splits $V$ in two.

\begin{proof}
	Suppose $V$ is reducible, $V = V_1 \cup V_2$ with $V_i \subsetneq V$ closed. Then $I(V_i) \supsetneq I(V)$ strictly, and $V_1 \not\subseteq V_2$, $V_2 \not\subseteq V_1$ (else the decomposition would be trivial), so $I(V_1) \not\subseteq I(V_2)$ and vice versa. Choose $f_1 \in I(V_1)\setminus I(V_2)$ and $f_2 \in I(V_2)\setminus I(V_1)$. Neither lies in $I(V) = I(V_1) \cap I(V_2)$, but $f_1f_2$ vanishes on $V_1 \cup V_2 = V$. So $I(V)$ is not prime.

	Conversely suppose $I(V)$ is not prime, with $f_1f_2 \in I(V)$ but $f_i \notin I(V)$. Set $V_i = V \cap \V(f_i)$, a closed subset of $V$. Since $f_i \notin I(V)$, $f_i$ does not vanish identically on $V$, so $V_i \subsetneq V$. But for any $P \in V$ we have $f_1(P)f_2(P) = 0$, so $P \in V_1 \cup V_2$. Hence $V = V_1 \cup V_2$ is reducible.
\end{proof}

\begin{example}
	$\V(X) \subseteq \A^2$ is irreducible, since $I(\V(X)) = (X)$ and $\C[X,Y]/(X) \cong \C[Y]$ is an integral domain, so $(X)$ is prime. Hence $\V(XY) = \V(X) \cup \V(Y)$ is the minimal decomposition of $\V(XY)$ into irreducible components.
\end{example}

\begin{example}
	More generally, if $f \in \C[\mathbf{X}]$ is irreducible then $(f)$ is prime (as $\C[\mathbf{X}]$ is a UFD), so the hypersurface $\V(f)$ is irreducible. If $f = f_1^{a_1}\cdots f_r^{a_r}$ is the factorisation of a general $f$ into irreducibles, then $\V(f) = \V(f_1) \cup \cdots \cup \V(f_r)$ is the decomposition into irreducible components.
\end{example}

\subsection{Hilbert's Nullstellensatz}

We now come to the theorem which makes the whole dictionary work. The name means `theorem of the zeros', and it comes in two forms.

\begin{theorem}[Weak Nullstellensatz]\label{thm:weaknss}
	\begin{enumerate}[label=(\roman*)]
		\item Every maximal ideal of $\C[X_1,\dots,X_n]$ has the form $\mm_P = (X_1-a_1,\dots,X_n-a_n)$ for some point $P = (a_1,\dots,a_n) \in \A^n$.
		\item If $I \trianglelefteq \C[\mathbf{X}]$ is a proper ideal then $\V(I) \neq \varnothing$.
	\end{enumerate}
\end{theorem}

Part (ii) is the statement that a system of polynomial equations with no `algebraic obstruction' -- that is, such that $1$ is not a combination of the equations -- really does have a solution. It is the many-variable generalisation of the fundamental theorem of algebra.

The proof we give uses that $\C$ is uncountable, and so is special to $\C$; the theorem is true over any algebraically closed field, but the general proof (via the Noether normalisation lemma, or via Zariski's lemma) is more work.

\begin{proof}
	Each ideal $\mm_P$ is maximal: the evaluation homomorphism $\C[\mathbf{X}] \to \C$, $f \mapsto f(P)$, is surjective with kernel exactly $\mm_P$, so $\C[\mathbf{X}]/\mm_P \cong \C$ is a field.

	Conversely let $\mm$ be a maximal ideal and set $K = \C[\mathbf{X}]/\mm$, a field containing $\C$. Write $a_i$ for the image of $X_i$ in $K$, so that $K = \C[a_1,\dots,a_n]$. If every $a_i$ lies in $\C \subseteq K$ then $X_i - a_i \in \mm$ for each $i$, so $\mm \supseteq \mm_P$ for $P = (a_1,\dots,a_n)$, and maximality of $\mm_P$ gives $\mm = \mm_P$ as required.

	So suppose for contradiction that $K \supsetneq \C$, and pick $t \in K \setminus \C$. Since $\C$ is algebraically closed, $t$ cannot be algebraic over $\C$, so $t$ is transcendental and the elements $\{(t-c)^{-1} : c \in \C\}$ make sense in $K$. We claim they are linearly independent over $\C$. Indeed, a relation $\sum_{i=1}^m \lambda_i (t-c_i)^{-1} = 0$ with distinct $c_i$, cleared of denominators, gives $\sum_i \lambda_i \prod_{j\neq i}(t - c_j) = 0$; substituting the transcendental $t$ into this polynomial identity and evaluating the polynomial at $c_i$ shows $\lambda_i \prod_{j \neq i}(c_i - c_j) = 0$, so $\lambda_i = 0$.

	Thus $K$ contains an uncountable linearly independent set over $\C$, so $\dim_\C K$ is uncountable. On the other hand, let $U_m \subseteq K$ be the $\C$-span of the monomials $a_1^{r_1}\cdots a_n^{r_n}$ with $\sum r_i \leq m$. Each $U_m$ is finite-dimensional and $K = \bigcup_{m \geq 0} U_m$, so $K$ has a countable spanning set and hence countable dimension. This is a contradiction.

	For (ii), a proper ideal $I$ is contained in some maximal ideal $\mm_P$ (using that $\C[\mathbf{X}]$ is Noetherian, or Zorn's lemma), and then $P \in \V(\mm_P) \subseteq \V(I)$.
\end{proof}

To state the strong form we need one more piece of algebra.

\begin{definition}[Radical]
	The \vocab{radical} of an ideal $I \trianglelefteq R$ is
	$$\sqrt{I} = \{f \in R : f^m \in I \text{ for some } m \geq 1\}.$$
	An ideal is \vocab{radical} if $I = \sqrt{I}$.
\end{definition}

That $\sqrt I$ is an ideal follows from the binomial theorem: if $f^a, g^b \in I$ then every term of $(f+g)^{a+b}$ contains $f^a$ or $g^b$. Note $\sqrt{\sqrt I} = \sqrt I$, and prime ideals are radical. Since $f$ and $f^m$ have the same zeros, $\V(I) = \V(\sqrt I)$ always.

\begin{theorem}[Strong Nullstellensatz]\label{thm:strongnss}
	Let $I \trianglelefteq \C[X_1,\dots,X_n]$ be an ideal. Then
	$$I(\V(I)) = \sqrt{I}.$$
\end{theorem}

\begin{proof}
	The inclusion $\sqrt I \subseteq I(\V(I))$ is clear: if $f^m \in I$ then $f^m$, hence $f$, vanishes on $\V(I)$.

	For the converse we use the \emph{Rabinowitsch trick}. Let $f \in I(\V(I))$, and let $g_1,\dots,g_r$ generate $I$. Introduce a new variable $T$ and consider the ideal
	$$J = (g_1,\dots,g_r,\ 1 - Tf) \trianglelefteq \C[X_1,\dots,X_n,T].$$
	I claim $\V(J) = \varnothing$ in $\A^{n+1}$. Indeed, if $(P,\tau) \in \V(J)$ then $P \in \V(I)$, so $f(P) = 0$, and then $1 - \tau f(P) = 1 \neq 0$, a contradiction.

	By the weak Nullstellensatz, $J = \C[\mathbf{X},T]$, so we may write
	$$1 = \sum_{i=1}^r h_i(\mathbf{X},T) g_i(\mathbf{X}) + h_0(\mathbf{X},T)(1 - Tf(\mathbf{X})).$$
	This is an identity in the polynomial ring, so we may substitute $T = 1/f$ and work in the field $\C(X_1,\dots,X_n)$ (note $f \neq 0$, else $f \in \sqrt I$ trivially). The last term vanishes, leaving
	$$1 = \sum_{i=1}^r h_i(\mathbf{X},1/f)\, g_i(\mathbf{X}).$$
	Multiplying by a sufficiently large power $f^N$ to clear all denominators gives $f^N = \sum_i \tilde{h}_i g_i$ with $\tilde h_i \in \C[\mathbf{X}]$, so $f^N \in I$ and $f \in \sqrt I$.
\end{proof}

\subsection{The Dictionary}

We can now say precisely what $I$ and $\V$ do.

\begin{corollary}\label{cor:dictionary}
	The maps $V \mapsto I(V)$ and $I \mapsto \V(I)$ are mutually inverse, inclusion-reversing bijections
	$$\{\text{affine varieties in } \A^n\} \longleftrightarrow \{\text{radical ideals of } \C[X_1,\dots,X_n]\}.$$
	Under this bijection, irreducible varieties correspond to prime ideals, and points correspond to maximal ideals.
\end{corollary}

\begin{proof}
	We know $\V(I(V)) = V$ for varieties $V$, and $I(\V(I)) = \sqrt I = I$ for radical $I$. Also $I(V)$ is always radical, since if $f^m$ vanishes on $V$ then so does $f$. The correspondence for prime ideals is Proposition~\ref{prop:irrprime}, and for maximal ideals it is Theorem~\ref{thm:weaknss}(i).
\end{proof}

This is the fundamental dictionary of the subject, and it is worth pausing over. Every geometric question about varieties in $\A^n$ can, in principle, be converted into an algebraic question about radical ideals of a polynomial ring. The rest of the course consists of doing exactly this, in increasingly refined ways.

\begin{example}
	Let us see how the correspondence handles the classic example. In $\C[X]$, the ideal $(X^2)$ is not radical, and $\sqrt{(X^2)} = (X)$. Correspondingly $\V(X^2) = \{0\} = \V(X)$: the variety cannot see the doubling. Later, when we come to intersection multiplicities and to schemes, we will want to remember the difference; for now we simply record that varieties forget it.
\end{example}

\begin{example}
	Consider $I = (X^2 - Y^2, X^2+Y^2 - 2) \trianglelefteq \C[X,Y]$. Adding and subtracting, $I$ contains $2X^2 - 2$ and $2Y^2-2$, so $I = (X^2-1, Y^2-1)$ and $\V(I) = \{(\pm 1, \pm 1)\}$, four points. The corresponding radical ideal is $I(\V(I)) = (X^2-1,Y^2-1)$, which is indeed radical -- the quotient is $\C^4$ by the Chinese remainder theorem, and a finite product of fields has no nilpotents.
\end{example}

\section{Regular Functions and Morphisms}

Having built our objects, we need the maps between them. As always in mathematics, the maps are more important than the objects, and here they will let us make precise the sense in which a variety `is' its ring of functions.

\subsection{Coordinate Rings}

Let $V \subseteq \A^n$ be a variety. Any polynomial $f \in \C[\mathbf{X}]$ restricts to a function $V \to \C$, but different polynomials may restrict to the same function. Precisely, $f$ and $g$ restrict to the same function on $V$ if and only if $f - g$ vanishes on $V$, i.e.\ $f - g \in I(V)$. So the ring of functions $V \to \C$ obtained by restricting polynomials is the quotient $\C[\mathbf{X}]/I(V)$.

\begin{definition}[Coordinate Ring]
	Let $V \subseteq \A^n$ be an affine variety. Its \vocab{coordinate ring}, or \vocab{ring of regular functions}, is
	$$\C[V] = \frac{\C[X_1,\dots,X_n]}{I(V)}.$$
	Elements of $\C[V]$ are called \vocab{regular functions} on $V$. We sometimes write $\Oh(V)$ for $\C[V]$.
\end{definition}

The coordinate ring is a finitely generated $\C$-algebra (generated by the images $x_i$ of the coordinate functions $X_i$), and it is \emph{reduced}: it has no nonzero nilpotents, because $I(V)$ is radical. Conversely, every reduced finitely generated $\C$-algebra arises this way, by writing it as $\C[\mathbf{X}]/I$ with $I$ radical.

\begin{corollary}
	A variety $V$ is irreducible if and only if $\C[V]$ is an integral domain.
\end{corollary}

\begin{proof}
	$\C[V]$ is a domain if and only if $I(V)$ is prime, which by Proposition~\ref{prop:irrprime} happens if and only if $V$ is irreducible.
\end{proof}

\begin{example}
	$\C[\A^n] = \C[X_1,\dots,X_n]$, since $I(\A^n) = 0$.
\end{example}

\begin{example}
	Let $V = \{0\} \subseteq \A^1$. We can present $V$ as $\V(X^2)$, but $I(V) = \sqrt{(X^2)} = (X)$, so
	$$\C[V] = \C[X]/(X) \cong \C.$$
	The presentation as a doubled point is invisible to the coordinate ring. This is exactly the information that the theory of schemes is designed to retain: the ring $\C[X]/(X^2)$ is a perfectly good ring, with a nonzero element squaring to zero, and it remembers the doubling. We will not pursue this.
\end{example}

\begin{example}
	Let $C = \V(Y - X^2) \subseteq \A^2$ be a parabola. Then $I(C) = (Y - X^2)$ (one checks that $Y - X^2$ is irreducible, so $(Y-X^2)$ is prime hence radical, and it certainly has the right vanishing locus). Hence
	$$\C[C] = \C[X,Y]/(Y-X^2) \cong \C[X],$$
	the isomorphism sending $Y \mapsto X^2$. So the parabola has the same coordinate ring as the affine line -- and we will see in a moment that this means they are isomorphic as varieties.
\end{example}

\begin{remark}
	One should be careful: $\C[V]$ on its own, as an abstract ring, does not remember the embedding $V \subseteq \A^n$. A surjection $\C[\mathbf{X}] \to \C[V]$ is the same thing as a presentation of $V$ as a subvariety of $\A^n$, and there are many such. The parabola example above is a case in point: $\C[X] \cong \C[X,Y]/(Y-X^2)$ presents the same ring both as $\A^1 \subseteq \A^1$ and as a curve in $\A^2$.
\end{remark}

\subsection{Morphisms}

\begin{definition}[Morphism]
	Let $V \subseteq \A^n$ and $W \subseteq \A^m$ be affine varieties. A \vocab{morphism} (or \vocab{regular map}) $\varphi \colon V \to W$ is a function such that there exist $f_1,\dots,f_m \in \C[V]$ with
	$$\varphi(P) = (f_1(P),\dots,f_m(P))$$
	for all $P \in V$. We write $\Mor(V,W)$ for the set of morphisms $V \to W$.
\end{definition}

In other words, a morphism is a map given in coordinates by polynomials. Note the requirement that the image actually lands in $W$, which is a genuine condition on the $f_i$.

\begin{example}
	\begin{enumerate}[label=(\roman*)]
		\item A morphism $V \to \A^1$ is precisely an element of $\C[V]$.
		\item Linear maps, translations, and more generally affine-linear maps $\A^n \to \A^m$ are morphisms.
		\item If $V \subseteq W$ are varieties in $\A^n$, the inclusion $V \hookrightarrow W$ is a morphism.
		\item The map $\A^1 \to \A^3$, $t \mapsto (t,t^2,t^3)$, is a morphism onto the twisted cubic.
	\end{enumerate}
\end{example}

\begin{proposition}
	Compositions of morphisms are morphisms, so affine varieties and morphisms form a category.
\end{proposition}

\begin{proof}
	A polynomial in polynomials is a polynomial.
\end{proof}

Morphisms are automatically continuous for the Zariski topology: if $Z = \V(g_1,\dots,g_r) \subseteq W$ then $\varphi^{-1}(Z) = \V(g_1 \circ \varphi, \dots, g_r \circ \varphi)$ is closed in $V$. The converse is very far from true -- for instance every bijection of $\A^1$ is a Zariski homeomorphism -- so being a morphism is much stronger than being continuous.

\subsection{Pullbacks}

The key construction is that a morphism induces a map on functions, going the other way.

\begin{definition}[Pullback]
	Let $\varphi \colon V \to W$ be a morphism and $g \in \C[W]$. The \vocab{pullback} of $g$ is
	$$\varphi^\star(g) = g \circ \varphi \colon V \to \C.$$
\end{definition}

If $g$ is (the class of) a polynomial in the coordinates on $\A^m$, then $g \circ \varphi$ is a polynomial in the coordinates on $\A^n$, so $\varphi^\star(g) \in \C[V]$. Thus $\varphi^\star \colon \C[W] \to \C[V]$, and it is clearly a ring homomorphism which is the identity on the constants $\C$.

\begin{definition}
	A ring homomorphism between $\C$-algebras which restricts to the identity on $\C$ is a \vocab{$\C$-algebra homomorphism}.
\end{definition}

\begin{theorem}\label{thm:morequiv}
	Let $V \subseteq \A^n$ and $W \subseteq \A^m$ be affine varieties. The map $\varphi \mapsto \varphi^\star$ is a bijection
	$$\Mor(V,W) \longleftrightarrow \{\C\text{-algebra homomorphisms } \C[W] \to \C[V]\}.$$
\end{theorem}

\begin{proof}
	Let $y_1,\dots,y_m \in \C[W]$ be the coordinate functions, i.e.\ the restrictions to $W$ of the standard coordinates on $\A^m$.

	\emph{Injectivity.} If $\varphi \colon V \to W$ is a morphism and $P \in V$, then
	$$\varphi(P) = \big(y_1(\varphi(P)),\dots,y_m(\varphi(P))\big) = \big(\varphi^\star(y_1)(P),\dots,\varphi^\star(y_m)(P)\big),$$
	so $\varphi$ is determined by $\varphi^\star(y_1),\dots,\varphi^\star(y_m)$, hence by $\varphi^\star$.

	\emph{Surjectivity.} Let $\lambda \colon \C[W] \to \C[V]$ be a $\C$-algebra homomorphism, and put $f_j = \lambda(y_j) \in \C[V]$. Define $\varphi = (f_1,\dots,f_m) \colon V \to \A^m$; this is given by polynomials, so we need only check its image lies in $W$. Let $P \in V$ and $g \in I(W)$. Since $\lambda$ is a $\C$-algebra homomorphism, it commutes with polynomial expressions, so
	$$g(f_1(P),\dots,f_m(P)) = \big(g(\lambda(y_1),\dots,\lambda(y_m))\big)(P) = \big(\lambda(g(y_1,\dots,y_m))\big)(P).$$
	But $g(y_1,\dots,y_m) = 0$ in $\C[W]$, because $g \in I(W)$. Hence $g(\varphi(P)) = 0$ for all $g \in I(W)$, so $\varphi(P) \in \V(I(W)) = W$. Finally $\varphi^\star(y_j) = f_j = \lambda(y_j)$, and since the $y_j$ generate $\C[W]$ as a $\C$-algebra, $\varphi^\star = \lambda$.
\end{proof}

\begin{definition}[Isomorphism]
	Varieties $V$ and $W$ are \vocab{isomorphic} if there are morphisms $\varphi \colon V \to W$ and $\psi \colon W \to V$ with $\psi\varphi = \mathrm{id}_V$ and $\varphi\psi = \mathrm{id}_W$.
\end{definition}

\begin{corollary}\label{cor:isocoord}
	$V \cong W$ as varieties if and only if $\C[V] \cong \C[W]$ as $\C$-algebras.
\end{corollary}

\begin{proof}
	Immediate from Theorem~\ref{thm:morequiv}, since the bijection is compatible with composition (it is contravariantly functorial): $(\psi\varphi)^\star = \varphi^\star\psi^\star$ and $\mathrm{id}^\star = \mathrm{id}$.
\end{proof}

This is the promised statement that a variety `is' its coordinate ring. All geometric information about an affine variety is encoded in a finitely generated reduced $\C$-algebra, and every such algebra arises. In the language of category theory, the functor $V \mapsto \C[V]$ is an anti-equivalence between the category of affine varieties and the category of finitely generated reduced $\C$-algebras.

\begin{example}
	The twisted cubic $C = \{(t,t^2,t^3)\}$ is isomorphic to $\A^1$. Indeed $I(C) = (X_2 - X_1^2, X_3 - X_1^3)$, and in $\C[C]$ we have $x_2 = x_1^2$, $x_3 = x_1^3$, so $\C[C] = \C[x_1] \cong \C[X]$.
\end{example}

\begin{example}
	Not every bijective morphism is an isomorphism. Consider $\varphi \colon \A^1 \to C = \V(Y^2 - X^3) \subseteq \A^2$ given by $t \mapsto (t^2,t^3)$. This is a bijection (the inverse sends $(x,y) \neq (0,0)$ to $y/x$, and $(0,0)$ to $0$). But $\varphi^\star \colon \C[C] \to \C[t]$ has image $\C[t^2,t^3] \subsetneq \C[t]$, since $t \notin \C[t^2,t^3]$. So $\varphi^\star$ is not surjective and $\varphi$ is not an isomorphism. Geometrically, the cuspidal cubic $C$ has a singular point at the origin which the smooth line $\A^1$ does not; we will make this precise in Section~\ref{sec:tangent}.
\end{example}

\begin{example}\label{ex:threelines}
	Let $V = \V(X_1X_2(X_1-X_2)) \subseteq \A^2$, a union of three concurrent lines in the plane, and let $W = \V(X_1X_2,\, X_2X_3,\, X_3X_1) \subseteq \A^3$, the union of the three coordinate axes. Both are unions of three lines meeting at a single point, and they are homeomorphic in the Euclidean topology. Nevertheless they are not isomorphic as varieties: the tangent space to $V$ at the origin is two-dimensional (it sits inside a plane) while that of $W$ is three-dimensional. We will be able to see this once we have defined tangent spaces.
\end{example}

\subsection{Rational Functions and Local Rings}

Regular functions are rather rigid. On $\A^1$ the only regular functions are the polynomials, and there are many natural functions -- $1/X$, say -- that we would like to allow, provided we keep track of where they are undefined. The right framework is that of rational functions.

Throughout this subsection, let $V$ be an \emph{irreducible} affine variety, so that $\C[V]$ is an integral domain.

\begin{definition}[Function Field]
	The \vocab{function field} of an irreducible affine variety $V$ is the field of fractions
	$$\C(V) = \Frac(\C[V]).$$
	Its elements are called \vocab{rational functions} on $V$.
\end{definition}

\begin{definition}
	A rational function $\varphi \in \C(V)$ is \vocab{regular} (or \vocab{defined}) at $P \in V$ if it can be written as $\varphi = f/g$ with $f,g \in \C[V]$ and $g(P) \neq 0$. The set of points at which $\varphi$ is regular is the \vocab{domain} of $\varphi$.
\end{definition}

If $\varphi = f/g$ then $\varphi$ genuinely is a function on the Zariski open set $V \setminus \V(g)$, and this set is nonempty (as $g \neq 0$ in $\C[V]$) hence dense, since $V$ is irreducible. The domain of $\varphi$ is the union of all such open sets, so it is a dense open subset of $V$. A rational function is thus a function defined `almost everywhere'.

\begin{example}
	Consider $\varphi = X_1^2/X_2 \in \C(\A^2)$. As written it is defined off the locus $X_2 = 0$. Note that $X_1^3/(X_1X_2)$ is the \emph{same} element of $\C(\A^2)$, even though the naive domain of the second expression is smaller. So an element of $\C(V)$ is not literally a function with a well-defined domain until we take the union over all representations. Here the domain really is $\A^2 \setminus \V(X_2)$.
\end{example}

\begin{remark}
	A slicker way to say all this: a rational function on $V$ is an equivalence class of pairs $(U,f)$ where $U \subseteq V$ is nonempty open and $f \colon U \to \C$ is a regular function on $U$, and $(U,f) \sim (U',f')$ if $f = f'$ on some nonempty open subset of $U \cap U'$. Irreducibility guarantees that $U \cap U'$ is nonempty, so this really is an equivalence relation.
\end{remark}

The functions regular at a fixed point form a ring, and its structure is the local algebraic invariant we will use over and over again.

\begin{definition}[Local Ring]
	A ring $R$ is a \vocab{local ring} if it has a unique maximal ideal.
\end{definition}

\begin{definition}
	Let $V$ be an irreducible variety and $P \in V$. The \vocab{local ring of $V$ at $P$} is
	$$\Oh_{V,P} = \{\varphi \in \C(V) : \varphi \text{ is regular at } P\}.$$
\end{definition}

\begin{lemma}\label{lem:localcrit}
	A ring $R$ is a local ring if and only if $R \setminus R^\times$ is an ideal, in which case it is the unique maximal ideal.
\end{lemma}

\begin{proof}
	An ideal $A \trianglelefteq R$ is proper if and only if $A$ contains no unit, i.e.\ $A \subseteq R\setminus R^\times$. So if $R \setminus R^\times$ is an ideal, it is proper and contains every proper ideal, hence is the unique maximal ideal.

	Conversely, suppose $R$ is local with maximal ideal $\mm$. Then $\mm \subseteq R\setminus R^\times$. If $x \in R \setminus R^\times$ then $(x) \neq R$, so $(x)$ is contained in a maximal ideal, which must be $\mm$. Hence $R\setminus R^\times = \mm$ is an ideal.
\end{proof}

\begin{proposition}
	$\Oh_{V,P}$ is a local ring, with maximal ideal
	$$\mm_{V,P} = \{\varphi \in \Oh_{V,P} : \varphi(P) = 0\}.$$
	Its units are exactly the $\varphi$ with $\varphi(P) \neq 0$, and $\Oh_{V,P}/\mm_{V,P} \cong \C$.
\end{proposition}

\begin{proof}
	If $\varphi = f/g \in \Oh_{V,P}$ with $g(P) \neq 0$, then $\varphi$ is invertible in $\Oh_{V,P}$ if and only if $g/f \in \Oh_{V,P}$, which happens exactly when $f(P) \neq 0$, i.e.\ $\varphi(P) \neq 0$. So $\Oh_{V,P}\setminus \Oh_{V,P}^\times = \mm_{V,P}$, which is an ideal (being the kernel of the evaluation homomorphism $\varphi \mapsto \varphi(P)$). Lemma~\ref{lem:localcrit} applies. Evaluation is surjective onto $\C$, giving the last statement.
\end{proof}

\begin{example}
	Take $V = \A^1$ and $P = 0$. Then
	$$\Oh_{\A^1,0} = \left\{ \frac{f(t)}{g(t)} : g(0) \neq 0 \right\} \subseteq \C(t),$$
	with maximal ideal $(t)$. This ring sits inside the ring $\C[[t]]$ of formal power series: expanding $1/g(t)$ as a power series (possible since $g(0) \neq 0$) embeds $\Oh_{\A^1,0} \hookrightarrow \C[[t]]$. The ring $\C[[t]]$ is also local, with maximal ideal $(t)$, and one should think of $\Oh_{\A^1,0}$ as the ring of `algebraic' functions with a power series expansion at $0$.
\end{example}

The local ring $\Oh_{V,P}$ is precisely the localisation of $\C[V]$ at the prime ideal $I(P)/I(V)$, in the sense of commutative algebra. It records everything about $V$ near $P$, and much of the second half of these notes consists of extracting geometry from these rings.

\section{Projective Varieties}

Affine varieties are not quite the right objects. Two problems recur: distinct parallel lines in $\A^2$ fail to meet, and affine varieties are never compact (in the Euclidean topology) unless they are finite. Both are cured by working in projective space, which adds `points at infinity' in a canonical way. The payoff is enormous: theorems become clean, intersection numbers become constant, and the geometry of curves becomes tractable.

\subsection{Projective Space}

\begin{definition}[Projectivisation]
	Let $U$ be a finite-dimensional complex vector space. The \vocab{projectivisation} $\PP(U)$ is the set of one-dimensional linear subspaces (`lines through the origin') of $U$. We write $\PP^n = \PP(\C^{n+1})$, \vocab{projective $n$-space}.
\end{definition}

We index coordinates on $\C^{n+1}$ by $0,1,\dots,n$. A line through the origin is spanned by some nonzero $(a_0,\dots,a_n)$, and two vectors span the same line exactly when they differ by a nonzero scalar. We write
$$(a_0 : a_1 : \cdots : a_n) \in \PP^n$$
for the point corresponding to the line spanned by $(a_0,\dots,a_n)$, and call these \vocab{homogeneous coordinates}. Thus
$$\PP^n = \frac{\C^{n+1}\setminus\{\mathbf{0}\}}{\C^\times},$$
where $\C^\times$ acts by scaling. For instance $(2:1:-2) = (4:2:-4) = (-1:-\tfrac12:1)$ in $\PP^2$.

\begin{center}
\begin{tikzpicture}[x=1cm,y=1cm]
	\draw[<->] (-2.9,0) -- (2.9,0) node[right] {$x_0$};
	\draw[<->] (0,-2.1) -- (0,2.4) node[above] {$x_1$};
	\foreach \xi in {-2.0,-1.0,-0.4,0.4,1.0,2.0}{
		\draw[color=blue, thin] (-1.15*\xi,-1.84) -- (1.15*\xi,1.84);
		\filldraw[color=blue] (\xi,1.6) circle (1.3pt);
	}
	\draw[very thick, color=blue] (-2.7,0) -- (2.7,0);
	\draw[thin, densely dashed] (-2.7,1.6) -- (2.7,1.6);
	\node[above left] at (-2.7,1.6) {$x_1 = 1$};
	\filldraw (0,0) circle (1.5pt);
	\node[below left] at (-0.02,-0.02) {$\mathbf{0}$};
	\node[below right, color=blue] at (2.2,0) {$\ell_\infty$};
\end{tikzpicture}
\end{center}

The picture above (drawn over $\R$ for legibility) shows $\PP^1$ as the set of lines through the origin in the plane. Almost every such line meets the horizontal line $x_1 = 1$ in exactly one point, which gives an identification of most of $\PP^1$ with a copy of $\A^1$; the exception is the horizontal axis itself, which is `the point at infinity'.

Let us record this decomposition in general. Splitting according to whether $a_0 = 0$,
\begin{align*}
	\PP^n &= \{(a_0 : \cdots : a_n) : a_0 \neq 0\} \sqcup \{(0 : a_1 : \cdots : a_n)\}\\
	&= \{(1 : b_1 : \cdots : b_n) : b_i \in \C\} \sqcup \PP^{n-1}\\
	&\cong \A^n \sqcup \PP^{n-1},
\end{align*}
where the first identification uses that we may scale $a_0$ to $1$, uniquely. Iterating,
$$\PP^n = \A^n \sqcup \A^{n-1} \sqcup \cdots \sqcup \A^1 \sqcup \{\text{pt}\}.$$
In particular $\PP^1 = \A^1 \sqcup \{\infty\}$ is the Riemann sphere, and $\PP^2$ is an affine plane together with a `line at infinity'.

\begin{definition}[Standard Affine Patches]
	For $0 \leq j \leq n$ let $H_j = \{(a_0:\cdots:a_n) : a_j = 0\}$, the \vocab{$j$th coordinate hyperplane}, and let $U_j = \PP^n \setminus H_j$ be the \vocab{$j$th standard affine patch}.
\end{definition}

There is a bijection $U_j \to \A^n$ given by
$$(a_0 : \cdots : a_n) \longmapsto \left(\frac{a_0}{a_j},\dots,\widehat{\frac{a_j}{a_j}},\dots,\frac{a_n}{a_j}\right),$$
where the hat means `omit', with inverse $(b_1,\dots,b_n) \mapsto (b_1 : \cdots : b_{j-1} : 1 : b_j : \cdots : b_n)$. Each $H_j$ is a copy of $\PP^{n-1}$. Since a point of $\PP^n$ has some nonzero coordinate, $\{U_0,\dots,U_n\}$ covers $\PP^n$: projective space is glued together out of $n+1$ copies of affine space.

\begin{definition}
	The \vocab{Zariski} (resp.\ \vocab{Euclidean}) topology on $\PP^n$ is the quotient topology induced from the Zariski (resp.\ Euclidean) topology on $\C^{n+1}\setminus\{\mathbf{0}\}$ by the quotient map $\pi \colon \C^{n+1}\setminus\{\mathbf 0\} \to \PP^n$.
\end{definition}

\begin{proposition}
	$\PP^n$ is compact in the Euclidean topology.
\end{proposition}

\begin{proof}
	The unit sphere $S^{2n+1} \subseteq \C^{n+1}\setminus\{\mathbf 0\}$ is compact, and every line through the origin meets it, so $\pi|_{S^{2n+1}}$ is surjective and continuous. A continuous image of a compact space is compact.
\end{proof}

This is the second reason to prefer $\PP^n$: it is compact, whereas $\A^n$ is not. Compactness will be responsible for a great deal of good behaviour, most notably the fact that the image of a projective variety under a morphism is closed.

\subsection{Homogeneous Ideals and Projective Varieties}

We cannot simply take vanishing loci of arbitrary polynomials in $\PP^n$.

\begin{example}
	The polynomial $X_0 + 1 \in \C[X_0,X_1]$ does not define a function on $\PP^1$, nor even a well-defined zero set: at the point $(-1:0) = (1:0)$ the representative $(-1,0)$ gives $0$ while $(1,0)$ gives $2$.
\end{example}

The fix is to restrict to polynomials whose vanishing is scale-invariant.

\begin{definition}[Homogeneous Polynomial]
	A \vocab{monomial} in $\C[X_0,\dots,X_n]$ is a product $X_0^{d_0}\cdots X_n^{d_n}$ with $d_i \geq 0$; its \vocab{degree} is $\sum_i d_i$. A polynomial is \vocab{homogeneous of degree $d$} if it is a $\C$-linear combination of monomials of degree $d$.
\end{definition}

Every $f \in \C[\mathbf X]$ decomposes uniquely as a finite sum $f = \sum_{r \geq 0} f_{[r]}$ with $f_{[r]}$ homogeneous of degree $r$; we call $f_{[r]}$ the \vocab{homogeneous parts} of $f$.

\begin{lemma}
	If $F$ is homogeneous of degree $d$ then $F(\lambda \mathbf{a}) = \lambda^d F(\mathbf a)$. In particular, whether $F(\mathbf a) = 0$ depends only on the point $(\mathbf a) \in \PP^n$, not on the representative.
\end{lemma}

\begin{proof}
	Each monomial of degree $d$ scales by $\lambda^d$.
\end{proof}

\begin{definition}[Homogeneous Ideal]
	An ideal $I \trianglelefteq \C[X_0,\dots,X_n]$ is \vocab{homogeneous} if it can be generated by homogeneous polynomials (possibly of different degrees).
\end{definition}

\begin{lemma}\label{lem:homid}
	An ideal $I$ is homogeneous if and only if for every $f \in I$, all homogeneous parts $f_{[r]}$ lie in $I$.
\end{lemma}

\begin{proof}
	($\Leftarrow$) Take any generating set and replace each generator by its homogeneous parts.

	($\Rightarrow$) Let $I$ be generated by homogeneous $g_1,\dots,g_s$ with $\deg g_j = d_j$, and let $f = \sum_j h_j g_j \in I$. Splitting each $h_j$ into homogeneous parts $h_{j[r]}$, the degree $r$ part of $f$ is $f_{[r]} = \sum_j h_{j[r-d_j]}g_j \in I$.
\end{proof}

\begin{definition}[Projective Variety]
	Let $I$ be a homogeneous ideal of $\C[X_0,\dots,X_n]$. Its \vocab{vanishing locus} is
	$$\V(I) = \{P \in \PP^n : F(P) = 0 \text{ for all homogeneous } F \in I\}.$$
	A \vocab{projective variety} is a subset of $\PP^n$ of this form.
\end{definition}

By Lemma~\ref{lem:homid} this is the same as requiring $f(P) = 0$ for all $f \in I$, interpreting `$f(P) = 0$' as `$f(\mathbf a) = 0$ for one, equivalently every, representative $\mathbf a$'. As in the affine case, the projective varieties are the closed sets of a topology; we omit the (identical) verification.

\begin{example}
	Let $U \leq \C^{n+1}$ be a linear subspace, say the kernel of a collection of linear forms $L_j$. Then $\PP(U) = \V(\{L_j\}) \subseteq \PP^n$ is a projective variety, called a \vocab{projective linear subspace}. Projective linear subspaces of $\PP^n$ are in bijection with linear subspaces of $\C^{n+1}$; those of dimension $1$ are called \vocab{lines}, and those of dimension $n-1$ are \vocab{hyperplanes}.
\end{example}

\begin{example}
	Two distinct lines in $\PP^2$ always meet, in exactly one point. Indeed, they correspond to distinct $2$-dimensional subspaces $U_1, U_2 \leq \C^3$, and $\dim(U_1 \cap U_2) \geq 2 + 2 - 3 = 1$, with equality since $U_1 \neq U_2$. This is the first fruit of working projectively: there are no parallel lines.
\end{example}

The group $\GL_{n+1}(\C)$ acts on $\PP^n$ by acting on homogeneous coordinates. Scalar matrices act trivially, so the action factors through
$$\PGL_{n+1}(\C) = \GL_{n+1}(\C)/\C^\times,$$
and this group acts transitively on $\PP^n$ (indeed on ordered $(n+2)$-tuples of points in general position). We refer to elements of $\PGL_{n+1}(\C)$ as \vocab{projective changes of coordinates}, and we will use them freely to put varieties into convenient position.

\subsection{Affine Charts and Projective Closure}

We have $\PP^n = U_0 \cup \cdots \cup U_n$ with each $U_j \cong \A^n$. The relationship between projective varieties and their affine pieces goes both ways.

Suppose $V = \V(I) \subseteq \PP^n$ is a projective variety. Identify $U_0 \cong \A^n$ with coordinates $Y_i = X_i/X_0$. Then
$$V \cap U_0 = \V(I_0), \qquad I_0 = \{F(1,Y_1,\dots,Y_n) : F \in I \text{ homogeneous}\},$$
an affine variety. So \emph{the intersection of a projective variety with an affine chart is an affine variety}: locally, projective varieties look affine.

Conversely, given an affine variety $W \subseteq \A^n \cong U_0$, we may take its closure in $\PP^n$. The following makes this explicit.

\begin{definition}[Homogenisation]
	Let $f \in \C[Y_1,\dots,Y_n]$ have total degree $d$. Its \vocab{homogenisation} is
	$$f^h(X_0,\dots,X_n) = X_0^{d}\, f\!\left(\frac{X_1}{X_0},\dots,\frac{X_n}{X_0}\right) \in \C[X_0,\dots,X_n],$$
	which is homogeneous of degree $d$ and not divisible by $X_0$. For an ideal $I \trianglelefteq \C[\mathbf Y]$ we write $I^h$ for the ideal generated by $\{f^h : f \in I\}$; this is homogeneous. The \vocab{projective closure} of an affine variety $V \subseteq \A^n$ is $\V(I(V)^h) \subseteq \PP^n$.
\end{definition}

\begin{example}
	If $f(Y_1,Y_2) = 1 + Y_1^2 + Y_1Y_2^2$ then $d = 3$ and $f^h = X_0^3 + X_0X_1^2 + X_1X_2^2$.
\end{example}

\begin{remark}
	It is essential to homogenise \emph{every} element of $I$, not merely a generating set. If $I = (f_1,\dots,f_r)$ then $(f_1^h,\dots,f_r^h) \subseteq I^h$ but the containment is often strict. (It is an equality when $I$ is principal.) For example, in $\C[Y_1,Y_2,Y_3]$ take the ideal of the affine twisted cubic $I = (Y_2 - Y_1^2, Y_3 - Y_1^3)$; homogenising the two generators gives $X_0X_2 - X_1^2$ and $X_0^2X_3 - X_1^3$, whose common zero locus contains the whole line $\{X_0 = X_1 = 0\}$, which is not in the projective closure of the twisted cubic. One needs also the homogenisation of $Y_1Y_3 - Y_2^2 \in I$, namely $X_1X_3 - X_2^2$.
\end{remark}

\begin{proposition}
	Let $V \subseteq \A^n = U_0$ be an affine variety. Then the Zariski closure $\overline V$ of $V$ in $\PP^n$ equals the projective closure $\V(I(V)^h)$.
\end{proposition}

\begin{proof}
	Write $I = I(V)$, which we may assume radical. Certainly $\V(I^h)$ is closed and contains $V$ (as $f^h$ restricted to $U_0$ is $f$), so $\overline V \subseteq \V(I^h)$.

	Conversely, let $Y = \V(J) \supseteq V$ be any closed set, $J$ homogeneous. Take a homogeneous $G \in J$; we may write $G = X_0^e\, f^h$ for some $f \in \C[\mathbf Y]$ and $e \geq 0$, by dehomogenising and re-homogenising. Now $G$ vanishes on $V$, so $f = G(1,\mathbf Y)$ vanishes on $V \subseteq \A^n$, so $f \in I(V) = I$. Hence $f^h \in I^h$ and so $G = X_0^e f^h \in I^h$. Therefore $J \subseteq I^h$ and $\V(I^h) \subseteq \V(J) = Y$. Taking $Y = \overline V$ gives the reverse inclusion.
\end{proof}

\begin{remark}
	The operations do not quite invert one another. If $V \subseteq \PP^n$ is projective and $W = V \cap U_0$, then $\overline W$ need not be $V$: take $V = \V(X_0)$, so $W = \varnothing$ and $\overline W = \varnothing \neq V$. The problem is that $V$ may have components entirely at infinity.
\end{remark}

\begin{example}\label{ex:conicpatches}
	Let $V = \V(X_0X_1 - X_2^2) \subseteq \PP^2$, a conic. Look at it in the three standard patches.
	\begin{itemize}
		\item In $U_0$, setting $X_0 = 1$ and writing $(y_1,y_2) = (X_1/X_0, X_2/X_0)$, we get $y_1 = y_2^2$: a parabola.
		\item In $U_1$ we similarly get a parabola.
		\item In $U_2$, setting $X_2 = 1$ and writing $(z_0,z_1) = (X_0/X_2, X_1/X_2)$, we get $z_0z_1 = 1$: a rectangular hyperbola.
	\end{itemize}
	So the parabola and the hyperbola are two views of the same projective curve, seen from different affine charts. The ellipse is a third. That the conics of school geometry become a single object is one of the most satisfying features of projective geometry.
\end{example}

It is worth drawing the picture that goes with this. Fix the chart $U_0 \cong \A^2$ and think of $\PP^2$ as this affine plane together with the line at infinity $\ell_\infty = \V(X_0)$, whose points record the directions in which one can head off to infinity. A parabola in the affine plane has both of its arms heading off in the same (vertical) direction, so in $\PP^2$ they meet at a single point of $\ell_\infty$, and the projective conic is a closed curve.

\begin{center}
\begin{tikzpicture}[x=1cm,y=1cm]
	\draw[thin, color=gray] (-3.1,-1.5) rectangle (3.1,2.35);
	\node[below right, color=gray] at (-3.05,-0.55) {$U_0 \cong \A^2$};
	\draw[very thick] (-3.5,2.8) -- (3.5,2.8);
	\node[above right] at (1.35,2.85) {$\ell_\infty = \V(X_0)$};
	\draw[thick, color=blue, smooth] plot coordinates {(-3.09,2.342) (-2.90,1.944) (-2.50,1.188) (-2.00,0.400) (-1.50,-0.213) (-1.00,-0.650) (-0.50,-0.913) (0.00,-1.000) (0.50,-0.913) (1.00,-0.650) (1.50,-0.213) (2.00,0.400) (2.50,1.188) (2.90,1.944) (3.09,2.342)};
	\draw[thick, color=blue, densely dashed] (-3.09,2.342) .. controls (-2.50,2.74) and (-1.20,2.80) .. (0,2.80);
	\draw[thick, color=blue, densely dashed] (3.09,2.342) .. controls (2.50,2.74) and (1.20,2.80) .. (0,2.80);
	\filldraw (0,2.8) circle (2pt);
	\node[below] at (0,2.72) {$P_\infty$};
	\node[color=blue, right] at (0.15,-1.05) {$C \cap U_0$};
\end{tikzpicture}
\end{center}

Here is the general classification, which is just the classification of quadratic forms from IB Linear Algebra.

\begin{theorem}
	Let $Q = \V(F) \subseteq \PP^n$ with $F$ homogeneous of degree $2$. After a projective change of coordinates, $F = X_0^2 + X_1^2 + \cdots + X_{r-1}^2$, where $r$ is the rank of the quadratic form $F$.
\end{theorem}

\begin{proof}
	Over $\C$ every quadratic form is equivalent to $\sum_{i<r} X_i^2$ where $r$ is its rank.
\end{proof}

\subsection{The Projective Nullstellensatz}

Two homogeneous ideals may cut out the same projective variety for a trivial reason: the \vocab{irrelevant ideal} $\mm_0 = (X_0,\dots,X_n)$ has $\V(\mm_0) = \varnothing$ in $\PP^n$ (its affine cone is just the origin, which we deleted). So the correct statement has to allow for this.

\begin{definition}
	For $V \subseteq \PP^n$, let $I^h(V) \trianglelefteq \C[X_0,\dots,X_n]$ be the ideal generated by all homogeneous polynomials vanishing on $V$.
\end{definition}

\begin{theorem}[Projective Nullstellensatz]
	Let $I$ be a homogeneous ideal of $\C[X_0,\dots,X_n]$.
	\begin{enumerate}[label=(\roman*)]
		\item $\V(I) = \varnothing$ in $\PP^n$ if and only if $\sqrt I \supseteq (X_0,\dots,X_n)$.
		\item If $V = \V(I) \neq \varnothing$ then $I^h(V) = \sqrt I$.
	\end{enumerate}
\end{theorem}

\begin{proof}
	Let $C(V) = \V(I) \subseteq \A^{n+1}$ denote the affine vanishing locus, the \vocab{affine cone}. Since $I$ is homogeneous, $C(V)$ is stable under scaling and contains $\mathbf 0$, and the quotient map restricts to a surjection $C(V)\setminus\{\mathbf 0\} \to \V(I) \subseteq \PP^n$. In particular $\V(I) = \varnothing$ if and only if $C(V) \subseteq \{\mathbf 0\}$.

	For (i), $C(V) \subseteq \{\mathbf 0\}$ means $I(C(V)) \supseteq (X_0,\dots,X_n)$, and by the affine Nullstellensatz $I(C(V)) = \sqrt I$.

	For (ii), the homogeneous polynomials vanishing on $\V(I) \subseteq \PP^n$ are exactly the homogeneous polynomials vanishing on $C(V) \subseteq \A^{n+1}$, so $I^h(V) = I(C(V)) = \sqrt I$ by the affine Nullstellensatz again (using Lemma~\ref{lem:homid} to see that $I(C(V))$ is homogeneous).
\end{proof}

Exactly as before, we get:

\begin{proposition}
	\begin{enumerate}[label=(\roman*)]
		\item Every projective variety is a finite union of irreducible projective varieties, uniquely if the decomposition is minimal.
		\item A projective variety $V$ is irreducible if and only if $I^h(V)$ is prime.
	\end{enumerate}
\end{proposition}

\begin{proof}
	(i) is verbatim the affine argument. For (ii), the only new point is that if a homogeneous ideal $I$ is not prime then there exist \emph{homogeneous} $F, G \notin I$ with $FG \in I$: take $f,g \notin I$ with $fg \in I$, and let $F$ (resp.\ $G$) be a homogeneous part of $f$ (resp.\ $g$) of largest degree not lying in $I$; then the top-degree part of $fg$ that could fail to lie in $I$ is $FG$. With that in hand, the affine proof applies.
\end{proof}

\begin{proposition}\label{prop:opendense}
	Let $V \subseteq \PP^n$ be irreducible and $W \subsetneq V$ a proper closed subvariety. Then $V \setminus W$ is dense in $V$.
\end{proposition}

\begin{proof}
	Let $F$ be homogeneous and vanish on $V \setminus W$. Since $W \neq V$, by the projective Nullstellensatz there is a homogeneous $G \in I^h(W) \setminus I^h(V)$. Then $FG$ vanishes on $W$ (because $G$ does) and on $V \setminus W$ (because $F$ does), hence on all of $V$. As $I^h(V)$ is prime and $G \notin I^h(V)$, we get $F \in I^h(V)$. So every homogeneous polynomial vanishing on $V\setminus W$ vanishes on $V$, which says exactly that $V \setminus W$ is dense.
\end{proof}

This is the projective form of the principle that nonempty open subsets are large. Morally, $W$ is lower-dimensional than $V$, and removing something lower-dimensional cannot disconnect or shrink an irreducible variety.

\subsection{Rational Functions on Projective Varieties}

There is no useful analogue of the coordinate ring in the projective setting: a homogeneous polynomial has a well-defined zero set on $\PP^n$ but not a well-defined value. However, a \emph{ratio} of two homogeneous polynomials of the same degree \emph{does} have a well-defined value wherever the denominator does not vanish, since the scaling factors $\lambda^d$ cancel.

\begin{definition}[Function Field]
	Let $V \subseteq \PP^n$ be an irreducible projective variety. Its \vocab{function field} is
	$$\C(V) = \left\{ \frac{F}{G} : F,G \in \C[\mathbf X] \text{ homogeneous of the same degree},\ G \notin I^h(V)\right\}\bigg/ \sim,$$
	where $F_1/G_1 \sim F_2/G_2$ if and only if $F_1G_2 - F_2G_1 \in I^h(V)$.
\end{definition}

\begin{lemma}
	The relation $\sim$ is an equivalence relation, and $\C(V)$ is a field.
\end{lemma}

\begin{proof}
	Reflexivity and symmetry are clear. For transitivity, suppose $F_1G_2 \equiv F_2G_1$ and $F_2G_3 \equiv F_3G_2$ modulo $I^h(V)$. Working in the integral domain $R = \C[\mathbf X]/I^h(V)$ (a domain since $I^h(V)$ is prime), we get $\bar F_1 \bar G_2 \bar G_3 = \bar F_2\bar G_1\bar G_3 = \bar F_3 \bar G_2 \bar G_1$, so $\bar G_2(\bar F_1 \bar G_3 - \bar F_3\bar G_1) = 0$. As $\bar G_2 \neq 0$ and $R$ is a domain, $F_1G_3 - F_3G_1 \in I^h(V)$.

	The field operations are the evident ones (note that $F_1/G_1 + F_2/G_2 = (F_1G_2+F_2G_1)/G_1G_2$ is again a ratio of forms of equal degree), and $F/G$ is invertible whenever $F \notin I^h(V)$.
\end{proof}

\begin{proposition}
	$\C(V)$ is a finitely generated field extension of $\C$.
\end{proposition}

\begin{proof}
	Some coordinate, say $X_0$, does not vanish identically on $V$. We claim $\C(V) = \C(x_1,\dots,x_n)$ where $x_i$ is the class of $X_i/X_0$. Given a ratio $F/G$ of forms of degree $d$, divide numerator and denominator by $X_0^d$; then $F/X_0^d$ is a polynomial in the $X_i/X_0$, and likewise for $G$. So $F/G$ lies in the field generated by the $x_i$.
\end{proof}

\begin{corollary}\label{cor:funcfieldpatch}
	Let $V \subseteq \PP^n$ be irreducible with $V \not\subseteq \{X_0 = 0\}$, and let $V_0 = V \cap U_0$, an irreducible affine variety. Then $\C(V) \cong \C(V_0)$.
\end{corollary}

\begin{proof}
	$\C[V_0] = \C[x_1,\dots,x_n]$ where $x_i$ is the restriction of $X_i/X_0$, so $\C(V_0) = \Frac \C[V_0]$ is generated by the $x_i$, and the proof above shows $\C(V)$ is too. One checks the relations match, giving an isomorphism.
\end{proof}

So the function field is a birational invariant that can be computed in any affine chart. Regularity and local rings are defined exactly as before: $\varphi \in \C(V)$ is \vocab{regular at $P$} if it can be written $F/G$ with $G(P) \neq 0$, and $\Oh_{V,P} \subseteq \C(V)$ is the local ring of germs at $P$, with maximal ideal $\mm_P$. Under the isomorphism of Corollary~\ref{cor:funcfieldpatch}, $\Oh_{V,P} \cong \Oh_{V_0,P}$ for $P \in V_0$.

\subsection{Rational Maps and Morphisms}

\begin{proposition}
	Let $F_0,\dots,F_m \in \C[X_0,\dots,X_n]$ be homogeneous of the same degree $d$, not all zero. Then $\mathbf F = (F_0,\dots,F_m)$ descends to a well-defined map
	$$\varphi \colon \PP^n \setminus \V(F_0,\dots,F_m) \to \PP^m, \qquad P \mapsto (F_0(\mathbf a) : \cdots : F_m(\mathbf a)),$$
	where $\mathbf a$ is any representative of $P$.
\end{proposition}

\begin{proof}
	Replacing $\mathbf a$ by $\lambda\mathbf a$ multiplies every $F_j(\mathbf a)$ by $\lambda^d$, which does not change the point of $\PP^m$. The condition $P \notin \V(F_0,\dots,F_m)$ ensures not all coordinates are zero.
\end{proof}

Two tuples can define the same map on the overlap of their domains: multiplying all $F_j$ by a fixed homogeneous $G$ gives a tuple defining the same map wherever both are defined, but on a smaller set. The following definition is designed to handle this.

\begin{definition}[Rational Map]
	Let $V \subseteq \PP^n$ be irreducible. A \vocab{rational map} $\varphi \colon V \dto \PP^m$ is an equivalence class of tuples $(F_0,\dots,F_m)$ of homogeneous polynomials of a common degree, not all in $I^h(V)$, where
	$$(F_0,\dots,F_m) \sim (G_0,\dots,G_m) \iff F_iG_j - F_jG_i \in I^h(V) \text{ for all } i,j.$$
	A point $P \in V$ is a \vocab{regular point} of $\varphi$ if some representative has $F_j(P) \neq 0$ for some $j$. The \vocab{domain} $\dom\varphi$ is the set of regular points. If $\dom\varphi = V$, we call $\varphi$ a \vocab{morphism} and write $\varphi \colon V \to \PP^m$.
\end{definition}

If $W \subseteq \PP^m$ is a projective variety, a rational map (or morphism) $V \dto W$ is one $V \dto \PP^m$ whose image lies in $W$. An \vocab{isomorphism} is a morphism with a two-sided inverse morphism.

The subtlety here -- that a rational map may be regular at a point without any \emph{obvious} representative witnessing it -- is important, and the next example is the one to remember.

\begin{example}[Projection from a point]\label{ex:projpoint}
	Let $P = (0:0:1) \in \PP^2$ and define $\pi \colon \PP^2 \dto \PP^1$ by $\pi(a_0:a_1:a_2) = (a_0:a_1)$. This is undefined exactly at $P$, and it is regular everywhere else. Geometrically, $\pi$ sends a point $Q \neq P$ to the line $PQ$, thought of as a point of the $\PP^1$ of lines through $P$.
\end{example}

\begin{example}\label{ex:conicP1}
	Let $C = \V(X_0X_2 - X_1^2) \subseteq \PP^2$ and restrict the projection $\pi$ of the previous example to $C$. Note $P = (0:0:1) \in C$, so \emph{a priori} $\pi|_C$ is undefined at $P$. But $\pi|_C$ is represented by $(X_0,X_1)$, and we may look for another representative $(F_0,F_1)$ with
	$$F_0X_1 - F_1X_0 \in I^h(C) = (X_0X_2 - X_1^2).$$
	The pair $(X_1,X_2)$ works, since $X_1 \cdot X_1 - X_2 \cdot X_0 = -(X_0X_2-X_1^2)$. At $P = (0:0:1)$ we have $X_2(P) = 1 \neq 0$, so $\pi|_C$ \emph{is} regular at $P$, with value $\pi(P) = (0:1)$. Hence $\pi|_C \colon C \to \PP^1$ is a morphism.
\end{example}

\begin{center}
\begin{tikzpicture}[x=1.15cm,y=1.15cm]
	\draw[thick, color=blue] (0,0.9) circle (0.9);
	\node[color=blue] at (1.25,1.5) {$C$};
	\draw[thick] (-2.5,-1.1) -- (2.5,-1.1);
	\node[right] at (2.55,-1.1) {$\PP^1$};
	\foreach \qx/\qy/\lx in {-0.846/0.592/-2.031, -0.516/0.163/-0.914, 0.516/0.163/0.914, 0.846/0.592/2.031}{
		\draw[thin, densely dashed, color=gray] (0,1.8) -- (\lx,-1.1);
		\filldraw[color=blue] (\qx,\qy) circle (1.4pt);
		\filldraw (\lx,-1.1) circle (1.4pt);
	}
	\filldraw (0,1.8) circle (1.7pt);
	\node[above] at (0,1.87) {$P$};
	\node[color=blue, above left] at (-0.9,0.62) {$Q$};
\end{tikzpicture}
\end{center}

The picture shows the construction: a point $Q \in C$ is sent to the point where the line $PQ$ meets a fixed line. Since a line meets a conic in two points, and one of them is $P$, the other is determined, and this is why the map is a bijection.

\begin{proposition}\label{prop:conicisP1}
	Let $C = \V(F) \subseteq \PP^2$ with $F$ an irreducible homogeneous quadratic. Then $C \cong \PP^1$.
\end{proposition}

\begin{proof}
	An irreducible quadratic form in three variables has rank $3$; after a change of coordinates we may take $F = X_0X_2 - X_1^2$. By Example~\ref{ex:conicP1} we have a morphism $\pi \colon C \to \PP^1$. Define $\mu \colon \PP^1 \to \PP^2$ by
	$$\mu(Y_0:Y_1) = (Y_0^2 : Y_0Y_1 : Y_1^2).$$
	This is a morphism (the three forms have no common zero on $\PP^1$), and its image lies in $C$ since $Y_0^2\cdot Y_1^2 - (Y_0Y_1)^2 = 0$. Now $\pi(\mu(Y_0:Y_1)) = (Y_0^2 : Y_0Y_1) = (Y_0:Y_1)$ whenever $Y_0 \neq 0$, and using the other representative $(X_1,X_2)$ of $\pi$ we get $(Y_0Y_1 : Y_1^2) = (Y_0:Y_1)$ whenever $Y_1 \neq 0$; so $\pi\mu = \mathrm{id}$. Conversely, for $(a_0:a_1:a_2) \in C$ we have $a_0a_2 = a_1^2$, so $\mu(a_0:a_1) = (a_0^2 : a_0a_1 : a_1^2) = (a_0^2 : a_0a_1 : a_0a_2) = (a_0:a_1:a_2)$ when $a_0 \neq 0$, and similarly on the other patch. So $\mu\pi = \mathrm{id}$.
\end{proof}

So up to isomorphism there is only one smooth conic in $\PP^2$, and it is a line. The distinction between lines and conics is a matter of how they are embedded, not of their intrinsic geometry. We will see in due course that this ceases to be true in degree $3$ and beyond.

\begin{example}[Veronese embedding]
	Fix $n,d \geq 1$ and let $F_0,\dots,F_m$ be all the monomials of degree $d$ in $X_0,\dots,X_n$, so $m = \binom{n+d}{d}-1$. The \vocab{Veronese map} is
	$$\nu_d \colon \PP^n \to \PP^m, \qquad \mathbf a \mapsto (F_0(\mathbf a) : \cdots : F_m(\mathbf a)).$$
	It is a morphism, since the $F_j$ include $X_i^d$ for each $i$ and these have no common zero. One can show $\nu_d$ is injective with image a projective variety isomorphic to $\PP^n$. The map $\mu$ used in Proposition~\ref{prop:conicisP1} is $\nu_2$ for $n=1$. The point of the construction is that $\nu_d$ turns degree $d$ hypersurfaces into hyperplane sections.
\end{example}

\begin{example}[Segre embedding]
	Products of projective spaces are not projective spaces: $\PP^n\times\PP^m \not\cong \PP^{n+m}$. But they are still projective varieties. The \vocab{Segre map} is
	$$\sigma_{n,m} \colon \PP^n \times \PP^m \to \PP^{nm+n+m}, \qquad \big((x_i),(y_j)\big) \mapsto (x_iy_j)_{i,j},$$
	where we index the target coordinates by $Z_{ij}$, $0 \leq i \leq n$, $0 \leq j \leq m$. Note $(n+1)(m+1)-1 = nm+n+m$.
\end{example}

\begin{theorem}
	The Segre map $\sigma_{n,m}$ is a bijection onto the projective variety $\V(I)$, where $I$ is generated by the quadrics $Z_{ij}Z_{pq} - Z_{iq}Z_{pj}$.
\end{theorem}

\begin{proof}
	Clearly the image lies in $\V(I)$. Work on the patch $Z_{00} \neq 0$ and set $y_{ij} = Z_{ij}/Z_{00}$. The relations $Z_{ij}Z_{00} = Z_{i0}Z_{0j}$ become $y_{ij} = y_{i0}y_{0j}$, and one checks the remaining relations follow. So on this patch $\V(I)$ is the graph of the map $\A^n \times \A^m \to \A^{nm}$, $\big((y_{i0}),(y_{0j})\big) \mapsto (y_{i0}y_{0j})$, and hence
	$$(y_{ij}) \longmapsto \big((y_{10},\dots,y_{n0}),(y_{01},\dots,y_{0m})\big)$$
	is a two-sided inverse to $\sigma_{n,m}$ there. The same works on each patch $Z_{ij} \neq 0$.
\end{proof}

Concretely, $\sigma_{1,1} \colon \PP^1\times\PP^1 \to \PP^3$ has image the quadric surface $\V(Z_{00}Z_{11} - Z_{01}Z_{10}) \subseteq \PP^3$: identifying $\C^4$ with $2\times 2$ matrices, $\sigma_{1,1}$ is the outer product $(v,w)\mapsto vw^{\mathsf T}$, whose image is the set of matrices of rank at most one, cut out by the vanishing of the determinant. As a corollary, products of projective varieties are projective varieties.

\subsection{Birational Maps}

Composing rational maps is delicate: the image of $\dom\varphi$ might land entirely inside the indeterminacy locus of $\psi$. The following condition rules this out.

\begin{definition}[Dominant]
	A rational map $\varphi \colon V \dto W$ is \vocab{dominant} if $\varphi(\dom\varphi)$ is Zariski dense in $W$.
\end{definition}

\begin{proposition}
	If $\varphi \colon V \dto W$ is dominant then $\psi \circ \varphi$ is a well-defined rational map for any rational $\psi \colon W \dto Z$.
\end{proposition}

\begin{proof}
	$\dom\psi$ is a nonempty open subset of $W$, and $\varphi(\dom\varphi)$ is dense, so $\varphi^{-1}(\dom\psi)$ is a nonempty open subset of $\dom\varphi$. On it, $\psi\varphi$ is given by composing tuples of homogeneous polynomials, hence is a rational map.
\end{proof}

A dominant rational map $\varphi \colon V \dto W$ induces a $\C$-algebra map on function fields $\varphi^\star \colon \C(W) \to \C(V)$ by $f \mapsto f \circ \varphi$; this is a map of fields, hence injective.

\begin{definition}[Birational]
	Irreducible varieties $V, W$ are \vocab{birational} if there are dominant rational maps $\varphi \colon V \dto W$ and $\psi \colon W \dto V$ with $\psi\varphi$ and $\varphi\psi$ equal to the identity as rational maps. Such $\varphi$ is called a \vocab{birational map}.
\end{definition}

Equivalently, $V$ and $W$ are birational if they contain isomorphic nonempty open subsets. Every isomorphism is birational, but not conversely.

\begin{example}[Cremona transformation]
	Consider $\kappa \colon \PP^2 \dto \PP^2$ given by
	$$\kappa(X_0:X_1:X_2) = (X_1X_2 : X_0X_2 : X_0X_1),$$
	which is `$(x_0:x_1:x_2) \mapsto (1/x_0 : 1/x_1 : 1/x_2)$' cleared of denominators. It is undefined exactly at the three coordinate points $(1:0:0), (0:1:0), (0:0:1)$. One computes $\kappa \circ \kappa = \mathrm{id}$ as a rational map, so $\kappa$ is birational; it is not an isomorphism, since it is not everywhere defined. The Cremona transformation sends a general line to a conic, so it certainly does not preserve degrees.
\end{example}

\begin{remark}
	The group $\Bir(\PP^2)$ of birational automorphisms of $\PP^2$ is called the \emph{Cremona group}. A theorem of Noether and Castelnuovo says it is generated by $\PGL_3(\C)$ together with $\kappa$. The analogous statement in higher dimensions is false and the structure of $\Bir(\PP^n)$ for $n \geq 3$ remains mysterious.
\end{remark}

\begin{theorem}\label{thm:biratfield}
	Let $V, W$ be irreducible projective varieties. Then $V$ and $W$ are birational if and only if $\C(V) \cong \C(W)$ as $\C$-algebras.
\end{theorem}

\begin{proof}
	($\Rightarrow$) A birational map $\varphi$ gives $\varphi^\star \colon \C(W) \to \C(V)$ and $\varphi^{-1}$ gives an inverse, so these are isomorphisms.

	($\Leftarrow$) Suppose $\Theta \colon \C(W) \to \C(V)$ is a $\C$-algebra isomorphism. We may assume $V \subseteq \PP^n$ is not contained in $\{X_0 = 0\}$ and $W\subseteq\PP^m$ is not contained in $\{Y_0=0\}$. As we showed, $\C(V) = \C(x_1,\dots,x_n)$ with $x_i = X_i/X_0$ and $\C(W) = \C(y_1,\dots,y_m)$ with $y_j = Y_j/Y_0$. Write $\Theta(y_j) = f_j(x_1,\dots,x_n)$ for rational functions $f_j$ of $n$ variables. Clearing denominators and homogenising with respect to $X_0$, there are homogeneous $F_0,\dots,F_m$ of a common degree with
	$$f_j\left(\frac{X_1}{X_0},\dots,\frac{X_n}{X_0}\right) = \frac{F_j(X_0,\dots,X_n)}{F_0(X_0,\dots,X_n)}.$$
	Then $(F_0:\cdots:F_m)$ defines a rational map $\varphi \colon V \dto \PP^m$; its image lies in $W$, since any homogeneous $G \in I^h(W)$ pulls back to $G(f_1,\dots,f_m)$ up to a factor, which is $\Theta$ applied to the relation $G(y_1,\dots,y_m) = 0$, hence zero. Running the same argument with $\Theta^{-1}$ gives $\psi \colon W \dto V$, and the composites induce the identity on function fields, hence are the identity as rational maps.
\end{proof}

This should be compared with Corollary~\ref{cor:isocoord}: coordinate rings classify affine varieties up to isomorphism, and function fields classify irreducible varieties up to birational equivalence. Classifying varieties up to birational equivalence is thus the same as classifying finitely generated field extensions of $\C$, and for curves we will be able to do this completely.

\section{Tangent Spaces, Smoothness and Dimension}\label{sec:tangent}

So far our varieties have been rather formless: we have no notion of how big one is, and no way to distinguish the smooth curve $y = x^2$ from the cuspidal one $y^2 = x^3$. Both defects are repaired at once by the tangent space, which is a purely algebraic construction imitating the derivative.

\subsection{The Tangent Space to a Hypersurface}

Let $V = \V(f) \subseteq \A^n$ be a hypersurface and $P = (a_1,\dots,a_n) \in V$. A line through $P$ has the form
$$L = \{(a_1 + b_1t,\dots,a_n+b_nt) : t \in \C\}$$
for some $\mathbf b \in \C^n \setminus\{\mathbf 0\}$. Substituting into $f$ gives a polynomial in one variable,
$$g(t) = f(a_1+b_1t,\dots,a_n+b_nt) = \sum_{r \geq 0} c_r t^r,$$
whose roots are the parameters at which $L$ meets $V$. Since $P \in V$ we have $c_0 = 0$, so $t = 0$ is a root; and by the chain rule
$$c_1 = g'(0) = \sum_{i=1}^n b_i \frac{\partial f}{\partial X_i}(P).$$
The line meets $V$ at $P$ `to order at least two' -- geometrically, is tangent -- exactly when $c_1 = 0$ as well.

\begin{definition}[Tangent Space to a Hypersurface]
	Let $V = \V(f) \subseteq \A^n$ and $P \in V$. The \vocab{(affine embedded) tangent space} to $V$ at $P$ is the affine subspace
	$$T^{\mathrm{aff}}_{V,P} = \V\left(\sum_{i=1}^n \frac{\partial f}{\partial X_i}(P)\,(X_i - a_i)\right) \subseteq \A^n.$$
	A line $L$ through $P$ is \vocab{tangent} to $V$ at $P$ if $L \subseteq T^{\mathrm{aff}}_{V,P}$.
\end{definition}

The polynomial cutting out $T^{\mathrm{aff}}_{V,P}$ is linear, so $T^{\mathrm{aff}}_{V,P}$ is either a hyperplane (dimension $n-1$), when some partial derivative is nonzero at $P$, or all of $\A^n$, when they all vanish.

\begin{definition}[Smooth and Singular Points]
	With the notation above, $P$ is a \vocab{singular point} of $V$ if $\frac{\partial f}{\partial X_i}(P) = 0$ for all $i$, equivalently $\dim T^{\mathrm{aff}}_{V,P} = n$. Otherwise $P$ is a \vocab{smooth} (or \vocab{nonsingular}) point. $V$ is \vocab{smooth} if all its points are.
\end{definition}

This is the \vocab{Jacobian criterion}: a point is singular precisely when the gradient of the defining equation vanishes there. Let us test it on the standard examples.

\begin{example}[Nodal cubic]
	Let $C = \V(f)$ with $f = Y^2 - X^2(X+1) = Y^2 - X^3 - X^2$. Then
	$$\frac{\partial f}{\partial X} = -3X^2 - 2X, \qquad \frac{\partial f}{\partial Y} = 2Y.$$
	Both vanish when $Y = 0$ and $X(3X+2) = 0$, i.e.\ at $(0,0)$ and $(-2/3,0)$; but $(-2/3,0) \notin C$. So the unique singular point is the origin. Near the origin, $Y^2 \approx X^2$, and indeed the real picture shows two branches crossing transversally, with tangent lines $Y = \pm X$: this is a \vocab{node}.
	\begin{center}
	\begin{tikzpicture}[x=1.6cm,y=1.15cm]
		\draw[ultra thin, color=lightgray] (-1.4,-2.1) grid[step=0.5] (1.4,2.1);
		\draw[<->] (-1.5,0) -- (1.5,0);
		\draw[<->] (0,-2.2) -- (0,2.2);
		\draw[thin, color=gray] (-0.85,-0.85) -- (0.85,0.85);
		\draw[thin, color=gray] (-0.85,0.85) -- (0.85,-0.85);
		\draw[thick, color=blue, smooth] plot coordinates {(0.00,0.00) (-0.13,-0.12) (-0.25,-0.22) (-0.36,-0.29) (-0.46,-0.34) (-0.56,-0.37) (-0.64,-0.38) (-0.72,-0.38) (-0.78,-0.37) (-0.84,-0.34) (-0.89,-0.30) (-0.93,-0.25) (-0.96,-0.19) (-0.98,-0.13) (-1.00,-0.07) (-1.00,0.00) (-1.00,0.07) (-0.98,0.13) (-0.96,0.19) (-0.93,0.25) (-0.89,0.30) (-0.84,0.34) (-0.78,0.37) (-0.72,0.38) (-0.64,0.38) (-0.56,0.37) (-0.46,0.34) (-0.36,0.29) (-0.25,0.22) (-0.13,0.12) (0.00,0.00)};
		\draw[thick, color=blue, smooth] plot coordinates {(0.00,0.00) (0.06,0.06) (0.12,0.13) (0.18,0.19) (0.24,0.26) (0.30,0.34) (0.36,0.42) (0.42,0.50) (0.48,0.58) (0.54,0.66) (0.60,0.75) (0.65,0.84) (0.71,0.94) (0.77,1.03) (0.83,1.13) (0.89,1.23) (0.95,1.33) (1.01,1.44) (1.07,1.54) (1.13,1.65) (1.19,1.76) (1.25,1.88)};
		\draw[thick, color=blue, smooth] plot coordinates {(1.25,-1.88) (1.19,-1.76) (1.13,-1.65) (1.07,-1.54) (1.01,-1.44) (0.95,-1.33) (0.89,-1.23) (0.83,-1.13) (0.77,-1.03) (0.71,-0.94) (0.65,-0.84) (0.60,-0.75) (0.54,-0.66) (0.48,-0.58) (0.42,-0.50) (0.36,-0.42) (0.30,-0.34) (0.24,-0.26) (0.18,-0.19) (0.12,-0.13) (0.06,-0.06) (0.00,0.00)};
		\filldraw (0,0) circle (1.6pt);
		\node[below right] at (0.04,-0.06) {$P$};
	\end{tikzpicture}
	\end{center}
	The two grey lines are the tangent directions to the two branches; note that neither is `the' tangent line, and that the tangent space at $P$ is the whole plane.
\end{example}

\begin{example}[Cuspidal cubic]
	Let $C = \V(Y^2 - X^3)$. Then $\partial f/\partial X = -3X^2$ and $\partial f/\partial Y = 2Y$, both vanishing only at the origin, which lies on $C$. So again the origin is the unique singular point, but the singularity is of a different kind: the curve has a single branch with a \vocab{cusp}, parametrised by $t \mapsto (t^2,t^3)$.
	\begin{center}
	\begin{tikzpicture}[x=1.4cm,y=1.0cm]
		\draw[ultra thin, color=lightgray] (-0.3,-2.4) grid[step=0.5] (1.9,2.4);
		\draw[<->] (-0.4,0) -- (2.0,0);
		\draw[<->] (0,-2.5) -- (0,2.5);
		\draw[thick, color=blue, smooth] plot coordinates {(1.69,-2.20) (1.47,-1.79) (1.27,-1.43) (1.08,-1.12) (0.91,-0.87) (0.75,-0.65) (0.61,-0.47) (0.48,-0.33) (0.37,-0.22) (0.27,-0.14) (0.19,-0.08) (0.12,-0.04) (0.07,-0.02) (0.03,-0.01) (0.01,-0.00) (0.00,0.00) (0.01,0.00) (0.03,0.01) (0.07,0.02) (0.12,0.04) (0.19,0.08) (0.27,0.14) (0.37,0.22) (0.48,0.33) (0.61,0.47) (0.75,0.65) (0.91,0.87) (1.08,1.12) (1.27,1.43) (1.47,1.79) (1.69,2.20)};
		\filldraw (0,0) circle (1.6pt);
		\node[left] at (-0.06,-0.16) {$P$};
	\end{tikzpicture}
	\end{center}
\end{example}

\begin{example}[A smooth point]
	Contrast the previous two examples with $C = \V(f)$, $f = Y^2 - X^3 + X$, at the point $P = \big(-1/\sqrt3,\ y_0\big) \approx (-0.577, 0.620)$, where $y_0 = \sqrt{2}\,3^{-3/4}$. Here $\partial f/\partial Y = 2y_0 \neq 0$, so $P$ is a smooth point, the tangent space is a line, and there is a genuine tangent line: differentiating implicitly, its slope is
	$$\frac{3X^2-1}{2Y}\bigg|_P = \frac{3\cdot\frac13 - 1}{2y_0} = 0,$$
	so the tangent at $P$ is the horizontal line $Y = y_0$.
	In fact this curve is smooth everywhere: $\partial f/\partial Y = 2Y = 0$ forces $Y = 0$, and then $X^3 = X$ and $3X^2 = 1$ cannot both hold.
	\begin{center}
	\begin{tikzpicture}[x=1.6cm,y=1.05cm]
		\draw[ultra thin, color=lightgray] (-1.4,-2.3) grid[step=0.5] (2.1,2.3);
		\draw[<->] (-1.5,0) -- (2.2,0);
		\draw[<->] (0,-2.4) -- (0,2.4);
		\draw[thick, color=blue, smooth] plot coordinates {(0.00,0.00) (-0.00,0.06) (-0.02,0.13) (-0.04,0.19) (-0.06,0.25) (-0.10,0.31) (-0.14,0.36) (-0.18,0.42) (-0.23,0.47) (-0.29,0.51) (-0.35,0.55) (-0.41,0.58) (-0.47,0.60) (-0.53,0.62) (-0.59,0.62) (-0.65,0.61) (-0.71,0.59) (-0.77,0.56) (-0.82,0.52) (-0.86,0.47) (-0.90,0.41) (-0.94,0.34) (-0.96,0.26) (-0.98,0.18) (-1.00,0.09) (-1.00,0.00) (-1.00,-0.09) (-0.98,-0.18) (-0.96,-0.26) (-0.94,-0.34) (-0.90,-0.41) (-0.86,-0.47) (-0.82,-0.52) (-0.77,-0.56) (-0.71,-0.59) (-0.65,-0.61) (-0.59,-0.62) (-0.53,-0.62) (-0.47,-0.60) (-0.41,-0.58) (-0.35,-0.55) (-0.29,-0.51) (-0.23,-0.47) (-0.18,-0.42) (-0.14,-0.36) (-0.10,-0.31) (-0.06,-0.25) (-0.04,-0.19) (-0.02,-0.13) (-0.00,-0.06) (0.00,0.00)};
		\draw[thick, color=blue, smooth] plot coordinates {(1.90,2.23) (1.82,2.06) (1.75,1.90) (1.68,1.75) (1.61,1.61) (1.55,1.48) (1.49,1.35) (1.44,1.23) (1.38,1.12) (1.33,1.02) (1.29,0.92) (1.24,0.83) (1.21,0.74) (1.17,0.66) (1.14,0.58) (1.11,0.50) (1.08,0.43) (1.06,0.37) (1.04,0.30) (1.03,0.24) (1.02,0.18) (1.01,0.12) (1.00,0.06) (1.00,0.00) (1.00,-0.06) (1.01,-0.12) (1.02,-0.18) (1.03,-0.24) (1.04,-0.30) (1.06,-0.37) (1.08,-0.43) (1.11,-0.50) (1.14,-0.58) (1.17,-0.66) (1.21,-0.74) (1.24,-0.83) (1.29,-0.92) (1.33,-1.02) (1.38,-1.12) (1.44,-1.23) (1.49,-1.35) (1.55,-1.48) (1.61,-1.61) (1.68,-1.75) (1.75,-1.90) (1.82,-2.06) (1.90,-2.23)};
		\draw[color=red, thick] (-1.12,0.6204) -- (0.15,0.6204);
		\filldraw (-0.5774,0.6204) circle (1.8pt);
		\node[below] at (-0.5774,0.545) {$P$};
	\end{tikzpicture}
	\end{center}
	Here the tangent line touches the curve and stays on one side of it near $P$ -- compare the node, where no single line is tangent, and the cusp, where the tangent `line' degenerates entirely.
\end{example}

\subsection{Tangent Spaces in General}

For a general variety, cut out by several equations, we take the intersection of the tangent hyperplanes and record it as a vector subspace rather than an affine one.

\begin{definition}[Tangent Space]
	Let $V \subseteq \A^n$ be a variety and $P \in V$. The \vocab{tangent space} is
	$$T_{V,P} = \left\{\mathbf v \in \C^n : \sum_{i=1}^n v_i \frac{\partial f}{\partial X_i}(P) = 0 \text{ for all } f \in I(V)\right\} \leq \C^n.$$
\end{definition}

This is a linear subspace of $\C^n$, and it suffices to impose the conditions for $f$ running over a generating set of $I(V)$, by the product and sum rules for derivatives. Concretely, if $I(V) = (f_1,\dots,f_r)$ then $T_{V,P}$ is the kernel of the \vocab{Jacobian matrix}
$$J(P) = \left(\frac{\partial f_j}{\partial X_i}(P)\right)_{j,i},$$
so that
$$\dim T_{V,P} = n - \rk J(P).$$

\begin{example}
	Revisiting Example~\ref{ex:threelines}: for $V = \V(X_1X_2(X_1-X_2)) \subseteq \A^2$ we have $I(V) = (X_1X_2(X_1-X_2))$, whose gradient is a vector of quadratics, vanishing at the origin. So $T_{V,\mathbf 0} = \C^2$. For $W = \V(X_1X_2,X_2X_3,X_3X_1) \subseteq \A^3$ every generator has vanishing gradient at the origin, so $T_{W,\mathbf 0} = \C^3$. If $V$ and $W$ were isomorphic, some isomorphism would carry the origin to the origin (it is the unique point of each where the local structure is singular), and would induce a linear isomorphism of tangent spaces. Since $2 \neq 3$, $V \not\cong W$.
\end{example}

To transport tangent spaces along maps we need the derivative. Let $\varphi \colon V \dto W$ be a rational map of irreducible varieties, $P \in \dom\varphi$ and $Q = \varphi(P)$. Working in affine charts, $\varphi$ is given by rational functions $f_1,\dots,f_m$ regular at $P$, and we set
$$\dee\varphi_P(\mathbf v) = \left(\sum_{i=1}^n v_i \frac{\partial f_j}{\partial X_i}(P)\right)_{j=1}^m,$$
a linear map $\C^n \to \C^m$.

\begin{proposition}\label{prop:dvarphi}
	With the above notation:
	\begin{enumerate}[label=(\roman*)]
		\item $\dee\varphi_P(T_{V,P}) \subseteq T_{W,Q}$;
		\item $\dee\varphi_P$ depends only on $\varphi$, not on the chosen representatives;
		\item if $\psi \colon W \dto Z$ is rational with $Q \in \dom\psi$ then $\dee(\psi\varphi)_P = \dee\psi_Q \circ \dee\varphi_P$;
		\item if $\varphi$ is birational and $\varphi^{-1}$ is regular at $Q$, then $\dee\varphi_P \colon T_{V,P} \to T_{W,Q}$ is an isomorphism.
	\end{enumerate}
\end{proposition}

\begin{proof}
	(i) Let $g \in I(W)$, with coordinates $Y_j$ on $W$. Then $\varphi^\star g = g(f_1,\dots,f_m)$ vanishes identically on $V$ near $P$, so by the chain rule, for $\mathbf v \in T_{V,P}$,
	$$0 = \sum_i v_i \frac{\partial (\varphi^\star g)}{\partial X_i}(P) = \sum_j \frac{\partial g}{\partial Y_j}(Q)\left(\sum_i v_i \frac{\partial f_j}{\partial X_i}(P)\right),$$
	which says exactly that $\dee\varphi_P(\mathbf v) \in T_{W,Q}$.

	(ii) If $f_j'$ is another representative, then $f_j - f_j'$ vanishes on $V$ where defined, so $f_j - f_j' = p_j/q_j$ with $p_j \in I(V)$ and $q_j(P) \neq 0$. Differentiating by the quotient rule and using $p_j(P) = 0$,
	$$\frac{\partial (f_j - f_j')}{\partial X_i}(P) = \frac{1}{q_j(P)}\frac{\partial p_j}{\partial X_i}(P),$$
	and since $p_j \in I(V)$, contracting with $\mathbf v \in T_{V,P}$ gives zero.

	(iii) is the chain rule, and (iv) follows from (iii) applied to $\varphi^{-1}\varphi = \mathrm{id}$.
\end{proof}

For projective varieties we define $T_{V,P}$ by choosing an affine patch $U_j \ni P$ and setting $T_{V,P} = T_{V\cap U_j, P}$. Part (iv) of Proposition~\ref{prop:dvarphi}, applied to the transition maps between patches (which are birational and regular on the overlaps), shows this is independent of $j$ up to canonical isomorphism.

There is also a convenient projective form of the tangent space to a hypersurface.

\begin{definition}
	Let $V = \V(F) \subseteq \PP^n$ with $F$ homogeneous. The \vocab{projective tangent space} at $P \in V$ is
	$$T^{\mathrm{proj}}_{V,P} = \V\left(\sum_{i=0}^n \frac{\partial F}{\partial X_i}(P)\,X_i\right) \subseteq \PP^n.$$
\end{definition}

That $P$ itself lies in this hyperplane follows from \vocab{Euler's relation}: for $F$ homogeneous of degree $d$,
$$\sum_{i=0}^n X_i\frac{\partial F}{\partial X_i} = d\,F,$$
which is checked on monomials. Dehomogenising shows $T^{\mathrm{proj}}_{V,P} \cap U_0 = T^{\mathrm{aff}}_{V\cap U_0, P}$, so the two notions agree.

\begin{proposition}\label{prop:singclosed}
	Let $V = \V(F) \subseteq \PP^n$ be an irreducible hypersurface. Then the singular locus of $V$ is a proper Zariski closed subset; in particular the smooth points are dense.
\end{proposition}

\begin{proof}
	The singular locus is $V \cap \bigcap_i \V(\partial F/\partial X_i)$, which is closed. Suppose it were all of $V$. Then each $\partial F/\partial X_i$ vanishes on $V$, so lies in $I^h(V) = (F)$ (as $F$ is irreducible, $(F)$ is prime hence radical). But $\partial F/\partial X_i$ is homogeneous of degree $\deg F - 1 < \deg F$, so divisibility by $F$ forces $\partial F/\partial X_i = 0$ for all $i$. In characteristic zero this means $F$ is constant, so $V = \PP^n$ or $\varnothing$, contradicting that $V$ is a hypersurface.
\end{proof}

\subsection{Dimension}

We can now define dimension. The idea is that the tangent space has the `expected' dimension at most points, and jumps up at bad points.

\begin{definition}[Dimension]
	Let $V$ be an irreducible affine or projective variety. Its \vocab{dimension} is
	$$\dim V = \min_{P \in V} \dim T_{V,P}.$$
	A point $P \in V$ is \vocab{smooth} if $\dim T_{V,P} = \dim V$, and \vocab{singular} otherwise. For reducible $V$, we set $\dim V$ to be the maximum of the dimensions of the irreducible components.
\end{definition}

A variety of dimension $1$ is a \vocab{curve}, of dimension $2$ a \vocab{surface}.

\begin{theorem}\label{thm:smoothopen}
	Let $V$ be a nonempty irreducible variety. The set of smooth points of $V$ is a nonempty Zariski open subset, hence dense.
\end{theorem}

\begin{proof}
	Nonemptiness is by definition, since the minimum is attained. We may work in an affine chart, so assume $V \subseteq \A^n$ with $I(V) = (f_1,\dots,f_r)$. Then $\dim T_{V,P} = n - \rk J(P)$ where $J$ is the Jacobian. For any $s$,
	$$\{P \in V : \dim T_{V,P} \geq s\} = \{P \in V : \rk J(P) \leq n - s\},$$
	and the locus where a matrix has rank at most $k$ is cut out by the vanishing of all its $(k+1)\times(k+1)$ minors, which are polynomials in the entries, hence in $P$. So this set is closed. Taking $s = \dim V + 1$, the singular locus is closed and proper, so the smooth locus is open and nonempty; density follows from Proposition~\ref{prop:opendense}.
\end{proof}

\begin{corollary}
	If $V$ and $W$ are irreducible and birational, then $\dim V = \dim W$.
\end{corollary}

\begin{proof}
	A birational map restricts to an isomorphism between nonempty open subsets $U \subseteq V$ and $U' \subseteq W$. Shrinking, we may assume both consist of smooth points. Then by Proposition~\ref{prop:dvarphi}(iv) the tangent spaces are isomorphic at corresponding points, so the minima agree.
\end{proof}

\subsection{Dimension via Transcendence Degree}

The definition above is geometric; there is an equivalent algebraic description in terms of the function field, which is often much easier to compute with. Recall from Galois theory that if $K \subseteq L$ are fields, an element $\alpha \in L$ is \vocab{transcendental} over $K$ if it satisfies no nonzero polynomial over $K$, and a subset $S \subseteq L$ is \vocab{algebraically independent} over $K$ if there is no nonzero polynomial relation among its elements with coefficients in $K$.

\begin{definition}
	An extension $K/\C$ is \vocab{purely transcendental} if $K = \C(x_1,\dots,x_m)$ for some algebraically independent $x_1,\dots,x_m$.
\end{definition}

\begin{proposition}\label{prop:transbasis}
	Let $K/\C$ be a finitely generated field extension. Then there is a purely transcendental subextension $K_0 = \C(x_1,\dots,x_m) \subseteq K$ with $K/K_0$ finite. Moreover $K = K_0(y)$ for a single $y \in K$.
\end{proposition}

\begin{proof}
	Write $K = \C(x_1,\dots,x_n)$. Choose a maximal algebraically independent subset of $\{x_1,\dots,x_n\}$, which after relabelling is $\{x_1,\dots,x_m\}$; put $K_0 = \C(x_1,\dots,x_m)$. By maximality, each of $x_{m+1},\dots,x_n$ is algebraic over $K_0$, so $K/K_0$ is finitely generated and algebraic, hence finite. Since $\C$ has characteristic zero, $K/K_0$ is separable, and the Primitive Element Theorem from II Galois Theory gives $K = K_0(y)$.
\end{proof}

The integer $m$ is independent of choices; it is the \vocab{transcendence degree} $\trdeg_\C K$, and the argument is an exchange lemma identical to the one for bases of vector spaces.

\begin{proposition}\label{prop:hypersurfideal}
	Let $x_1,\dots,x_n$ be algebraically independent over $\C$, put $K = \C(x_1,\dots,x_n)$, and let $x_{n+1}$ be algebraic over $K$. Then
	$$I = \{g \in \C[X_1,\dots,X_{n+1}] : g(x_1,\dots,x_{n+1}) = 0\}$$
	is a principal ideal generated by an irreducible polynomial $f$.
\end{proposition}

\begin{proof}
	Since the $x_i$ are algebraically independent, $R = \C[x_1,\dots,x_n] \subseteq K$ is isomorphic to a polynomial ring, hence a UFD, with $\Frac R = K$. As $x_{n+1}$ is algebraic over $K$, it has a minimal polynomial $h \in K[T]$, monic and irreducible. Clearing denominators, let $b \in R$ be a common denominator of the coefficients of $h$ and set $f = bh \in R[T]$. By Gauss's Lemma, $f$ is irreducible in $R[T] \cong \C[X_1,\dots,X_{n+1}]$.

	If $g \in I$ then $g(x_1,\dots,x_n,T) \in K[T]$ has $x_{n+1}$ as a root, so is divisible by $h$ in $K[T]$; by Gauss's Lemma again, $f \mid g$ in $R[T]$. Hence $I = (f)$.
\end{proof}

\begin{corollary}\label{cor:birathyp}
	Every irreducible variety is birational to a hypersurface.
\end{corollary}

\begin{proof}
	Let $K = \C(V)$. By Proposition~\ref{prop:transbasis} we may write $K = \C(x_1,\dots,x_n,x_{n+1})$ with $x_1,\dots,x_n$ algebraically independent and $x_{n+1}$ algebraic over $\C(x_1,\dots,x_n)$. By Proposition~\ref{prop:hypersurfideal},
	$$\C[x_1,\dots,x_{n+1}] \cong \frac{\C[X_1,\dots,X_{n+1}]}{(f)}$$
	for an irreducible $f$, and $K$ is its field of fractions. So $\C(V) \cong \C(\V(f))$ for the hypersurface $\V(f) \subseteq \A^{n+1}$, and Theorem~\ref{thm:biratfield} gives the result.
\end{proof}

\begin{corollary}\label{cor:dimtrdeg}
	For an irreducible variety $V$ we have $\dim V = \trdeg_\C \C(V)$. In particular an irreducible hypersurface in $\A^n$ or $\PP^n$ has dimension $n-1$, and $\dim \A^n = \dim\PP^n = n$.
\end{corollary}

\begin{proof}
	Both sides are birational invariants -- the left by the corollary above, the right because birational varieties have isomorphic function fields. By Corollary~\ref{cor:birathyp} we may assume $V = \V(f) \subseteq \A^{n}$ is an irreducible hypersurface with $f$ non-constant, say involving $X_n$. Then $\C(V) = \C(x_1,\dots,x_n)$ where $x_1,\dots,x_{n-1}$ are algebraically independent and $x_n$ is algebraic over them, so $\trdeg = n-1$. On the other hand some partial derivative of $f$ is not identically zero on $V$ (Proposition~\ref{prop:singclosed}), so at a general point the Jacobian has rank $1$ and $\dim T_{V,P} = n-1$. Hence both sides equal $n-1$.
\end{proof}

This is the definition of dimension that we will actually use: \emph{the dimension of an irreducible variety is the transcendence degree of its function field}. In particular a \vocab{curve} is an irreducible variety whose function field is a finite extension of $\C(t)$ for some transcendental $t$, and this is the description that will drive everything in the second half of these notes.

\section{Plane Curves and B\'ezout's Theorem}

We now specialise to curves in $\PP^2$, where we can be extremely concrete. The organising principle is B\'ezout's theorem: two plane curves of degrees $m$ and $n$ meet in exactly $mn$ points, provided we work projectively, over $\C$, and count with multiplicity. Each of these three provisos is essential, and each corresponds to something we have set up.

\subsection{Plane Curves and their Degree}

\begin{definition}[Plane Curve]
	A \vocab{projective plane curve} of \vocab{degree} $d$ is $C = \V(F) \subseteq \PP^2$ where $F \in \C[X_0,X_1,X_2]$ is homogeneous of degree $d$ with no repeated irreducible factors. Curves of degree $1,2,3,4$ are called \vocab{lines}, \vocab{conics}, \vocab{cubics} and \vocab{quartics}.
\end{definition}

Requiring no repeated factors makes $F$ recoverable from $C$: by the Nullstellensatz $I^h(C) = \sqrt{(F)} = (F_{\mathrm{red}})$, so the degree is well defined. $C$ is irreducible if and only if $F$ is.

\begin{example}
	If $C = \V(f) \subseteq \A^2$ is an affine plane curve of degree $d$, its projective closure is $\V(f^h)$, a projective plane curve of degree $d$. The points at infinity are found by setting $X_0 = 0$: they are the zeros of the top-degree part $f_{[d]}$ of $f$, homogenised. For instance the affine cubic $Y^2 = X^3 - X$ has $f^h = X_0X_2^2 - X_1^3 + X_1X_0^2$, and setting $X_0 = 0$ gives $X_1^3 = 0$, so there is exactly one point at infinity, $(0:0:1)$, and it is a triple point of the intersection with the line at infinity. This point is where the two ends of the real picture in Section~1 join up.
\end{example}

\subsection{Intersection Multiplicity}

The naive count of intersection points is not constant: a line meets a conic in two points usually, but in only one if it is tangent. To get a constant answer we must count tangential intersections twice.

\begin{definition}[Intersection Multiplicity with a Line]
	Let $C = \V(F)$ be a plane curve of degree $d$ and $L$ a line, $L \not\subseteq C$. Parametrise $L$ by $s \mapsto \mathbf a + s\mathbf b$ (in homogeneous coordinates: $(s:t)\mapsto s\mathbf a + t \mathbf b$). Then $G(s,t) = F(s\mathbf a + t\mathbf b)$ is a nonzero binary form of degree $d$, so factors as
	$$G(s,t) = c\prod_{k}(\beta_k s - \alpha_k t)^{m_k}$$
	with $(\alpha_k:\beta_k)$ distinct. The \vocab{intersection multiplicity} of $C$ and $L$ at the point $P_k$ corresponding to $(\alpha_k : \beta_k)$ is $I_{P_k}(C,L) = m_k$, and $I_P(C,L) = 0$ for $P \notin C \cap L$.
\end{definition}

Since a binary form of degree $d$ over $\C$ factors into $d$ linear factors, we get for free:

\begin{proposition}[B\'ezout for a line]\label{prop:bezoutline}
	Let $C$ be a plane curve of degree $d$ and $L$ a line with $L \not\subseteq C$. Then
	$$\sum_{P \in C\cap L} I_P(C,L) = d.$$
	In particular $|C \cap L| \leq d$, with equality for all but finitely many lines.
\end{proposition}

\begin{proof}
	The displayed identity is the factorisation of $G$ above, as $\deg G = d$ (the coefficient of $s^d$ is $F(\mathbf a)$, which is nonzero for a suitable choice of $\mathbf a \in L$ unless $L \subseteq C$).
\end{proof}

One checks directly that $I_P(C,L) \geq 2$ if and only if $L$ is the tangent line to $C$ at $P$, or $P$ is a singular point of $C$; this is exactly the computation with which we opened Section~\ref{sec:tangent}.

\begin{definition}[Intersection Multiplicity, general]
	For plane curves $C = \V(F)$ and $D = \V(G)$ with no common component and $P \in \PP^2$, the \vocab{intersection multiplicity} is
	$$I_P(C,D) = \dim_\C \frac{\Oh_{\PP^2,P}}{(f,g)},$$
	where $f,g$ are the dehomogenisations of $F,G$ in an affine chart containing $P$.
\end{definition}

This is finite and independent of the chart, and it agrees with the multiplicity defined above when $D$ is a line. We will not develop its properties in detail -- the local algebra required is beyond this course -- but it is worth recording the characterisation that makes it usable: $I_P(C,D)$ is the unique assignment of non-negative integers such that
\begin{enumerate}[label=(\roman*)]
	\item $I_P(C,D) = I_P(D,C)$, and $I_P(C,D) = 0$ if and only if $P \notin C\cap D$;
	\item $I_P(C,D) = 1$ if and only if $P$ is a smooth point of both curves and their tangent lines at $P$ differ (a \vocab{transverse} intersection);
	\item $I_P$ is additive: $I_P(C, D\cup D') = I_P(C,D) + I_P(C,D')$;
	\item $I_P(C,D) = I_P(C, D')$ whenever $G' = G + HF$ for some form $H$ of the appropriate degree.
\end{enumerate}

\subsection{B\'ezout's Theorem}

\begin{theorem}[B\'ezout]\label{thm:bezout}
	Let $C, D \subseteq \PP^2$ be plane curves of degrees $m$ and $n$ with no common component. Then $C\cap D$ is finite and
	$$\sum_{P \in C\cap D} I_P(C,D) = mn.$$
	In particular $|C\cap D| \leq mn$, with equality if and only if all intersections are transverse.
\end{theorem}

We will prove the weak form -- the inequality $|C\cap D| \leq mn$ -- in Section~\ref{sec:divisors} using divisors, once we have the machinery of orders of vanishing. The strong form requires the full theory of intersection multiplicity. For now let us see what it says.

\begin{example}
	A line and a conic meet in $2$ points, a line and a cubic in $3$, and two conics in $4$. Here is the last of these: the circle $X^2+Y^2 = 4$ and the ellipse $X^2/9 + Y^2 = 1$ meet in the four points $(\pm\sqrt{27}/\!\sqrt{8},\pm\sqrt{5/8}) \approx (\pm 1.837, \pm 0.791)$.
	\begin{center}
	\begin{tikzpicture}[x=1cm,y=1cm]
		\draw[<->] (-3.6,0) -- (3.6,0);
		\draw[<->] (0,-2.5) -- (0,2.5);
		\draw[thick, color=blue] (0,0) circle (2);
		\draw[thick, color=blue] (0,0) ellipse (3 and 1);
		\foreach \x/\y in {1.837/0.791, -1.837/0.791, 1.837/-0.791, -1.837/-0.791}{
			\filldraw (\x,\y) circle (1.7pt);
		}
		\node[color=blue, above right] at (1.35,1.55) {$C$};
		\node[color=blue, below] at (2.95,-0.15) {$D$};
	\end{tikzpicture}
	\end{center}
\end{example}

\begin{example}
	The line $Y = 0.55X + 0.45$ meets the cubic $Y^2 = X^3 - X$ in the three points with $X \approx 1.439, -0.995, -0.141$. All three are real here, and they are visible in the picture.
	\begin{center}
	\begin{tikzpicture}[x=1.5cm,y=1.05cm]
		\draw[ultra thin, color=lightgray] (-1.4,-2.3) grid[step=0.5] (2.1,2.3);
		\draw[<->] (-1.5,0) -- (2.2,0);
		\draw[<->] (0,-2.4) -- (0,2.4);
		\draw[thick, color=blue, smooth] plot coordinates {(0.00,0.00) (-0.00,0.06) (-0.02,0.13) (-0.04,0.19) (-0.06,0.25) (-0.10,0.31) (-0.14,0.36) (-0.18,0.42) (-0.23,0.47) (-0.29,0.51) (-0.35,0.55) (-0.41,0.58) (-0.47,0.60) (-0.53,0.62) (-0.59,0.62) (-0.65,0.61) (-0.71,0.59) (-0.77,0.56) (-0.82,0.52) (-0.86,0.47) (-0.90,0.41) (-0.94,0.34) (-0.96,0.26) (-0.98,0.18) (-1.00,0.09) (-1.00,0.00) (-1.00,-0.09) (-0.98,-0.18) (-0.96,-0.26) (-0.94,-0.34) (-0.90,-0.41) (-0.86,-0.47) (-0.82,-0.52) (-0.77,-0.56) (-0.71,-0.59) (-0.65,-0.61) (-0.59,-0.62) (-0.53,-0.62) (-0.47,-0.60) (-0.41,-0.58) (-0.35,-0.55) (-0.29,-0.51) (-0.23,-0.47) (-0.18,-0.42) (-0.14,-0.36) (-0.10,-0.31) (-0.06,-0.25) (-0.04,-0.19) (-0.02,-0.13) (-0.00,-0.06) (0.00,0.00)};
		\draw[thick, color=blue, smooth] plot coordinates {(1.90,2.23) (1.82,2.06) (1.75,1.90) (1.68,1.75) (1.61,1.61) (1.55,1.48) (1.49,1.35) (1.44,1.23) (1.38,1.12) (1.33,1.02) (1.29,0.92) (1.24,0.83) (1.21,0.74) (1.17,0.66) (1.14,0.58) (1.11,0.50) (1.08,0.43) (1.06,0.37) (1.04,0.30) (1.03,0.24) (1.02,0.18) (1.01,0.12) (1.00,0.06) (1.00,0.00) (1.00,-0.06) (1.01,-0.12) (1.02,-0.18) (1.03,-0.24) (1.04,-0.30) (1.06,-0.37) (1.08,-0.43) (1.11,-0.50) (1.14,-0.58) (1.17,-0.66) (1.21,-0.74) (1.24,-0.83) (1.29,-0.92) (1.33,-1.02) (1.38,-1.12) (1.44,-1.23) (1.49,-1.35) (1.55,-1.48) (1.61,-1.61) (1.68,-1.75) (1.75,-1.90) (1.82,-2.06) (1.90,-2.23)};
		\draw[color=red, thick] (-1.35,-0.2925) -- (2.05,1.5775);
		\foreach \x/\y in {1.439/1.2415, -0.995/-0.0974, -0.141/0.3722}{
			\filldraw (\x,\y) circle (1.7pt);
		}
	\end{tikzpicture}
	\end{center}
\end{example}

\begin{corollary}[Five points determine a conic]
	Let $P_1,\dots,P_5 \in \PP^2$ be five points, no three collinear. Then there is a unique conic through all five, and it is smooth.
\end{corollary}

\begin{proof}
	A conic is $\V(F)$ for a homogeneous quadratic $F$ in three variables; the space of such $F$ is a $6$-dimensional vector space, and $F(P_i) = 0$ is a linear condition. Five linear conditions on a $6$-dimensional space have a nonzero solution, so a conic through the five points exists.

	If two distinct conics $C, D$ passed through all five points, then by B\'ezout they share a component (else $|C\cap D| \leq 4 < 5$). A shared component is a line or the whole conic; the latter is excluded as $C \neq D$, so $C = L \cup L'$ and $D = L \cup L''$ with $L$ a common line. Then at least three of the five points lie on one line, contradiction. So the conic is unique.

	Finally, a singular conic is a union of two lines (or a double line), which again would force three collinear points. So the conic is smooth.
\end{proof}

\subsection{Inflection Points}

\begin{definition}[Flex]
	Let $C \subseteq \PP^2$ be a plane curve and $P$ a smooth point with tangent line $L$. Then $P$ is a \vocab{point of inflection}, or \vocab{flex}, if $I_P(C,L) \geq 3$.
\end{definition}

So a flex is a point where the tangent line meets the curve to order at least three -- it `crosses over' rather than touching. For a curve of degree $2$ there are no flexes, since $I_P(C,L) \leq 2$ by B\'ezout with a line. For cubics, flexes are exactly what we need to normalise the equation.

\begin{example}
	On $C \colon Y^2 = X^3 - X + 1$ there is a real flex at $P \approx (0.0871, 0.9558)$, with tangent line of slope $\approx -0.5112$. Substituting the tangent into the cubic gives a polynomial with a triple root at $x = 0.0871$.
	\begin{center}
	\begin{tikzpicture}[x=1.35cm,y=0.95cm]
		\draw[ultra thin, color=lightgray] (-1.6,-2.8) grid[step=0.5] (2.1,2.8);
		\draw[<->] (-1.7,0) -- (2.2,0);
		\draw[<->] (0,-2.9) -- (0,2.9);
		\draw[thick, color=blue, smooth] plot coordinates {(-1.32,0.00) (-1.32,0.13) (-1.31,0.26) (-1.29,0.38) (-1.26,0.50) (-1.23,0.62) (-1.18,0.73) (-1.13,0.83) (-1.07,0.92) (-1.00,1.00) (-0.93,1.06) (-0.85,1.11) (-0.76,1.15) (-0.66,1.17) (-0.55,1.18) (-0.44,1.16) (-0.31,1.13) (-0.18,1.08) (-0.04,1.02) (0.10,0.95) (0.26,0.87) (0.42,0.81) (0.59,0.78) (0.77,0.83) (0.95,0.95) (1.15,1.17) (1.35,1.45) (1.56,1.79) (1.77,2.19) (2.00,2.65)};
		\draw[thick, color=blue, smooth] plot coordinates {(2.00,-2.65) (1.77,-2.19) (1.56,-1.79) (1.35,-1.45) (1.15,-1.17) (0.95,-0.95) (0.77,-0.83) (0.59,-0.78) (0.42,-0.81) (0.26,-0.87) (0.10,-0.95) (-0.04,-1.02) (-0.18,-1.08) (-0.31,-1.13) (-0.44,-1.16) (-0.55,-1.18) (-0.66,-1.17) (-0.76,-1.15) (-0.85,-1.11) (-0.93,-1.06) (-1.00,-1.00) (-1.07,-0.92) (-1.13,-0.83) (-1.18,-0.73) (-1.23,-0.62) (-1.26,-0.50) (-1.29,-0.38) (-1.31,-0.26) (-1.32,-0.13) (-1.32,-0.00)};
		\draw[color=red, thick] (-1.5,1.767) -- (1.9,0.029);
		\filldraw (0.0871,0.9558) circle (1.8pt);
		\node[below right] at (0.15,0.87) {$P$};
	\end{tikzpicture}
	\end{center}
	Note that near $P$ the tangent line crosses from one side of the curve to the other.
\end{example}

There is a beautiful way of detecting flexes algebraically.

\begin{definition}[Hessian]
	Let $F \in \C[X_0,X_1,X_2]$ be homogeneous of degree $d$. Its \vocab{Hessian} is
	$$H_F = \det\left(\frac{\partial^2 F}{\partial X_i \partial X_j}\right)_{0\leq i,j \leq 2},$$
	a homogeneous polynomial of degree $3(d-2)$.
\end{definition}

\begin{theorem}\label{thm:hessian}
	Let $C = \V(F) \subseteq \PP^2$ be a smooth curve of degree $d \geq 2$ and $P \in C$. Then $P$ is a flex of $C$ if and only if $H_F(P) = 0$. Consequently, a smooth plane curve of degree $d \geq 3$ has at least one flex, and at most $3d(d-2)$ of them; a smooth cubic has exactly $9$.
\end{theorem}

\begin{proof}[Proof sketch]
	Move $P$ to $(1:0:0)$ and arrange that the tangent line there is $\V(X_2)$. Dehomogenising in the chart $X_0 = 1$ with coordinates $(x,y)$, the curve is $y = a x^2 + (\text{higher order})$ near the origin, and $P$ is a flex precisely when $a = 0$. A direct computation using Euler's relation shows that $H_F(P)$ is a nonzero multiple of $a$, which gives the equivalence.

	For the count: $H_F$ has degree $3(d-2)$, so if $C \not\subseteq \V(H_F)$ then B\'ezout gives $|C \cap \V(H_F)| \leq 3d(d-2)$, with equality when all intersections are transverse. That $C \not\subseteq \V(H_F)$ for smooth $C$ (in characteristic zero) is a separate check. For $d = 3$ one shows all the intersections are transverse, giving exactly $9$ flexes.
\end{proof}

\subsection{Weierstrass Form}

We can now put any smooth plane cubic into a standard shape. This is the calculation that connects the abstract theory to the concrete equations you have seen in number theory.

\begin{theorem}[Weierstrass Normal Form]\label{thm:weierstrass}
	Let $C \subseteq \PP^2$ be a smooth plane cubic. After a projective change of coordinates,
	$$C = \V\big(X_0X_2^2 - X_1^3 - aX_0^2X_1 - bX_0^3\big)$$
	for some $a,b \in \C$ with $4a^3 + 27b^2 \neq 0$. Equivalently, in the affine chart $X_0 = 1$,
	$$C \colon\quad y^2 = x^3 + ax + b = (x-\lambda_1)(x-\lambda_2)(x-\lambda_3)$$
	with $\lambda_1,\lambda_2,\lambda_3$ distinct, together with the single point $(0:0:1)$ at infinity.
\end{theorem}

\begin{proof}
	By Theorem~\ref{thm:hessian}, $C$ has a flex $P_0$. Change coordinates so that $P_0 = (0:0:1)$ and the tangent line to $C$ at $P_0$ is the line at infinity $\V(X_0)$. Write $F$ for the defining cubic. Since $P_0$ is a flex, $I_{P_0}(C,\V(X_0)) = 3$, so by Proposition~\ref{prop:bezoutline} the line at infinity meets $C$ only at $P_0$; restricting $F$ to $X_0 = 0$ therefore gives (up to scalar) $X_1^3$. So $F$ has no $X_1^2X_2$, $X_1X_2^2$ or $X_2^3$ terms except via $X_0$, and after scaling
	$$F = cX_0X_2^2 + X_0 X_2 \ell(X_0,X_1) - X_1^3 + X_0 q(X_0,X_1)$$
	for a linear form $\ell$ and quadratic $q$. Smoothness at $P_0 = (0:0:1)$ forces $c \neq 0$ (else $\partial F/\partial X_0$, $\partial F/\partial X_1$, $\partial F/\partial X_2$ all vanish there), so we may scale $c = 1$. Completing the square in $X_2$ -- that is, replacing $X_2$ by $X_2 - \tfrac12 \ell(X_0,X_1)$, a linear change of coordinates -- removes the middle term, leaving
	$$F = X_0X_2^2 - X_1^3 + X_0\tilde q(X_0,X_1)$$
	for a new quadratic $\tilde q$. Finally, replacing $X_1$ by $X_1 - \tfrac13(\text{coefficient of } X_1^2 \text{ in } \tilde q)X_0$ kills the $X_0X_1^2$ term, giving the stated form.

	It remains to see that smoothness is equivalent to $x^3+ax+b$ having distinct roots. In the affine chart, $f = y^2 - x^3-ax-b$ has $\partial f/\partial y = 2y$ and $\partial f/\partial x = -(3x^2+a)$. A singular point has $y = 0$ and $x$ a common root of $x^3+ax+b$ and its derivative, i.e.\ a repeated root; and $4a^3+27b^2$ is (up to sign) the discriminant of $x^3+ax+b$, vanishing exactly when there is a repeated root.
\end{proof}

We call the pair $(C,P_0)$, where $P_0 = (0:0:1)$ is the flex at infinity, an \vocab{elliptic curve}. (This is not quite the general definition -- a genuinely intrinsic one will be given in Section~\ref{sec:lowgenus} -- but for plane cubics it is equivalent.)

\subsection{The Group Law on a Cubic}

Here is the most striking fact about plane cubics: the points of a smooth cubic form an abelian group, and the group operation is given by drawing lines.

By B\'ezout, a line $L$ meets a smooth cubic $C$ in exactly three points counted with multiplicity. Given $P,Q \in C$, let $L$ be the line through them (the tangent line at $P$ if $P = Q$), and let $P * Q$ denote the third point of $C \cap L$.

\begin{definition}[Chord-and-Tangent Law]\label{def:chordtangent}
	Let $(C,P_0)$ be a smooth plane cubic with a chosen flex $P_0$. For $P,Q \in C$ define
	$$P + Q = P_0 * (P*Q).$$
\end{definition}

In Weierstrass form, with $P_0 = (0:0:1)$ the point at infinity, the line through $P_0$ and any point $R = (x,y)$ is the vertical line through $R$, so $P_0 * R = (x,-y)$. Thus the recipe is: \emph{draw the line through $P$ and $Q$, find the third intersection with $C$, and reflect in the $x$-axis.}

\begin{center}
\begin{tikzpicture}[x=1.35cm,y=0.9cm]
	\draw[<->] (-1.7,0) -- (2.25,0);
	\draw[<->] (0,-2.9) -- (0,2.9);
	\draw[thick, color=blue, smooth] plot coordinates {(-1.32,0.00) (-1.32,0.13) (-1.31,0.26) (-1.29,0.38) (-1.26,0.50) (-1.23,0.62) (-1.18,0.73) (-1.13,0.83) (-1.07,0.92) (-1.00,1.00) (-0.93,1.06) (-0.85,1.11) (-0.76,1.15) (-0.66,1.17) (-0.55,1.18) (-0.44,1.16) (-0.31,1.13) (-0.18,1.08) (-0.04,1.02) (0.10,0.95) (0.26,0.87) (0.42,0.81) (0.59,0.78) (0.77,0.83) (0.95,0.95) (1.15,1.17) (1.35,1.45) (1.56,1.79) (1.77,2.19) (2.00,2.65)};
	\draw[thick, color=blue, smooth] plot coordinates {(2.00,-2.65) (1.77,-2.19) (1.56,-1.79) (1.35,-1.45) (1.15,-1.17) (0.95,-0.95) (0.77,-0.83) (0.59,-0.78) (0.42,-0.81) (0.26,-0.87) (0.10,-0.95) (-0.04,-1.02) (-0.18,-1.08) (-0.31,-1.13) (-0.44,-1.16) (-0.55,-1.18) (-0.66,-1.17) (-0.76,-1.15) (-0.85,-1.11) (-0.93,-1.06) (-1.00,-1.00) (-1.07,-0.92) (-1.13,-0.83) (-1.18,-0.73) (-1.23,-0.62) (-1.26,-0.50) (-1.29,-0.38) (-1.31,-0.26) (-1.32,-0.13) (-1.32,-0.00)};
	\draw[color=red, thick] (-1.35,-1.386) -- (1.80,2.091);
	\draw[color=DarkGreen, thick, densely dashed] (0.618,0.786) -- (0.618,-0.786);
	\filldraw (-1.0,-1.0) circle (1.8pt) node[left] {$P$};
	\filldraw (1.6,1.870) circle (1.8pt) node[right] {$Q$};
	\filldraw (0.618,0.786) circle (1.8pt);
	\node[below right] at (0.66,0.78) {$P*Q$};
	\filldraw (0.618,-0.786) circle (1.8pt);
	\node[below right] at (0.66,-0.82) {$P+Q$};
\end{tikzpicture}
\end{center}

\begin{center}
\begin{tikzpicture}[x=1.35cm,y=0.9cm]
	\draw[<->] (-1.7,0) -- (2.25,0);
	\draw[<->] (0,-2.9) -- (0,2.9);
	\draw[thick, color=blue, smooth] plot coordinates {(-1.32,0.00) (-1.32,0.13) (-1.31,0.26) (-1.29,0.38) (-1.26,0.50) (-1.23,0.62) (-1.18,0.73) (-1.13,0.83) (-1.07,0.92) (-1.00,1.00) (-0.93,1.06) (-0.85,1.11) (-0.76,1.15) (-0.66,1.17) (-0.55,1.18) (-0.44,1.16) (-0.31,1.13) (-0.18,1.08) (-0.04,1.02) (0.10,0.95) (0.26,0.87) (0.42,0.81) (0.59,0.78) (0.77,0.83) (0.95,0.95) (1.15,1.17) (1.35,1.45) (1.56,1.79) (1.77,2.19) (2.00,2.65)};
	\draw[thick, color=blue, smooth] plot coordinates {(2.00,-2.65) (1.77,-2.19) (1.56,-1.79) (1.35,-1.45) (1.15,-1.17) (0.95,-0.95) (0.77,-0.83) (0.59,-0.78) (0.42,-0.81) (0.26,-0.87) (0.10,-0.95) (-0.04,-1.02) (-0.18,-1.08) (-0.31,-1.13) (-0.44,-1.16) (-0.55,-1.18) (-0.66,-1.17) (-0.76,-1.15) (-0.85,-1.11) (-0.93,-1.06) (-1.00,-1.00) (-1.07,-0.92) (-1.13,-0.83) (-1.18,-0.73) (-1.23,-0.62) (-1.26,-0.50) (-1.29,-0.38) (-1.31,-0.26) (-1.32,-0.13) (-1.32,-0.00)};
	\draw[color=red, thick] (-1.35,-1.35) -- (1.85,1.85);
	\draw[color=DarkGreen, thick, densely dashed] (-1.0,-1.0) -- (-1.0,1.0);
	\filldraw (1.0,1.0) circle (1.8pt) node[right] {$T$};
	\filldraw (-1.0,-1.0) circle (1.8pt);
	\node[below right] at (-0.96,-1.05) {$T*T$};
	\filldraw (-1.0,1.0) circle (1.8pt);
	\node[above right] at (-0.96,1.05) {$2T$};
\end{tikzpicture}
\end{center}

The second picture shows the doubling: on $y^2 = x^3-x+1$ the tangent at $T = (1,1)$ has slope $1$ and meets the curve again at $(-1,-1)$, so $2T = (-1,1)$.

\begin{theorem}\label{thm:grouplaw}
	Let $(C,P_0)$ be a smooth plane cubic with $P_0$ a flex. Then $(C,+)$ is an abelian group with identity $P_0$, and $-P = P_0 * P$. Moreover $P + Q + R = P_0$ if and only if $P, Q, R$ are collinear.
\end{theorem}

\begin{proof}[Proof sketch]
	Commutativity is clear, since $P*Q = Q*P$. The identity: $P + P_0 = P_0*(P*P_0)$, and $P_0 * (P*P_0) = P$ because the line through $P_0$ and $P*P_0$ is the line through $P_0$ and $P$, whose third point is $P$. Inverses: set $-P = P_0*P$, so $P + (-P) = P_0*(P * (P_0*P)) = P_0 * P_0 = P_0$, the last equality because $P_0$ is a flex, so the tangent at $P_0$ meets $C$ only at $P_0$.

	The collinearity statement follows: if $P,Q,R$ are collinear then $R = P*Q$, so $P+Q = P_0*R = -R$ and $P+Q+R = P_0$.

	Associativity is the only difficulty. It can be proved by a direct (and unpleasant) computation with the Weierstrass equation, or by an application of B\'ezout's theorem for cubics -- if two cubics meet in nine points and a third cubic passes through eight of them, it passes through the ninth -- or, most cleanly, by identifying $(C,+)$ with the divisor class group, which we do in Theorem~\ref{thm:elliptic}.
\end{proof}

The last of these is the argument we will actually give, and it explains \emph{why} there is a group law: the points of $C$ are in bijection with the degree zero divisor classes, and those obviously form a group. Everything geometric in the chord-and-tangent construction is the statement that three points of $C$ sum to zero in the class group exactly when they are cut out by a line.

\section{Curves and their Local Rings}

From now on we specialise to curves, where the theory becomes both deeper and more complete than in higher dimensions. Recall that a curve is an irreducible variety of dimension $1$; equivalently, by Corollary~\ref{cor:dimtrdeg}, one whose function field has transcendence degree $1$ over $\C$.

\begin{notation}
	Unless stated otherwise, from now on \emph{curve} means a smooth irreducible projective curve. We will always have in mind some embedding $C \subseteq \PP^n$, but the ambient space will play no essential role: everything we do is intrinsic.
\end{notation}

\subsection{Closed Subsets of Curves}

The first thing to establish is that curves really are one-dimensional in the naive sense.

\begin{proposition}\label{prop:closedfinite}
	Let $C$ be an irreducible curve and $D \subsetneq C$ a proper Zariski closed subset. Then $D$ is a finite set of points.
\end{proposition}

\begin{proof}
	Decomposing into irreducible components, it suffices to show that a proper irreducible closed $W \subsetneq V$ is a point, where we may pass to an affine chart and take $V \subseteq \A^n$ affine irreducible with $\trdeg_\C \C(V) = 1$.

	The inclusion $\iota \colon W \hookrightarrow V$ gives a surjection $\iota^\star \colon \C[V] \twoheadrightarrow \C[W]$ with kernel $I(W)/I(V) \neq 0$, since $I(V) \subsetneq I(W)$ strictly by the Nullstellensatz. Choose $x \in \C[V]$ nonzero with $\iota^\star(x) = 0$.

	Suppose for contradiction that $\C[W] \neq \C$, and pick $t \in \C[W]\setminus\C$; choose $y \in \C[V]$ with $\iota^\star(y) = t$. I claim $x$ and $y$ are algebraically independent over $\C$. If not, there is a nonzero $g \in \C[S,T]$ with $g(x,y) = 0$; choose such a $g$ of minimal degree in $S$. Applying $\iota^\star$, $g(0,t) = 0$ in $\C[W]$. Now $t$ is transcendental over $\C$ (as $\C$ is algebraically closed and $\C[W]$ is a domain containing $\C$ properly, so any algebraic element would lie in $\C$), so $g(0,T)$ is the zero polynomial, i.e.\ $S \mid g$. Then $g = Sg_1$ and $xg_1(x,y) = 0$; as $\C[V]$ is a domain and $x \neq 0$, $g_1(x,y) = 0$, contradicting minimality.

	So $\C(V)$ contains two algebraically independent elements, contradicting $\trdeg_\C\C(V) = 1$. Hence $\C[W] = \C$, so $W$ is a single point.
\end{proof}

So the Zariski topology on a curve is the cofinite topology, exactly as on $\A^1$. The interesting structure is all in the local rings.

\subsection{Local Parameters}

\begin{theorem}\label{thm:localparam}
	Let $P$ be a smooth point of an irreducible curve $C$. Then the maximal ideal $\mm_P \trianglelefteq \Oh_{C,P}$ is principal.
\end{theorem}

Any generator $\pi_P$ of $\mm_P$ is called a \vocab{local parameter}, \vocab{local coordinate} or \vocab{uniformiser} at $P$. It is unique up to multiplication by a unit of $\Oh_{C,P}$.

The proof rests on the following piece of commutative algebra, which one should think of as an algebraic substitute for the implicit function theorem.

\begin{lemma}[Nakayama]\label{lem:nakayama}
	Let $R$ be a local ring with maximal ideal $\mm$, and let $M$ be a finitely generated $R$-module.
	\begin{enumerate}[label=(\roman*)]
		\item If $\mm M = M$ then $M = 0$.
		\item If $N \leq M$ is a submodule with $\mm M + N = M$, then $M = N$.
	\end{enumerate}
\end{lemma}

\begin{proof}
	(i) Suppose $M \neq 0$ and let $b_1,\dots,b_k$ be a generating set of minimal size. Since $M = \mm M$, we may write $b_k = \sum_j a_j b_j$ with $a_j \in \mm$. Then $(1-a_k)b_k = \sum_{j<k} a_jb_j$. As $a_k \in \mm$ and $R$ is local, $1 - a_k$ is a unit, so $b_k$ lies in the span of $b_1,\dots,b_{k-1}$, contradicting minimality.

	(ii) Apply (i) to $M/N$, noting $\mm(M/N) = (\mm M + N)/N = M/N$.
\end{proof}

\begin{proof}[Proof of Theorem~\ref{thm:localparam}]
	Pass to an affine chart and change coordinates so that $P = \mathbf 0 \in \A^n$. Write $\C[C] = \C[x_1,\dots,x_n]$ for the coordinate ring, where $x_i$ is the image of $X_i$. Then
	$$\Oh_P = \left\{\tfrac fg : f,g \in \C[C],\ g(P) \neq 0\right\}, \qquad \mm_P = x_1\Oh_P + \cdots + x_n\Oh_P.$$
	Note that every ideal $J \trianglelefteq \Oh_P$ is finitely generated: $J$ is determined by $J \cap \C[C]$, which is finitely generated as $\C[C]$ is Noetherian.

	Since $P$ is smooth and $\dim C = 1$, the tangent space $T_{C,P}$ is a line; after a further linear change of coordinates we may assume $T_{C,P} = \{X_2 = \cdots = X_n = 0\}$. The tangent space is cut out by the linear parts at $P$ of the elements of $I(C)$, so for each $j \geq 2$ there is $f_j \in I(C)$ whose linear part is $X_j$, say
	$$f_j = X_j - h_j(X_1,\dots,X_n)$$
	with $h_j$ having no terms of degree $< 2$. Passing to $\C[C]$, where $f_j = 0$,
	$$x_j = h_j(x_1,\dots,x_n) \in \mm_P^2 \qquad (j \geq 2).$$
	Hence
	$$\mm_P = \sum_{i=1}^n x_i\Oh_P \subseteq x_1\Oh_P + \mm_P^2 \subseteq \mm_P.$$
	Now apply Nakayama's lemma with $R = \Oh_P$, $M = \mm_P$ (finitely generated), $N = x_1\Oh_P$, noting $\mm_P M = \mm_P^2$: we get $\mm_P = x_1\Oh_P$, so $\pi_P = x_1$ is a local parameter.
\end{proof}

\begin{corollary}\label{cor:localparamplane}
	Let $C = \V(f) \subseteq \A^2$ be an irreducible plane curve and $P \in C$ a smooth point. Then the function $Q \mapsto X(Q) - X(P)$ is a local parameter at $P$ if and only if $\frac{\partial f}{\partial Y}(P) \neq 0$; and $Q \mapsto Y(Q) - Y(P)$ is one if and only if $\frac{\partial f}{\partial X}(P) \neq 0$.
\end{corollary}

\begin{proof}
	This is the proof above: if $\partial f/\partial Y (P) \neq 0$ then the tangent line is not vertical, so after translating $P$ to the origin the roles of $x_1$ and $x_2$ in the argument are as stated.
\end{proof}

Note that by smoothness at least one of the two partial derivatives is nonzero, so at least one of $x - x(P)$, $y - y(P)$ is always a local parameter. This will be used constantly in computations with plane curves.

\subsection{Orders of Vanishing}

Once $\mm_P$ is principal, every element of the function field has a well-defined order of vanishing at $P$, and this is the single most useful tool we have.

\begin{corollary}[Valuation]\label{cor:valuation}
	Let $P$ be a smooth point of a curve $C$. There is a surjective group homomorphism
	$$\nu_P \colon \C(C)^\times \to \Z$$
	such that
	\begin{enumerate}[label=(\roman*)]
		\item $\Oh_{C,P} = \{0\} \cup \{f \in \C(C)^\times : \nu_P(f) \geq 0\}$;
		\item $\mm_P = \{0\} \cup \{f \in \C(C)^\times : \nu_P(f) > 0\}$;
		\item for any local parameter $\pi_P$ and any $f \in \C(C)^\times$ we may write $f = \pi_P^{\nu_P(f)}u$ with $u \in \Oh_{C,P}^\times$.
	\end{enumerate}
	Moreover $\nu_P(f+g) \geq \min\{\nu_P(f),\nu_P(g)\}$.
\end{corollary}

\begin{proof}
	Write $\mm_P = (\pi_P)$, so $\mm_P^n = (\pi_P^n)$. Let $J = \bigcap_{n\geq 0}\mm_P^n$. This is a finitely generated ideal with $\mm_P J = J$ (since $\pi_P \mm_P^n = \mm_P^{n+1}$), so $J = 0$ by Nakayama. Thus no nonzero element vanishes to infinite order, and for $f \in \Oh_{C,P}\setminus\{0\}$ there is a unique $n \geq 0$ with $f \in \mm_P^n\setminus\mm_P^{n+1}$; set $\nu_P(f) = n$. Writing $f = \pi_P^n g$, we have $g \notin \mm_P$, so $g$ is a unit, giving (iii) for $f \in \Oh_{C,P}$.

	For $f \in \C(C)^\times\setminus\Oh_{C,P}$, write $f = g/h$ with $g,h \in \Oh_{C,P}$; then $g = \pi_P^ku$, $h = \pi_P^\ell u'$ with $u,u'$ units, so $f = \pi_P^{k-\ell}(u/u')$ with $k - \ell < 0$. Set $\nu_P(f) = k - \ell$. Multiplicativity, and hence the homomorphism property, is now clear from (iii), and surjectivity from $\nu_P(\pi_P^n) = n$. Properties (i) and (ii) restate the construction. The last inequality is immediate: if $\nu_P(f),\nu_P(g)\geq m$ then $f,g \in \mm_P^m$ and so $f+g \in \mm_P^m$.
\end{proof}

\begin{definition}
	If $\nu_P(f) = n > 0$ we say $f$ has a \vocab{zero of order $n$} at $P$; if $\nu_P(f) = -n < 0$, a \vocab{pole of order $n$}.
\end{definition}

\begin{example}
	Let $C = \A^1$ and $P = 0$, so $\Oh_{\A^1,0} = \{f/g : g(0)\neq 0\}$ and $t$ is a local parameter. Then $\nu_0(f/g)$ is the multiplicity of $t = 0$ as a root of $f$. Concretely,
	$$\nu_0\!\left(\frac{t^2+8t^3}{1+t^4}\right) = 2.$$
	Thinking of $f/g$ as a power series in $t$ (possible since $g(0)\neq 0$), $\nu_0$ is the index of the first nonzero coefficient. This is exactly the order of vanishing familiar from IB Complex Analysis, and the analogy is not a coincidence.
\end{example}

\begin{example}
	On $\PP^1$ with coordinate $t = X_1/X_0$, a local parameter at $a \in \C$ is $t - a$, and a local parameter at $\infty = (0:1)$ is $1/t$. So for $f = \prod_i (t-a_i)^{n_i}$ we have $\nu_{a_i}(f) = n_i$ and $\nu_\infty(f) = -\sum_i n_i$.
\end{example}

\begin{corollary}\label{cor:foneover}
	If $C$ is a curve and $P \in C$ is smooth, then for any $f \in \C(C)^\times$, at least one of $f$, $f^{-1}$ is regular at $P$.
\end{corollary}

\begin{proof}
	$\nu_P(f)$ and $\nu_P(f^{-1}) = -\nu_P(f)$ cannot both be negative.
\end{proof}

This fails badly in higher dimensions: on $\A^2$ neither $X/Y$ nor $Y/X$ is regular at the origin. It is a genuinely one-dimensional phenomenon, and it has the following striking consequence.

\begin{corollary}\label{cor:ratismorph}
	Let $C$ be a smooth curve. Then every rational map $\varphi \colon C \dto \PP^m$ is a morphism.
\end{corollary}

\begin{proof}
	Fix $P \in C$. Reordering coordinates, assume the image of $\varphi$ is not contained in $\{X_0 = 0\}$, so we may write $\varphi = (1 : g_1 : \cdots : g_m)$ with $g_j \in \C(C)$. If every $g_j \in \Oh_{C,P}$ then $\varphi$ is visibly regular at $P$. Otherwise let $t = \min_j \nu_P(g_j) < 0$ and choose a local parameter $\pi_P$. Then
	$$\varphi = \left(\pi_P^{-t} : \pi_P^{-t}g_1 : \cdots : \pi_P^{-t}g_m\right),$$
	an equivalent representative, and now every entry lies in $\Oh_{C,P}$ while at least one has $\nu_P = 0$, hence is nonzero at $P$. So $\varphi$ is regular at $P$.
\end{proof}

Since every projective variety sits inside some $\PP^m$, we conclude: \emph{a rational map from a smooth curve to a projective variety is automatically a morphism}. There are no indeterminacy loci in dimension one. This is the technical fact that makes the theory of curves so much cleaner than that of surfaces.

\subsection{Morphisms of Curves, Degree and Ramification}

\begin{proposition}\label{prop:degfinite}
	Let $\varphi \colon C \to D$ be a non-constant morphism of irreducible projective curves. Then every fibre $\varphi^{-1}(Q)$ is finite, $\varphi$ is dominant, and $\varphi^\star \colon \C(D)\hookrightarrow\C(C)$ makes $\C(C)$ a finite extension of $\C(D)$.
\end{proposition}

\begin{proof}
	$\varphi^{-1}(Q)$ is closed and, since $\varphi$ is non-constant, proper, hence finite by Proposition~\ref{prop:closedfinite}. As $C$ is infinite, $\varphi(C)$ is infinite, so its closure in $D$ is all of $D$ by Proposition~\ref{prop:closedfinite} again; so $\varphi$ is dominant and $\varphi^\star$ is defined and injective.

	Pick $t \in \C(D)\setminus\C$ and put $x = \varphi^\star t$. Both $\C(D)/\C(t)$ and $\C(C)/\C(x)$ are finite, since both fields have transcendence degree $1$ and $t$, $x$ are transcendental. Hence $\C(C)$ is finite over $\C(x) = \varphi^\star\C(t) \subseteq \varphi^\star\C(D)$, so a fortiori finite over $\varphi^\star \C(D)$.
\end{proof}

\begin{definition}[Degree]
	The \vocab{degree} of a non-constant morphism $\varphi \colon C \to D$ of curves is
	$$\deg\varphi = [\C(C) : \varphi^\star\C(D)].$$
\end{definition}

\begin{definition}[Ramification]
	Let $\varphi \colon C \to D$ be a non-constant morphism of curves, $P \in C$ smooth, $Q = \varphi(P)$ smooth. The \vocab{ramification index} of $\varphi$ at $P$ is
	$$e_P = e(\varphi,P) = \nu_P(\varphi^\star\pi_Q),$$
	where $\pi_Q$ is a local parameter at $Q$. We say $\varphi$ is \vocab{ramified} at $P$ if $e_P > 1$, and \vocab{unramified} otherwise.
\end{definition}

This is well defined: another local parameter differs from $\pi_Q$ by a unit of $\Oh_{D,Q}$, whose pullback is a unit of $\Oh_{C,P}$ and hence has $\nu_P = 0$.

\begin{example}\label{ex:zd}
	Consider $\varphi \colon \PP^1\to\PP^1$ given in the affine coordinate by $z \mapsto z^d$. On function fields, $\varphi^\star \colon \C(Y) \to \C(X)$ sends $Y \mapsto X^d$, so $\varphi^\star\C(Y) = \C(X^d)$ and $\deg\varphi = [\C(X):\C(X^d)] = d$.

	At $P = 0$ we have $Q = 0$ with local parameter $Y$, and $\varphi^\star Y = X^d$, so $e_0 = \nu_0(X^d) = d$: the map is totally ramified at $0$. The same holds at $\infty$. At $P = 1$ we have $Q = 1$ with local parameter $Y-1$, and $\varphi^\star(Y-1) = X^d - 1$, which has a simple zero at $X = 1$; so $e_1 = 1$. Indeed the fibre over $1$ consists of the $d$ distinct $d$th roots of unity, each with $e = 1$. Note $\sum_{P \in \varphi^{-1}(1)} e_P = d = \deg\varphi$, and also $\sum_{P\in\varphi^{-1}(0)}e_P = d$. This is a general phenomenon.
\end{example}

\begin{center}
\begin{tikzpicture}[x=1.05cm,y=0.85cm]
	\draw[thin, densely dotted, color=gray] (-1.6,3.35) -- (-1.6,-1.0);
	\draw[thin, densely dotted, color=gray] (-1.6,2.586) -- (-1.6,-1.0);
	\draw[thin, densely dotted, color=gray] (-1.6,1.414) -- (-1.6,-1.0);
	\draw[thin, densely dotted, color=gray] (1.7,3.30) -- (1.7,-1.0);
	\draw[thin, densely dotted, color=gray] (1.7,2.00) -- (1.7,-1.0);
	% sheet 1
	\draw[thick, color=blue, smooth] plot coordinates {(-3.0,3.20) (-1.6,3.35) (0,3.15) (1.7,3.30) (3.0,3.20)};
	% the folded pair of sheets
	\draw[thick, color=blue, smooth] plot coordinates {(-3.00,2.700) (-2.50,2.662) (-2.00,2.621) (-1.60,2.586) (-1.00,2.530) (-0.50,2.479) (0.00,2.421) (0.50,2.354) (0.90,2.289) (1.20,2.228) (1.45,2.161) (1.60,2.102) (1.68,2.046) (1.70,2.000) (1.68,1.954) (1.60,1.898) (1.45,1.839) (1.20,1.772) (0.90,1.711) (0.50,1.646) (0.00,1.579) (-0.50,1.521) (-1.00,1.470) (-1.60,1.414) (-2.00,1.379) (-2.50,1.338) (-3.00,1.300)};
	\node[color=blue] at (3.45,3.20) {$C$};
	% base curve D
	\draw[thick, color=blue] (-3.0,-1.0) -- (3.0,-1.0);
	\node[color=blue] at (3.45,-1.0) {$D$};
	\draw[->, thick] (-3.4,2.0) -- (-3.4,-0.6);
	\node[left] at (-3.45,0.7) {$\varphi$};
	\foreach \y in {3.35,2.586,1.414}{\filldraw[color=blue] (-1.6,\y) circle (1.5pt);}
	\filldraw[color=blue] (1.7,3.30) circle (1.5pt);
	\filldraw[color=blue] (1.7,2.00) circle (1.5pt);
	\filldraw (-1.6,-1.0) circle (1.8pt);
	\filldraw (1.7,-1.0) circle (1.8pt);
	\node[below] at (-1.6,-1.15) {$Q$};
	\node[below] at (1.7,-1.15) {$Q'$};
	\node[above] at (-1.6,3.42) {\small $e=1$};
	\node[above] at (-1.6,2.66) {\small $e=1$};
	\node[above] at (-1.6,1.49) {\small $e=1$};
	\node[above right] at (1.78,3.28) {\small $e=1$};
	\node[right] at (1.85,2.00) {\small $e=2$};
\end{tikzpicture}
\end{center}

The picture is the one to keep in mind: a degree $3$ map has three preimages over a general point, but the preimages can collide, and when two sheets come together the ramification index absorbs the deficit. The next theorem says the bookkeeping always balances.

\begin{theorem}\label{thm:degformula}
	Let $\varphi \colon C \to D$ be a non-constant morphism of irreducible projective curves.
	\begin{enumerate}[label=(\roman*)]
		\item $\varphi$ is surjective.
		\item If $C$ and $D$ are smooth then for every $Q \in D$,
		$$\deg\varphi = \sum_{P \in \varphi^{-1}(Q)} e_P.$$
		\item $e_P = 1$ for all but finitely many $P \in C$.
	\end{enumerate}
\end{theorem}

We prove (i) now, and (iii) will follow from the theory of differentials in Section~\ref{sec:differentials}. Part (ii) requires a piece of commutative algebra (that $\Oh_{D,Q}$ is a principal ideal domain and the semilocal ring $\bigcap_{P \in \varphi^{-1}(Q)}\Oh_{C,P}$ is free of rank $\deg\varphi$ over it) which we take on trust. The proof of (i) rests on the following, which is the projective analogue of compactness.

\begin{proposition}[Fundamental theorem of elimination theory]\label{prop:elimination}
	The projection $\PP^n\times\A^m \to \A^m$ is a closed map for the Zariski topology.
\end{proposition}

We do not prove this; it is a matter of showing that the condition on $\mathbf b \in \A^m$ that a system of homogeneous equations with coefficients depending on $\mathbf b$ has a nontrivial solution is itself given by polynomial equations in $\mathbf b$ (resultants). Note that it genuinely uses projectivity: $\A^1 \times \A^1 \to \A^1$ is not closed, since $\V(XY - 1)$ has image $\A^1\setminus\{0\}$.

\begin{proposition}\label{prop:projclosed}
	Let $\varphi \colon V \to W$ be a morphism of quasi-projective varieties with $V$ projective. Then $\varphi$ is a closed map.
\end{proposition}

Here a \vocab{quasi-projective variety} is a Zariski open subset of a projective variety; all affine and all projective varieties are quasi-projective, and all our definitions (rational functions, morphisms, local rings) extend verbatim.

\begin{proof}
	Factor $\varphi$ as $V \to \Gamma_\varphi \subseteq V\times W \to W$, where $\Gamma_\varphi = \{(x,\varphi(x))\}$ is the graph. The graph is closed, being the preimage of the diagonal $\Delta_W \subseteq W\times W$ (which is closed) under $\varphi\times\mathrm{id}$. So it suffices to show the projection $V \times W \to W$ is closed. As $V \subseteq \PP^n$ is closed, it suffices to show $\PP^n\times W \to W$ is closed; covering $W$ by affine opens, it suffices to show $\PP^n\times U \to U$ is closed for $U$ affine; and as $U$ is closed in some $\A^m$, it suffices to show $\PP^n \times \A^m \to \A^m$ is closed, which is Proposition~\ref{prop:elimination}.
\end{proof}

\begin{proof}[Proof of Theorem~\ref{thm:degformula}(i)]
	$\varphi(C)$ is closed in $D$ by Proposition~\ref{prop:projclosed}, so by Proposition~\ref{prop:closedfinite} it is either finite or all of $D$. As $\varphi$ is non-constant and $C$ is irreducible, $\varphi(C)$ is infinite, so $\varphi(C) = D$.
\end{proof}

Here is the payoff, and it is the algebraic geometer's version of Liouville's theorem.

\begin{corollary}\label{cor:liouville}
	Let $C$ be a smooth projective irreducible curve and $f \in \C(C)^\times$.
	\begin{enumerate}[label=(\roman*)]
		\item If $f$ is regular at every point of $C$, then $f$ is constant.
		\item $\{P \in C : \nu_P(f)\neq 0\}$ is finite, and $\displaystyle\sum_{P\in C}\nu_P(f) = 0$.
	\end{enumerate}
\end{corollary}

\begin{proof}
	Given $f$, consider the map $\varphi = (1 : f) \colon C \to \PP^1$, which is a morphism by Corollary~\ref{cor:ratismorph}. Write $0 = (1:0)$ and $\infty = (0:1)$ in $\PP^1$. Then $\varphi(P) = 0$ iff $f(P) = 0$, and $\varphi(P) = \infty$ iff $f$ has a pole at $P$.

	(i) If $f$ is everywhere regular then $\varphi$ misses $\infty$, so $\varphi$ is not surjective, so by Theorem~\ref{thm:degformula}(i) it is constant, i.e.\ $f$ is constant.

	(ii) Assume $f$ is non-constant, so $\varphi$ is a non-constant morphism, and let $n = \deg\varphi$. The function $t = X_1/X_0$ is a local parameter at $0 \in \PP^1$ and $\varphi^\star t = f$; so for $P$ with $f(P) = 0$ we get $e_P = \nu_P(f)$. Similarly $1/t$ is a local parameter at $\infty$ and $\varphi^\star(1/t) = 1/f$, so for $P$ a pole we get $e_P = -\nu_P(f)$. At all other $P$, $\nu_P(f) = 0$. The zeros and poles are the fibres over $0$ and $\infty$, hence finite, and Theorem~\ref{thm:degformula}(ii) applied to $Q = 0$ and $Q = \infty$ gives
	$$\sum_{f(P) = 0}\nu_P(f) = n = \sum_{f(P) = \infty}(-\nu_P(f)),$$
	so the total sum is zero.
\end{proof}

So on a projective curve there are no non-constant `holomorphic' functions, and every rational function has as many zeros as poles, counted with multiplicity. Both statements will be familiar to anyone who has met compact Riemann surfaces, and the analogy is exact.

\section{Divisors on Curves}\label{sec:divisors}

The field $\C(C)$ is enormous and hard to study directly. Divisor theory is a device for cutting it into finite-dimensional pieces: instead of studying all rational functions at once, we study those whose poles are bounded by a prescribed amount at prescribed points.

Throughout, $C$ is a smooth irreducible projective curve.

\subsection{Divisors, Degree and Linear Equivalence}

\begin{definition}[Divisor]
	A \vocab{divisor} on $C$ is an element of the free abelian group on the points of $C$,
	$$\Div(C) = \bigoplus_{P \in C}\Z[P].$$
	So a divisor is a finite formal sum $D = \sum_P n_P[P]$ with $n_P \in \Z$, almost all zero. Its \vocab{degree} is $\deg D = \sum_P n_P$. We write $\nu_P(D) = n_P$.
\end{definition}

The degree map $\deg \colon \Div(C)\to\Z$ is a surjective group homomorphism; its kernel is written $\Div^0(C)$.

\begin{definition}
	A divisor $D$ is \vocab{effective}, written $D \geq 0$, if $\nu_P(D)\geq 0$ for all $P$. We write $D \geq D'$ for $D - D' \geq 0$.
\end{definition}

\begin{definition}[Divisor of a Function]
	For $f \in \C(C)^\times$, the \vocab{divisor of $f$} is
	$$\Divv(f) = \sum_{P \in C}\nu_P(f)[P] \in \Div(C).$$
	Divisors of this form are called \vocab{principal}, and they form a subgroup $\Prin(C) \leq \Div(C)$.
\end{definition}

That the sum is finite, and that $\deg\Divv(f) = 0$, is Corollary~\ref{cor:liouville}. That $\Prin(C)$ is a subgroup follows from $\Divv(fg) = \Divv(f) + \Divv(g)$, which is the multiplicativity of $\nu_P$. So $\Prin(C) \leq \Div^0(C)$.

\begin{definition}[Picard Group]
	The \vocab{Picard group} or \vocab{divisor class group} of $C$ is
	$$\Pic(C) = \Cl(C) = \Div(C)/\Prin(C),$$
	and $\Pic^0(C) = \Cl^0(C) = \Div^0(C)/\Prin(C)$. Divisors $D, D'$ are \vocab{linearly equivalent}, written $D \sim D'$, if $D - D'$ is principal.
\end{definition}

The degree descends to $\Pic(C) \to \Z$, so linearly equivalent divisors have the same degree. The group $\Pic^0(C)$ measures the failure of the converse, and it turns out to be a subtle and extremely interesting invariant.

\begin{proposition}\label{prop:picP1}
	On $\PP^1$, every degree zero divisor is principal; that is, $\Pic^0(\PP^1) = 0$ and $\Pic(\PP^1)\cong\Z$ via the degree.
\end{proposition}

\begin{proof}
	Identify $\PP^1 = \C\cup\{\infty\}$ with coordinate $t$, and let $D = \sum_{a\in\C}n_a[a] + n_\infty[\infty]$ have degree $0$, so $n_\infty = -\sum_a n_a$. Put
	$$f = \prod_{a \in \C}(t-a)^{n_a} \in \C(t)^\times,$$
	a finite product. Then $\nu_a(f) = n_a$ for each $a \in \C$, since $t-a$ is a local parameter at $a$ and the other factors are units there. At $\infty$, using the local parameter $1/t$, each factor $(t-a)$ has $\nu_\infty = -1$, so $\nu_\infty(f) = -\sum_a n_a = n_\infty$. Hence $\Divv(f) = D$.
\end{proof}

That this fails for other curves is the source of the whole theory. We will see, for instance, that on a smooth cubic $\Pic^0(C) \cong C$ as a set, which is a very large group indeed.

\subsection{Hyperplane Sections and the Degree of a Curve}

There is a second, geometric, source of divisors.

\begin{definition}[Hyperplane Section]
	Let $C \subseteq \PP^n$ be a curve and $L$ a nonzero linear form with $C \not\subseteq \V(L)$. The \vocab{hyperplane section} of $C$ by $\V(L)$ is
	$$\Divv(L) = \sum_{P\in C} n_P[P], \qquad n_P = \nu_P\!\left(\frac{L}{X_i}\right) \text{ for any } i \text{ with } X_i(P)\neq 0.$$
\end{definition}

This is well defined: if $X_i(P), X_j(P) \neq 0$ then $X_i/X_j \in \Oh_{C,P}^\times$, so $\nu_P(X_i/X_j) = 0$ and $\nu_P(L/X_i) = \nu_P(L/X_j)$. Note $\Divv(L)$ is effective, since $L/X_i$ is regular at every such $P$, and it is supported on $C \cap \V(L)$, a finite set.

\begin{proposition}\label{prop:hyperplanediv}
	If $L,L'$ are linear forms with $C \not\subseteq \V(L)\cup\V(L')$, then
	$$\Divv(L) - \Divv(L') = \Divv\!\left(\frac{L}{L'}\right),$$
	which is principal. In particular $\deg\Divv(L) = \deg\Divv(L')$.
\end{proposition}

\begin{proof}
	At a point $P$ with $X_i(P)\neq0$ we have $\nu_P(L/L') = \nu_P(L/X_i) - \nu_P(L'/X_i)$, since $\nu_P$ is a homomorphism. The degrees agree because principal divisors have degree $0$.
\end{proof}

\begin{definition}[Degree of a Curve]
	The \vocab{degree} of a curve $C \subseteq \PP^n$ is $\deg C = \deg\Divv(L)$ for any linear form $L$ not vanishing on $C$.
\end{definition}

So the degree of a curve counts, with multiplicity, the points in which a general hyperplane meets it. A line in $\PP^2$ has degree $1$, a smooth conic degree $2$, and in general a plane curve $\V(F)$ with $\deg F = d$ has degree $d$ by Proposition~\ref{prop:bezoutline}.

The construction generalises in two useful directions. First, if $\varphi \colon C \to \PP^n$ is any non-constant morphism with $\varphi(C)\not\subseteq \V(L)$, we may set $\Divv(L) = \sum_P \nu_P(\varphi^\star L/X_i)[P]$; taking $\varphi$ to be an inclusion recovers the above. Secondly, if $G$ is homogeneous of degree $m \geq 1$ and does not vanish on $C$, we set
$$\Divv(G) = \sum_{P \in C}\nu_P\!\left(\frac{\varphi^\star G}{X_i^m}\right)[P],$$
again independent of the choice of $i$ with $X_i(P)\neq 0$. The same argument as above shows $\Divv(G) \sim m\,\Divv(L)$ for any linear form $L$, because $G/L^m$ is a rational function on $C$. This is exactly what we need to prove B\'ezout.

\begin{theorem}[Weak B\'ezout]\label{thm:weakbezout}
	Let $C, C' \subseteq \PP^2$ be distinct smooth irreducible projective plane curves of degrees $m$ and $n$. Then $|C\cap C'| \leq mn$.
\end{theorem}

\begin{proof}
	Write $C = \V(F)$, $C' = \V(G)$ with $\deg F = m$, $\deg G = n$. Restricting $G$ to $C$ gives a divisor $\Divv(G)$ on $C$, which is effective and supported exactly on $C \cap C'$ -- if $\nu_P(G/X_i^n) > 0$ then $G$ vanishes at $P$. Since $C \neq C'$ and both are irreducible, $G$ does not vanish identically on $C$, so $\Divv(G)$ is defined.

	By the remark above, $\Divv(G) \sim n\,\Divv(L)$ for a linear form $L$, so
	$$\deg\Divv(G) = n\deg\Divv(L) = n\deg C = nm.$$
	As $\Divv(G)$ is effective with support $C\cap C'$, each point contributes at least $1$ to the degree, so $|C\cap C'| \leq mn$.
\end{proof}

The full B\'ezout theorem is the assertion that $\nu_P(G/X_i^n)$ is exactly the intersection multiplicity $I_P(C,C')$, so that the count is an equality when weighted correctly. With the definition of $I_P$ we gave, this is essentially a tautology; the work is in showing that definition has the expected properties.

\subsection{Riemann--Roch Spaces}

Now we come to the central construction.

\begin{definition}[Riemann--Roch Space]
	Let $D \in \Div(C)$. The \vocab{Riemann--Roch space} of $D$ is
	$$L(D) = \{f \in \C(C)^\times : \Divv(f) + D \geq 0\}\cup\{0\},$$
	that is, the set of $f$ with $\nu_P(f) \geq -\nu_P(D)$ for all $P$. We write $\ell(D) = \dim_\C L(D)$.
\end{definition}

Unwinding: if $\nu_P(D) = n > 0$, then $f$ is allowed a pole of order at most $n$ at $P$; if $\nu_P(D) = -n < 0$, then $f$ is required to vanish to order at least $n$ at $P$; and if $\nu_P(D) = 0$, then $f$ must be regular at $P$. So $L(D)$ is the space of rational functions with poles bounded by $D$.

\begin{proposition}
	$L(D)$ is a $\C$-vector subspace of $\C(C)$.
\end{proposition}

\begin{proof}
	Closure under scalars is clear, and $\nu_P(f+g)\geq\min\{\nu_P(f),\nu_P(g)\}$ from Corollary~\ref{cor:valuation} gives closure under addition.
\end{proof}

\begin{example}\label{ex:LDP1}
	Let $C = \PP^1$, $\infty = (0:1)$, and $D = m[\infty]$ with $m \geq 0$. An element of $L(D)$ is a rational function with no poles in $\C$ -- hence a polynomial in $t$, by Proposition~\ref{prop:picP1} type reasoning -- with a pole of order at most $m$ at infinity, hence of degree at most $m$. So $L(D)$ is spanned by $1,t,t^2,\dots,t^m$ and
	$$\ell(m[\infty]) = m+1 = \deg D + 1.$$
\end{example}

\begin{proposition}\label{prop:ellbounds}
	Let $D$ be a divisor on $C$.
	\begin{enumerate}[label=(\roman*)]
		\item $L(0) = \C$, so $\ell(0) = 1$.
		\item If $\deg D < 0$ then $L(D) = 0$.
		\item For any $P \in C$, $\ell(D) \leq \ell(D - [P]) + 1$.
		\item If $\deg D \geq 0$ then $\ell(D)\leq \deg D + 1$. In particular $L(D)$ is always finite-dimensional.
	\end{enumerate}
\end{proposition}

\begin{proof}
	(i) $L(0)$ is the set of everywhere-regular functions, which is $\C$ by Corollary~\ref{cor:liouville}(i).

	(ii) If $f \in L(D)$ is nonzero then $\Divv(f) + D \geq 0$, so $0 \leq \deg(\Divv f + D) = \deg D$.

	(iii) Let $n = \nu_P(D)$ and $\pi_P$ a local parameter. Define
	$$\mathrm{ev}_P \colon L(D) \to \C, \qquad f \mapsto (\pi_P^{n}f)(P).$$
	This makes sense: $f \in L(D)$ has $\nu_P(f)\geq -n$, so $\nu_P(\pi_P^nf)\geq 0$ and $\pi_P^nf$ is regular at $P$. The map is $\C$-linear, and its kernel is $\{f : \nu_P(\pi_P^nf) > 0\} = \{f : \nu_P(f) \geq -n+1\} \cap L(D) = L(D - [P])$. So $\ell(D) \leq \ell(D-[P]) + 1$.

	(iv) Let $d = \deg D \geq 0$ and pick any $P$. Applying (iii) $d+1$ times,
	$$\ell(D) \leq \ell(D - (d+1)[P]) + d + 1 = d+1,$$
	since $\deg(D - (d+1)[P]) = -1 < 0$ and (ii) applies.
\end{proof}

Intuitively, (iii) says that relaxing the pole bound at one point by one can only gain one dimension: you gain the possibility of one extra term in the local expansion.

\begin{proposition}\label{prop:Lsim}
	If $D \sim E$ then $L(D)\cong L(E)$ as $\C$-vector spaces; in particular $\ell(D) = \ell(E)$.
\end{proposition}

\begin{proof}
	Write $D - E = \Divv(g)$. Multiplication by $g$ is a $\C$-linear map $L(D)\to L(E)$: if $\Divv(f)+D\geq0$ then $\Divv(fg) + E = \Divv(f) + \Divv(g) + E = \Divv(f)+D \geq 0$. Multiplication by $g^{-1}$ is its inverse.
\end{proof}

So $\ell$ is a function on $\Pic(C)$, and computing it is the central problem that Riemann--Roch solves.

\section{Differentials}\label{sec:differentials}

We have produced divisors in two ways so far: from rational functions, and from hyperplane sections. There is a third source, which will turn out to be the most important of all: differentials. Their divisor class is intrinsic to the curve -- it does not depend on any embedding -- and the dimension of its Riemann--Roch space is the genus.

\subsection{Differentials of a Field Extension}

The construction is formal, imitating the rules of calculus.

\begin{definition}[Space of Differentials]
	Let $K/\C$ be a field extension. The \vocab{space of differentials} $\Omega_{K/\C}$ is the quotient $M/N$, where $M$ is the $K$-vector space with basis the symbols $\delta x$ for $x \in K$, and $N \leq M$ is the subspace generated by
	$$\delta(x+y) - \delta x - \delta y, \qquad \delta(xy) - x\,\delta y - y\,\delta x, \qquad \delta a$$
	for $x,y \in K$ and $a \in \C$. We write $\dee x$ for the image of $\delta x$, and $\dee \colon K \to \Omega_{K/\C}$ for the resulting map, the \vocab{exterior derivative}. We abbreviate $\Omega_{K/\C}$ to $\Omega_K$.
\end{definition}

So $\Omega_K$ is the universal recipient of a `derivative'.

\begin{definition}[Derivation]
	Let $U$ be a $K$-vector space. A $\C$-linear map $D \colon K \to U$ is a \vocab{derivation} if $D(xy) = xD(y) + yD(x)$ for all $x,y$.
\end{definition}

\begin{lemma}[Universal property]\label{lem:univderiv}
	A $\C$-linear map $D \colon K \to U$ is a derivation if and only if there is a (necessarily unique) $K$-linear map $\lambda \colon \Omega_K \to U$ with $D = \lambda\circ\dee$.
\end{lemma}

\begin{proof}
	If such a $\lambda$ exists then $D$ inherits the Leibniz rule from $\dee$, which satisfies it by construction. Conversely, given a derivation $D$, define $\tilde\lambda \colon M \to U$ on the basis by $\delta x \mapsto D(x)$ and extend $K$-linearly. The relations generating $N$ are exactly the statements that $D$ is additive, Leibniz and kills $\C$ (note $D(1) = D(1\cdot1) = 2D(1)$, so $D(1) = 0$, and $\C$-linearity gives $D(a) = 0$). So $\tilde\lambda(N) = 0$ and $\tilde\lambda$ descends.
\end{proof}

\begin{remark}
	Every derivation obeys the quotient rule: if $y \neq 0$ then $D(x) = D(y\cdot\frac xy) = yD(\frac xy) + \frac xy D(y)$, so
	$$D\!\left(\frac xy\right) = \frac{yD(x) - xD(y)}{y^2}.$$
\end{remark}

\begin{lemma}\label{lem:chainrule}
	\begin{enumerate}[label=(\roman*)]
		\item Let $f = h/g \in \C(X_1,\dots,X_n)$ and let $x_1,\dots,x_n \in K$ be such that $y = f(x_1,\dots,x_n)$ is defined. Then
		$$\dee y = \sum_{i=1}^n \frac{\partial f}{\partial X_i}(x_1,\dots,x_n)\,\dee x_i.$$
		\item If $K = \C(x_1,\dots,x_n)$ then $\Omega_K$ is spanned over $K$ by $\dee x_1,\dots,\dee x_n$.
	\end{enumerate}
\end{lemma}

\begin{proof}
	(i) By $\C$-linearity it suffices to treat monomials, where it is the Leibniz rule and induction, and then quotients, where it is the quotient rule. (ii) is immediate from (i), since every element of $K$ is such an $f(x_1,\dots,x_n)$.
\end{proof}

\begin{theorem}\label{thm:omega1dim}
	Let $K/\C(t)$ be a finite extension with $t$ transcendental over $\C$. Then $\Omega_K$ is a one-dimensional $K$-vector space, spanned by $\dee t$.
\end{theorem}

\begin{proof}
	First suppose $K = \C(t)$. By Lemma~\ref{lem:chainrule}(ii), $\Omega_K$ is spanned by $\dee t$, so it has dimension $0$ or $1$; we must rule out $0$. By Lemma~\ref{lem:univderiv} it suffices to exhibit one nonzero derivation on $K$, and $\frac{\dee}{\dee t} \colon \C(t)\to\C(t)$ is one.

	Now suppose $K \supsetneq K_0 = \C(t)$. Since $\C$ has characteristic zero, $K/K_0$ is separable, so $K = K_0(\alpha)$ for some $\alpha$; let $h \in K_0[T]$ be its minimal polynomial. As $h$ is irreducible and $K_0$ has characteristic zero, $h$ has no repeated roots, so $h'(\alpha)\neq 0$.

	By Lemma~\ref{lem:chainrule}, $\Omega_K$ is spanned by $\dee t$ and $\dee\alpha$. Writing $D_th$ for the polynomial obtained by differentiating the coefficients of $h$ with respect to $t$, applying $\dee$ to the relation $h(\alpha) = 0$ gives
	$$0 = (D_th)(\alpha)\,\dee t + h'(\alpha)\,\dee\alpha,$$
	so $\dee\alpha = -\frac{(D_th)(\alpha)}{h'(\alpha)}\dee t$ and $\Omega_K$ is spanned by $\dee t$ alone.

	It remains to show $\Omega_K \neq 0$, for which we exhibit a nonzero derivation. Define $D \colon K_0[T]\to K$ by
	$$D\!\left(\sum_i c_iT^i\right) = \sum_i (D_tc_i)\alpha^i + \left(\sum_i ic_i\alpha^{i-1}\right)\cdot\left(\frac{-(D_th)(\alpha)}{h'(\alpha)}\right).$$
	One checks that $D$ is a derivation and that $D(h K_0[T]) = 0$, using the displayed formula for $D(h)$; hence $D$ descends to a nonzero derivation $K = K_0[T]/(h)\to K$.
\end{proof}

Since the function field of a curve is exactly a finite extension of some $\C(t)$ with $t$ transcendental (Proposition~\ref{prop:transbasis}), the theorem applies to all our curves.

\subsection{Rational Differentials on a Curve}

\begin{definition}
	Let $C$ be a curve. Write $\Omega_C = \Omega_{\C(C)/\C}$; its elements are \vocab{rational differentials} on $C$. A differential $\omega \in \Omega_C$ is \vocab{regular} at $P$ if it can be written $\omega = \sum_i f_i\,\dee g_i$ with all $f_i,g_i \in \Oh_{C,P}$. We write
	$$\Omega_{C,P} = \{\omega \in \Omega_C : \omega \text{ regular at } P\}.$$
\end{definition}

Note $\Omega_{C,P}$ is an $\Oh_{C,P}$-module but not a $\C(C)$-subspace: multiplying a regular differential by a function with a pole destroys regularity.

\begin{theorem}\label{thm:omegafree}
	Let $P \in C$ be a point with local parameter $\pi_P$. Then $\Omega_{C,P}$ is a free $\Oh_{C,P}$-module of rank one generated by $\dee\pi_P$; that is,
	$$\Omega_{C,P} = \{f\,\dee\pi_P : f \in \Oh_{C,P}\}.$$
\end{theorem}

\begin{proof}
	First, $\Omega_{C,P}$ is finitely generated over $\Oh_{C,P}$: choose an affine patch $C_0 \ni P$ with $\C[C_0] = \C[x_1,\dots,x_n]$. Given $f \in \Oh_{C,P}$, write $f = g/h$ with $g,h$ polynomials in the $x_i$ and $h(P)\neq0$; by the quotient rule and Lemma~\ref{lem:chainrule},
	$$\dee f = \sum_{i=1}^n \frac{h\frac{\partial g}{\partial X_i} - g\frac{\partial h}{\partial X_i}}{h^2}(x_1,\dots,x_n)\,\dee x_i,$$
	and every coefficient is regular at $P$ since $h(P)\neq0$. So $\Omega_{C,P}$ is generated by $\dee x_1,\dots,\dee x_n$.

	Now clearly $\Oh_{C,P}\dee\pi_P \subseteq \Omega_{C,P}$. Conversely, for $f \in \Oh_{C,P}$ we have $f - f(P) \in \mm_P = (\pi_P)$, so $f = f(P) + \pi_P g$ with $g \in \Oh_{C,P}$; the Leibniz rule gives
	$$\dee f = g\,\dee\pi_P + \pi_P\,\dee g \in \Oh_{C,P}\dee\pi_P + \mm_P\,\Omega_{C,P}.$$
	Hence $\Omega_{C,P} = \Oh_{C,P}\dee\pi_P + \mm_P\Omega_{C,P}$, and Nakayama's lemma (Lemma~\ref{lem:nakayama}(ii)) gives $\Omega_{C,P} = \Oh_{C,P}\dee\pi_P$. Freeness holds because $\dee\pi_P \neq 0$ in the one-dimensional space $\Omega_C$ and $\Oh_{C,P}$ is a domain.
\end{proof}

Since $\Omega_C$ is one-dimensional over $\C(C)$, every $\omega \in \Omega_C$ can be written $\omega = f\,\dee\pi_P$ for a unique $f \in \C(C)$. This lets us define orders of vanishing for differentials.

\begin{definition}
	For $\omega \in \Omega_C\setminus\{0\}$ and $P \in C$, write $\omega = f\,\dee\pi_P$ and set $\nu_P(\omega) = \nu_P(f)$.
\end{definition}

This is well defined: if $\pi'$ is another local parameter then $\pi' = u\pi_P$ with $u$ a unit, so $\dee\pi' = u\,\dee\pi_P + \pi_P\,\dee u$, and $\dee u = h\,\dee\pi_P$ with $h$ regular, so $\dee\pi' = (u + \pi_Ph)\dee\pi_P$ with $u + \pi_Ph$ a unit. Hence $\nu_P$ is unchanged.

\begin{lemma}\label{lem:omegafinite}
	Let $\omega \in \Omega_C$ be nonzero. Then $\nu_P(\omega)\neq 0$ for only finitely many $P$.
\end{lemma}

\begin{proof}
	Since $\nu_P(f\dee g) = \nu_P(f) + \nu_P(\dee g)$ and $\nu_P(f) = 0$ for all but finitely many $P$, it suffices to treat $\omega = \dee g$. Here $g$ is non-constant (else $\dee g = 0$), so $g$ is transcendental over $\C$ and $\C(C)/\C(g)$ is finite. Consider the morphism $\varphi = (1:g)\colon C \to \PP^1$. By Corollary~\ref{cor:liouville} there are finitely many $P$ with $g(P) = \infty$, and by Theorem~\ref{thm:degformula}(iii) finitely many with $e_P > 1$.

	At any other $P$, put $Q = \varphi(P) = g(P) \in \C$. Then $t - g(P)$ is a local parameter at $Q$, and since $e_P = 1$, $\varphi^\star(t - g(P)) = g - g(P)$ is a local parameter at $P$. Hence $\dee g = \dee(g - g(P))$ is $\dee$ of a local parameter, so $\nu_P(\dee g) = 0$.
\end{proof}

(Strictly, this argument uses Theorem~\ref{thm:degformula}(iii), which we have not yet proved. The logic is not circular: (iii) follows from the case $\omega = \dee t$ on $\PP^1$, computed explicitly in Example~\ref{ex:divdt} below, together with the ramification formula of Lemma~\ref{lem:pullbackdiff}. We record the dependency and move on.)

\begin{definition}[Canonical Divisor]
	For $\omega \in \Omega_C\setminus\{0\}$, set $\Divv(\omega) = \sum_{P\in C}\nu_P(\omega)[P] \in \Div(C)$.
\end{definition}

\begin{proposition}
	If $\omega,\omega'$ are nonzero rational differentials on $C$, then $\Divv(\omega) - \Divv(\omega')$ is principal.
\end{proposition}

\begin{proof}
	$\Omega_C$ is one-dimensional over $\C(C)$, so $\omega = f\omega'$ for some $f \in \C(C)^\times$, and then $\nu_P(\omega) = \nu_P(f) + \nu_P(\omega')$ at every $P$.
\end{proof}

\begin{definition}[Canonical Class]
	The class $K_C = [\Divv(\omega)] \in \Pic(C)$, for any nonzero $\omega \in \Omega_C$, is the \vocab{canonical class} of $C$. Any divisor in this class is a \vocab{canonical divisor}.
\end{definition}

\begin{example}\label{ex:divdt}
	Let $C = \PP^1$ with affine coordinate $t$. For $a \in \C$, $t - a$ is a local parameter at $a$ and $\dee t = \dee(t-a)$, so $\nu_a(\dee t) = 0$. At $\infty$, the local parameter is $s = 1/t$, and $t = 1/s$ gives $\dee t = -s^{-2}\dee s$, so
	$$\nu_\infty(\dee t) = \nu_\infty(-s^{-2}) = -2.$$
	Hence
	$$\Divv(\dee t) = -2[\infty], \qquad \deg K_{\PP^1} = -2.$$
	In particular $K_{\PP^1}$ has nonzero degree, so it is not principal: there is no rational differential on $\PP^1$ with neither zeros nor poles.
\end{example}

\subsection{The Genus}

\begin{definition}[Genus]
	The \vocab{genus} of a curve $C$ is $g(C) = \ell(K_C)$.
\end{definition}

By Proposition~\ref{prop:Lsim} this is well defined, even though $L(K_C)$ itself depends on the choice of representative. Concretely, $L(K_C)$ is (isomorphic to) the space of \vocab{regular differentials} on $C$: for $\omega_0$ a fixed nonzero differential, $f \mapsto f\omega_0$ identifies $L(\Divv\omega_0)$ with $\{\omega \in \Omega_C : \Divv\omega \geq 0\}$, the differentials with no poles. So the genus counts the everywhere-regular differentials.

\begin{example}
	$g(\PP^1) = 0$. Indeed $K_{\PP^1} = -2[\infty]$ has degree $-2 < 0$, so $\ell(K_{\PP^1}) = 0$ by Proposition~\ref{prop:ellbounds}(ii).
\end{example}

The genus is the invariant that will organise everything that follows. Topologically, a smooth projective curve over $\C$ is a compact orientable real surface, and $g$ is the number of handles: $\PP^1$ is a sphere ($g = 0$), a smooth cubic is a torus ($g=1$), and in general the genus $g$ surface looks like this.

\begin{center}
\begin{tikzpicture}[x=1cm,y=1cm]
	% outline of a genus 3 surface
	\draw[thick, color=blue] (-4.1,0) .. controls (-4.1,1.35) and (-2.75,1.35) .. (-2.4,1.05)
		.. controls (-2.05,0.78) and (-1.75,0.78) .. (-1.4,1.05)
		.. controls (-1.05,1.35) and (-0.35,1.35) .. (0,1.05)
		.. controls (0.35,0.78) and (0.65,0.78) .. (1.0,1.05)
		.. controls (1.35,1.35) and (2.05,1.35) .. (2.4,1.05)
		.. controls (2.75,0.78) and (3.4,0.9) .. (3.7,0.55)
		.. controls (4.1,0.1) and (4.1,-0.1) .. (3.7,-0.55)
		.. controls (3.4,-0.9) and (2.75,-0.78) .. (2.4,-1.05)
		.. controls (2.05,-1.35) and (1.35,-1.35) .. (1.0,-1.05)
		.. controls (0.65,-0.78) and (0.35,-0.78) .. (0,-1.05)
		.. controls (-0.35,-1.35) and (-1.05,-1.35) .. (-1.4,-1.05)
		.. controls (-1.75,-0.78) and (-2.05,-0.78) .. (-2.4,-1.05)
		.. controls (-2.75,-1.35) and (-4.1,-1.35) .. (-4.1,0);
	% handles
	\foreach \c in {-2.75,-0.35,2.05}{
		\draw[thick, color=blue] (\c-0.42,0.12) .. controls (\c-0.22,-0.24) and (\c+0.22,-0.24) .. (\c+0.42,0.12);
		\draw[thick, color=blue] (\c-0.31,0.00) .. controls (\c-0.16,0.16) and (\c+0.16,0.16) .. (\c+0.31,0.00);
	}
	\node at (0,-1.75) {a curve of genus $3$};
\end{tikzpicture}
\end{center}

This is not a definition we can use -- we have defined $g$ algebraically -- but it is the right picture, and the numerology will confirm it: we will show $\deg K_C = 2g-2$, and $2 - 2g$ is precisely the Euler characteristic of a genus $g$ surface.

\subsection{The Canonical Class of a Plane Curve}

We can compute the canonical class of a smooth plane curve completely explicitly. The following calculation is the model for everything: choose a convenient differential, and compute its order of vanishing at every point using Corollary~\ref{cor:localparamplane}.

\begin{example}[Smooth plane cubics]\label{ex:cubicgenus}
	Let $C = \V(F) \subseteq \PP^2$ with $F = X_0X_2^2 - \prod_{i=1}^3(X_1 - \lambda_iX_0)$ and $\lambda_1,\lambda_2,\lambda_3$ distinct; this is the Weierstrass form of Theorem~\ref{thm:weierstrass}, and it is smooth. In the affine chart $X_0 = 1$ with coordinates $x = X_1/X_0$, $y = X_2/X_0$, the equation is
	$$f(x,y) = y^2 - p(x) = 0, \qquad p(x) = \prod_{i=1}^3(x - \lambda_i).$$
	Applying $\dee$ gives the relation $2y\,\dee y = p'(x)\,\dee x$ in $\Omega_C$. Take
	$$\omega = \frac{\dee x}{y} = \frac{2\,\dee y}{p'(x)}.$$
	We compute $\Divv(\omega)$.

	If $P = (a,b)$ with $b \neq 0$, then $\partial f/\partial y = 2b \neq 0$, so $x - a$ is a local parameter at $P$ by Corollary~\ref{cor:localparamplane}, and $\nu_P(\dee x) = 0$; also $\nu_P(1/y) = 0$. So $\nu_P(\omega) = 0$.

	If $P = (\lambda_i,0)$, then $\partial f/\partial x = -p'(\lambda_i)\neq 0$ (the $\lambda_i$ are distinct), so $y$ is a local parameter at $P$ and $\nu_P(\dee y) = 0$; also $\nu_P(p'(x)) = 0$ since $p'(\lambda_i)\neq0$. Using the second expression, $\nu_P(\omega) = 0$.

	Finally, at the point at infinity $P_0 = (0:0:1)$: in the chart $X_2 = 1$, with $z = X_0/X_2 = 1/y$ and $w = X_1/X_2 = x/y$, the equation becomes $z = \prod_i(w-\lambda_iz)$, and one computes that $w$ is a local parameter at $P_0$ with $\nu_{P_0}(z) = 3$, $\nu_{P_0}(x) = \nu_{P_0}(w/z) = -2$ and $\nu_{P_0}(y) = \nu_{P_0}(1/z) = -3$. Then $\nu_{P_0}(\dee x) = \nu_{P_0}(x) - 1 = -3$, so
	$$\nu_{P_0}(\omega) = \nu_{P_0}(\dee x) - \nu_{P_0}(y) = -3 - (-3) = 0.$$
	Hence $\Divv(\omega) = 0$ and $K_C = 0$. Therefore
	$$g(C) = \ell(K_C) = \ell(0) = 1.$$
\end{example}

\begin{theorem}\label{thm:cubicgenus1}
	Every smooth plane cubic has genus $1$. In particular no smooth plane cubic is isomorphic to $\PP^1$.
\end{theorem}

\begin{proof}
	By Theorem~\ref{thm:weierstrass} we may change coordinates into the shape of Example~\ref{ex:cubicgenus}, and the genus is an isomorphism invariant. Since $g(\PP^1) = 0 \neq 1$, the last statement follows.
\end{proof}

The general computation runs along exactly the same lines.

\begin{theorem}\label{thm:canonicalplane}
	Let $C = \V(F)\subseteq\PP^2$ be a smooth plane curve of degree $d$. Then
	$$K_C = (d-3)H,$$
	where $H$ is the class of a hyperplane section of $C$. In particular $\deg K_C = d(d-3)$.
\end{theorem}

\begin{proof}
	Change coordinates so that $(0:1:0)\notin C$. Let $x = X_1/X_0$ and $y = X_2/X_0$ in $\C(C)$, and set $f(X,Y) = F(1,X,Y)$, so $f(x,y) = 0$. Applying $\dee$,
	$$\frac{\partial f}{\partial X}(x,y)\,\dee x + \frac{\partial f}{\partial Y}(x,y)\,\dee y = 0 \quad\text{in }\Omega_C.$$
	Choose
	$$\omega = \frac{\dee x}{\frac{\partial f}{\partial Y}(x,y)} = \frac{-\dee y}{\frac{\partial f}{\partial X}(x,y)}.$$
	If $P \in C\cap U_0$ has $\frac{\partial f}{\partial Y}(P)\neq 0$, then $x - x(P)$ is a local parameter and $\nu_P(\omega) = -\nu_P\big(\frac{\partial f}{\partial Y}\big) = 0$. Otherwise smoothness gives $\frac{\partial f}{\partial X}(P)\neq0$, and the second expression gives $\nu_P(\omega) = 0$ likewise. So $\omega$ has neither zeros nor poles on the affine patch $U_0$.

	Since $(0:1:0)\notin C$, the points of $C$ at infinity (i.e.\ with $X_0 = 0$) all have $X_2 \neq 0$, so lie in $U_2$. Work there with $z = X_0/X_2$ and $w = X_1/X_2$, and $g(Z,W) = F(Z,W,1)$, and set
	$$\eta = \frac{\dee z}{\frac{\partial g}{\partial W}(z,w)} = \frac{-\dee w}{\frac{\partial g}{\partial Z}(z,w)}.$$
	The same argument gives $\nu_P(\eta) = 0$ for all $P \in C\cap U_2$.

	It remains to compare $\omega$ and $\eta$. We have $z = 1/y$, $w = x/y$, and $F$ homogeneous of degree $d$ gives $f(X,Y) = Y^dg(1/Y, X/Y)$. Differentiating this identity and substituting, a short computation yields
	$$\omega = z^{d-3}\,\eta.$$
	Hence for $P$ with $X_2(P)\neq0$,
	$$\nu_P(\omega) = (d-3)\nu_P(z) + \nu_P(\eta) = (d-3)\nu_P(z),$$
	and since $z = X_0/X_2$, the divisor $\sum_P \nu_P(z)[P]$ restricted to $U_2$ is exactly $\Divv(X_0)$, the hyperplane section by $\V(X_0)$ (which is supported at infinity, where $\omega$ was not already computed to be zero). Combining the two patches,
	$$\Divv(\omega) = (d-3)\Divv(X_0),$$
	which is $(d-3)H$.
\end{proof}

\begin{corollary}\label{cor:notiso}
	Smooth plane curves of distinct degrees $d,d'\geq 2$ are never isomorphic.
\end{corollary}

\begin{proof}
	$\deg K_C = d(d-3)$ depends only on the isomorphism class of $C$, and $d \mapsto d(d-3)$ is injective on $\{d \geq 2\}$ (it is strictly increasing for $d \geq 2$ apart from $d(d-3)$ taking the values $-2,0,4,10,\dots$ at $d = 2,3,4,5$, which are distinct).
\end{proof}

In particular there are infinitely many pairwise non-isomorphic smooth projective curves, which is already far from obvious.

\begin{proposition}\label{prop:canonicalbasis}
	Let $C$ be a smooth plane curve of degree $d \geq 3$ with affine equation $f(x,y) = 0$ as above. Then
	$$\left\{\frac{x^ry^s\,\dee x}{\frac{\partial f}{\partial Y}(x,y)} : r,s \geq 0,\ r+s \leq d-3\right\}$$
	is a basis for $L(K_C)$, taking the representative $(d-3)H$ with $H$ the hyperplane at infinity. Consequently
	$$g(C) = \#\{(r,s) : r,s\geq0,\ r+s\leq d-3\} = \binom{d-1}{2} = \frac{(d-1)(d-2)}{2}.$$
\end{proposition}

\begin{proof}[Proof sketch]
	The proof runs exactly as in Theorem~\ref{thm:canonicalplane}: multiplying $\omega$ by the monomial $x^ry^s$ shifts the order of vanishing at infinity by a controlled amount, and the condition $r+s\leq d-3$ is exactly what keeps $\Divv(x^ry^s\omega) + (d-3)H$ effective. Linear independence holds because the monomials $x^ry^s$ with $r+s\leq d-3$ are linearly independent in $\C(C)$, as $\deg f = d$. We omit the details.
\end{proof}

The formula $g = \frac{(d-1)(d-2)}{2}$ is the \vocab{genus--degree formula}. For $d = 1,2$ it gives $g = 0$, consistent with lines and conics being isomorphic to $\PP^1$; for $d = 3$ it gives $g = 1$, matching Theorem~\ref{thm:cubicgenus1}; for $d = 4$ it gives $g = 3$. We will re-derive it from Riemann--Roch in a moment, which is the more illuminating route.

\section{The Riemann--Roch Theorem}

We have two invariants of a divisor $D$ on a curve $C$: the degree $\deg D$, which is easy to compute, and the dimension $\ell(D)$, which is what we actually want. Proposition~\ref{prop:ellbounds} gives the crude bound $\ell(D)\leq\deg D + 1$. Riemann--Roch says that this is an equality up to a correction term measured by the genus and a second Riemann--Roch space.

\begin{theorem}[Riemann--Roch]\label{thm:RR}
	Let $C$ be a smooth irreducible projective curve of genus $g$, and let $D \in \Div(C)$. Let $K = K_C$ be a canonical divisor. Then
	$$\ell(D) - \ell(K - D) = \deg D - g + 1.$$
\end{theorem}

The proof is beyond the scope of this course. It is worth knowing that it is a genuinely deep statement, and that it has several very different proofs: a sheaf-cohomological one, in which $\ell(D) = \dim H^0(C,\Oh(D))$ and $\ell(K-D) = \dim H^1(C,\Oh(D))$ and the theorem is a computation of Euler characteristics; an analytic one via Hodge theory; and Weil's original adelic proof. It is closely related to Stokes' theorem and the Gauss--Bonnet theorem, both of which relate a local differential-geometric quantity to a global topological one.

Let us extract its consequences, which is where all the fun is.

\subsection{First Consequences}

\begin{corollary}\label{cor:degK}
	$\deg K_C = 2g - 2$.
\end{corollary}

\begin{proof}
	Apply Riemann--Roch with $D = K$:
	$$\ell(K) - \ell(0) = \deg K - g + 1.$$
	Now $\ell(K) = g$ by definition of the genus and $\ell(0) = 1$ by Proposition~\ref{prop:ellbounds}(i), so $g - 1 = \deg K - g + 1$, giving $\deg K = 2g-2$.
\end{proof}

Note that $2g-2 = -\chi(C)$ where $\chi$ is the Euler characteristic of the associated real surface. This is the promised confirmation that our algebraic genus is the topological one.

\begin{corollary}\label{cor:RRlarge}
	If $\deg D > 2g - 2$ then $\ell(D) = \deg D - g + 1$.
\end{corollary}

\begin{proof}
	$\deg(K - D) = 2g-2-\deg D < 0$, so $\ell(K-D) = 0$ by Proposition~\ref{prop:ellbounds}(ii).
\end{proof}

So for divisors of large degree, $\ell$ is completely determined by the degree and the genus. The `interesting' range $0 \leq \deg D \leq 2g-2$ is exactly where the geometry of the particular curve enters.

\begin{corollary}[Genus--degree formula]\label{cor:genusdegree}
	Let $C \subseteq \PP^2$ be a smooth plane curve of degree $d$. Then
	$$g(C) = \frac{(d-1)(d-2)}{2}.$$
\end{corollary}

\begin{proof}
	For $d = 1,2$ we know $C \cong \PP^1$ (Proposition~\ref{prop:conicisP1}), so $g = 0$, matching the formula. For $d \geq 3$, Theorem~\ref{thm:canonicalplane} gives $K_C = (d-3)H$ with $\deg H = \deg C = d$, so $\deg K_C = d(d-3)$. By Corollary~\ref{cor:degK},
	$$2g - 2 = d(d-3) \implies g = \frac{d^2-3d+2}{2} = \frac{(d-1)(d-2)}{2}. \qedhere$$
\end{proof}

\begin{example}
	A smooth plane quartic has genus $3$, a quintic genus $6$, a sextic genus $10$. Note that not every integer occurs: there is no smooth plane curve of genus $2$. (There are plenty of genus $2$ curves; they just do not embed smoothly in $\PP^2$.)
\end{example}

\begin{corollary}\label{cor:g0}
	If $g(C) = 0$ then $C \cong \PP^1$.
\end{corollary}

\begin{proof}
	Pick any point $P \in C$ and let $D = [P]$, of degree $1 > -2 = 2g-2$. By Corollary~\ref{cor:RRlarge}, $\ell(D) = 1 - 0 + 1 = 2$. So $L(D)$ contains a non-constant function $f$ (as the constants give only a one-dimensional subspace), whose only possible pole is a simple pole at $P$. Consider $\varphi = (1:f)\colon C \to \PP^1$, a morphism by Corollary~\ref{cor:ratismorph}. Its degree equals the number of poles of $f$ counted with multiplicity, which by the proof of Corollary~\ref{cor:liouville} is $-\nu_P(f) = 1$. A degree $1$ morphism induces an isomorphism of function fields, so $C \cong \PP^1$ by Theorem~\ref{thm:biratfield}; and a birational morphism of smooth projective curves is an isomorphism (its inverse is a rational map, hence a morphism by Corollary~\ref{cor:ratismorph}).
\end{proof}

So genus $0$ curves are exactly the rational ones, and there is only one of them.

\subsection{Riemann--Hurwitz}

We now compare the genera of two curves related by a morphism. The key is that differentials pull back.

Let $\varphi \colon C \to D$ be a non-constant morphism of curves. If $\omega = f\,\dee t \in \Omega_D$, define
$$\varphi^\star\omega = (\varphi^\star f)\,\dee(\varphi^\star t) \in \Omega_C.$$
This is well defined by Theorem~\ref{thm:omega1dim} applied on both sides.

\begin{lemma}\label{lem:pullbackdiff}
	Let $P \in C$, $Q = \varphi(P)$, with local parameters $\pi_P,\pi_Q$ and ramification index $e_P$. Then for any nonzero $\omega \in \Omega_D$,
	$$\nu_P(\varphi^\star\omega) = e_P\,\nu_Q(\omega) + e_P - 1.$$
\end{lemma}

\begin{proof}
	By definition of $e_P$ we may write $\varphi^\star\pi_Q = u\pi_P^{e_P}$ with $u \in \Oh_{C,P}^\times$. Differentiating,
	$$\varphi^\star(\dee\pi_Q) = \dee(u\pi_P^{e_P}) = \left(e_Pu\pi_P^{e_P-1} + \pi_P^{e_P}\frac{\dee u}{\dee \pi_P}\right)\dee\pi_P.$$
	The bracket has $\nu_P = e_P - 1$: the first term has order exactly $e_P-1$ (here we use characteristic zero, so $e_P \neq 0$ in $\C$), and the second has order at least $e_P$. So $\nu_P(\varphi^\star\dee\pi_Q) = e_P-1$, which is the case $\nu_Q(\omega) = 0$.

	In general write $\omega = v\pi_Q^m\dee\pi_Q$ with $v \in \Oh_{D,Q}^\times$ and $m = \nu_Q(\omega)$, using Theorem~\ref{thm:omegafree}. Then
	$$\varphi^\star\omega = (\varphi^\star v)(\varphi^\star\pi_Q)^m\,\varphi^\star(\dee\pi_Q),$$
	and taking $\nu_P$ gives $0 + me_P + (e_P-1)$ as claimed.
\end{proof}

The formula generalises the elementary fact that $\frac{\dee}{\dee x}x^n = nx^{n-1}$ drops the order by one; ramification of order $e$ costs $e-1$.

\begin{theorem}[Riemann--Hurwitz]\label{thm:RH}
	Let $\varphi \colon C \to D$ be a non-constant morphism of smooth projective irreducible curves of degree $n$. Then
	$$2g(C) - 2 = n\big(2g(D) - 2\big) + \sum_{P \in C}(e_P - 1).$$
\end{theorem}

\begin{proof}
	Choose a nonzero $\omega \in \Omega_D$. Then $\varphi^\star\omega$ is a nonzero element of $\Omega_C$, so $\Divv(\varphi^\star\omega)$ is a canonical divisor on $C$ and by Corollary~\ref{cor:degK} its degree is $2g(C)-2$. Now compute, grouping the points of $C$ into fibres and using Lemma~\ref{lem:pullbackdiff}:
	\begin{align*}
		2g(C)-2 &= \sum_{P\in C}\nu_P(\varphi^\star\omega)\\
		&= \sum_{Q\in D}\ \sum_{P \in \varphi^{-1}(Q)}\big(e_P\nu_Q(\omega) + e_P-1\big)\\
		&= \sum_{Q\in D}\left(\nu_Q(\omega)\sum_{P\in\varphi^{-1}(Q)}e_P + \sum_{P\in\varphi^{-1}(Q)}(e_P-1)\right)\\
		&= \sum_{Q\in D}\Big(n\,\nu_Q(\omega)\Big) + \sum_{P \in C}(e_P-1)\\
		&= n\deg\Divv(\omega) + \sum_{P\in C}(e_P-1)\\
		&= n\big(2g(D)-2\big) + \sum_{P\in C}(e_P-1),
	\end{align*}
	where the fourth line uses Theorem~\ref{thm:degformula}(ii).
\end{proof}

Note that the ramification term is a sum of non-negative integers, and (by Theorem~\ref{thm:degformula}(iii)) all but finitely many terms vanish, so the sum is finite. Conversely, running the argument in the case $D = \PP^1$ shows that the sum \emph{must} be finite, which is how one proves Theorem~\ref{thm:degformula}(iii).

\begin{corollary}\label{cor:genusincrease}
	With notation as above, $g(C) \geq g(D)$. If $g(C) < g(D)$ then every rational map $C \dto D$ is constant.
\end{corollary}

\begin{proof}
	If $g(D) \geq 1$ then $2g(D)-2\geq 0$ and Riemann--Hurwitz gives $2g(C)-2 \geq n(2g(D)-2)\geq 2g(D)-2$ since $n\geq1$, so $g(C)\geq g(D)$. If $g(D) = 0$ there is nothing to prove. For the second statement, recall that a rational map from a smooth curve is a morphism (Corollary~\ref{cor:ratismorph}), so a non-constant one would contradict the first part.
\end{proof}

\begin{example}
	There is no non-constant morphism $\PP^1 \to C$ for any curve $C$ of genus $\geq 1$. In particular a smooth plane cubic cannot be parametrised by rational functions -- which is a purely algebraic statement with a genuinely geometric proof.
\end{example}

\begin{example}[Hyperelliptic curves]\label{ex:hyperelliptic}
	Consider the smooth projective model $C$ of the affine curve $y^2 = p(x)$ where $p$ has degree $2m$ with distinct roots, and let $\varphi \colon C \to \PP^1$ be $(x,y)\mapsto x$. This has degree $2$, since $\C(C) = \C(x)[y]/(y^2-p(x))$ is a degree $2$ extension of $\C(x)$.

	Ramification occurs exactly where the two preimages coincide, i.e.\ at the $2m$ roots of $p$, and there $e_P = 2$; one checks the two points at infinity are unramified when $\deg p$ is even. So Riemann--Hurwitz gives
	$$2g(C)-2 = 2(-2) + 2m\cdot 1 = 2m-4,$$
	so $g(C) = m-1$. Taking $m = 2$ (a quartic $p$) gives $g = 1$; $m = 3$ gives $g = 2$; and in general every genus is achieved. So genus $2$ curves exist even though no plane curve has genus $2$.
\end{example}

\begin{example}
	The map $\varphi \colon \PP^1\to\PP^1$, $z \mapsto z^d$, has degree $d$ and is ramified exactly at $0$ and $\infty$, each with $e = d$. Riemann--Hurwitz reads
	$$-2 = d(-2) + 2(d-1),$$
	which is indeed an identity. This is a useful sanity check on the signs.
\end{example}

\begin{definition}
	A curve $C$ with $g(C)\geq 2$ is \vocab{hyperelliptic} if it admits a morphism $C \to \PP^1$ of degree $2$.
\end{definition}

The curves of Example~\ref{ex:hyperelliptic} with $m \geq 3$ are hyperelliptic. One shows (it is on the last example sheet) that $C$ is hyperelliptic if and only if there is a divisor $D$ with $\deg D = 2$ and $\ell(D) = 2$; the degree $2$ map is then $\varphi_D$, constructed below.

\section{Curves of Low Genus}\label{sec:lowgenus}

We finish by using Riemann--Roch to answer concrete questions: which curves are there, and how do we embed them in projective space?

\subsection{Maps to Projective Space from Divisors}

Riemann--Roch produces functions; functions produce maps to projective space. Here is the construction.

\begin{definition}
	Let $D$ be a divisor on $C$ with $\ell(D) = n+1 \geq 2$, and let $f_0,\dots,f_n$ be a basis of $L(D)$. The \vocab{morphism associated to $D$} is
	$$\varphi_D \colon C \to \PP^n, \qquad P \mapsto (f_0(P) : \cdots : f_n(P)).$$
\end{definition}

This is a rational map, hence (Corollary~\ref{cor:ratismorph}) a morphism, and it is independent of the choice of basis up to an element of $\PGL_{n+1}(\C)$. We say $\varphi_D$ is an \vocab{embedding} if it is an isomorphism onto its image.

\begin{definition}
	A divisor $D$ satisfies condition $(\star)$ if for all $P,Q \in C$ (not necessarily distinct),
	$$\ell(D - [P] - [Q]) = \ell(D) - 2.$$
\end{definition}

By Proposition~\ref{prop:ellbounds}(iii), $\ell(D)-2 \leq \ell(D-[P]-[Q]) \leq \ell(D)$ always, so $(\star)$ says the drop is as large as possible. Geometrically: taking $P \neq Q$, condition $(\star)$ says there is a function in $L(D)$ vanishing at $P$ but not at $Q$, i.e.\ $\varphi_D$ separates $P$ and $Q$; taking $P = Q$, it says $\varphi_D$ separates tangent directions at $P$, i.e.\ $\dee(\varphi_D)_P$ is injective.

\begin{theorem}\label{thm:embedding}
	$\varphi_D$ is an embedding if and only if $D$ satisfies $(\star)$.
\end{theorem}

We omit the proof, which is a careful unwinding of the two conditions in the previous paragraph.

\begin{corollary}\label{cor:embedbig}
	If $\deg D > 2g$ then $\varphi_D$ is an embedding.
\end{corollary}

\begin{proof}
	Both $\deg D$ and $\deg(D-[P]-[Q]) = \deg D - 2$ exceed $2g-2$, so Corollary~\ref{cor:RRlarge} applies to both, giving
	$$\ell(D) - \ell(D-[P]-[Q]) = \deg D - (\deg D - 2) = 2. \qedhere$$
\end{proof}

\begin{corollary}
	Every curve of genus $g$ embeds in $\PP^N$ for some $N$ depending only on $g$.
\end{corollary}

\begin{proof}
	If $g \geq 2$, take $D$ in the class $2K_C$, of degree $4g-4 > 2g$. If $g = 1$, take $D = 3[P_0]$ for any $P_0$, of degree $3 > 2$. If $g = 0$, $C \cong \PP^1$ already. In each case $\deg D > 2g$, so Corollary~\ref{cor:embedbig} applies, and $N = \ell(D)-1$ is determined by $g$ via Corollary~\ref{cor:RRlarge}.
\end{proof}

\begin{theorem}[Canonical embedding]
	Let $C$ be a non-hyperelliptic curve of genus $g \geq 3$. Then $\varphi_{K_C} \colon C \to \PP^{g-1}$ is an embedding.
\end{theorem}

\begin{proof}
	Note $\ell(K) = g \geq 3 \geq 2$, so $\varphi_K$ is defined and maps to $\PP^{g-1}$. Suppose $K$ fails $(\star)$, so there are $P,Q$ with $\ell(K - [P]-[Q]) \geq g - 1$. Apply Riemann--Roch to $D = [P]+[Q]$:
	$$\ell(D) - \ell(K-D) = \deg D - g + 1 = 3 - g,$$
	so $\ell(D) \geq (g-1) + 3 - g = 2$. But $\ell(D) = 2$ with $\deg D = 2$ says exactly that $C$ is hyperelliptic, contradiction. (And $\ell(D)\leq \deg D + 1 = 3$; if $\ell(D) = 3$ then by Corollary~\ref{cor:g0}-type reasoning $C\cong\PP^1$, excluded since $g \geq 3$.)
\end{proof}

So every non-hyperelliptic curve has a canonical model in $\PP^{g-1}$, determined intrinsically with no choices. For $g = 3$ this is a plane quartic, matching the genus--degree formula.

\subsection{Genus Zero: Rational Curves}

We have already done this case, but let us collect the statements.

\begin{theorem}
	For a smooth projective irreducible curve $C$, the following are equivalent:
	\begin{enumerate}[label=(\roman*)]
		\item $g(C) = 0$;
		\item $C \cong \PP^1$;
		\item $\C(C) \cong \C(t)$;
		\item $C$ has a divisor of degree $1$ that is linearly equivalent to an effective divisor with $\ell = 2$;
		\item $\Pic^0(C) = 0$.
	\end{enumerate}
\end{theorem}

\begin{proof}
	(i)$\Rightarrow$(ii) is Corollary~\ref{cor:g0}. (ii)$\Rightarrow$(iii) is clear, and (iii)$\Rightarrow$(ii) is Theorem~\ref{thm:biratfield} together with Corollary~\ref{cor:ratismorph}. (ii)$\Rightarrow$(v) is Proposition~\ref{prop:picP1}. For (v)$\Rightarrow$(i), if $g \geq 1$ then for distinct $P,Q$ we have $\ell([P]) = 1$ by Corollary~\ref{cor:RRlarge} if $g = 1$ (or Proposition~\ref{prop:ellbounds} in general), so $[P]\not\sim[Q]$ and $[P]-[Q]$ is a nonzero element of $\Pic^0(C)$. The equivalence with (iv) is the computation in Corollary~\ref{cor:g0}.
\end{proof}

\begin{example}
	The conic $C = \V(X_0X_2-X_1^2)$ has genus $0$, so is isomorphic to $\PP^1$; we checked this directly in Proposition~\ref{prop:conicisP1}, and the map $\PP^1\to C$ is $\varphi_D$ for $D = 2[\infty]$, since $L(2[\infty])$ has basis $1,t,t^2$.
\end{example}

\begin{example}
	The nodal cubic $\V(Y^2-X^2(X+1))$ is \emph{not} smooth, so our theory does not directly apply. But it is birational to $\PP^1$: parametrise by the slope $t = Y/X$ of the line through the node, giving $t^2 = X+1$, so $X = t^2-1$ and $Y = t(t^2-1)$. This is a rational parametrisation, and the two values $t = \pm1$ both map to the node. So the smooth projective curve birational to it is $\PP^1$, and the node is what you get by gluing two points of $\PP^1$ together. The same works for the cusp, with $t = Y/X$ giving $(t^2,t^3)$.
\end{example}

\subsection{Genus One: Elliptic Curves}

\begin{definition}[Elliptic Curve]
	An \vocab{elliptic curve} is a pair $E = (C,P_0)$ where $C$ is a smooth projective irreducible curve of genus $1$ and $P_0 \in C$.
\end{definition}

The choice of base point is essential: in genus $1$ there is no canonical point, and different choices give isomorphic curves but different groups structures (well, isomorphic ones, but the isomorphism is a translation).

Genus $1$ is the first case where the interesting range $0\leq \deg D \leq 2g-2$ is nonempty -- but only just, since $2g-2 = 0$. This makes the arithmetic particularly clean.

\begin{corollary}\label{cor:g1RR}
	Let $C$ have genus $1$ and let $D$ be a divisor with $\deg D > 0$. Then $\ell(D) = \deg D$.
\end{corollary}

\begin{proof}
	$K_C$ has degree $2g-2 = 0$, so $\deg(K-D) < 0$ and $\ell(K-D) = 0$. Riemann--Roch gives $\ell(D) = \deg D - 1 + 1 = \deg D$.
\end{proof}

Now here is the construction of the group law, done intrinsically.

\begin{theorem}\label{thm:elliptic}
	Let $E = (C,P_0)$ be an elliptic curve. Then the map
	$$\beta \colon C \to \Pic^0(C), \qquad P \mapsto [\,[P] - [P_0]\,]$$
	is a bijection. Transporting the group structure of $\Pic^0(C)$ back along $\beta$ makes $C$ into an abelian group with identity $P_0$.
\end{theorem}

\begin{proof}
	\emph{Injectivity.} If $\beta(P) = \beta(Q)$ then $[P]\sim[Q]$, so by Proposition~\ref{prop:Lsim}, $\ell([P]) = \ell([Q])$ and there is $f \in \C(C)^\times$ with $\Divv(f) = [P]-[Q]$. If $P \neq Q$ then $f$ has a single simple pole, so $\varphi = (1:f) \colon C \to \PP^1$ has degree $1$ and is therefore an isomorphism, forcing $g(C) = 0$, a contradiction. Hence $P = Q$.

	\emph{Surjectivity.} Let $D \in \Div^0(C)$. Then $\deg(D + [P_0]) = 1 > 0$, so $\ell(D+[P_0]) = 1$ by Corollary~\ref{cor:g1RR}. Pick $0 \neq f \in L(D+[P_0])$; then $\Divv(f) + D + [P_0]$ is effective of degree $1$, hence equals $[P]$ for a unique point $P$. So $D + [P_0]\sim [P]$, i.e.\ $D \sim [P]-[P_0]$ and $D = \beta(P)$ in $\Pic^0$.
\end{proof}

\begin{definition}
	For $P,Q \in E$ define $P +_E Q$ to be the unique point $R$ with $\beta(R) = \beta(P)+\beta(Q)$, i.e.\ with
	$$[P] + [Q] - [P_0] \sim [R].$$
\end{definition}

\begin{theorem}
	Let $E = (C,P_0)$ be an elliptic curve realised as a smooth plane cubic with $P_0$ a flex. Then for $P,Q,R \in C$,
	$$P +_E Q +_E R = P_0 \iff P, Q, R \text{ are collinear.}$$
	Consequently $+_E$ agrees with the chord-and-tangent law of Definition~\ref{def:chordtangent}, and in particular that law is associative.
\end{theorem}

\begin{proof}
	Let $L$ be a line, with $\Divv(L) = [P]+[Q]+[R]$ its divisor on $C$ (a degree $3$ effective divisor by B\'ezout). Let $L_0$ be the tangent line to $C$ at the flex $P_0$, so $\Divv(L_0) = 3[P_0]$. By Proposition~\ref{prop:hyperplanediv},
	$$\Divv(L) - \Divv(L_0) = \Divv(L/L_0)$$
	is principal, so $[P]+[Q]+[R] \sim 3[P_0]$, that is,
	$$\beta(P)+\beta(Q)+\beta(R) = \big[[P]+[Q]+[R] - 3[P_0]\big] = 0.$$
	Hence $P+_EQ+_ER = P_0$.

	Conversely, suppose $P +_E Q +_E R = P_0$, and let $L$ be the line through $P$ and $Q$, meeting $C$ again at $R'$. By the above, $P+_EQ+_ER' = P_0$, so $R = R'$ and $P,Q,R$ are collinear.

	Finally, the chord-and-tangent law sets $P + Q = P_0 * (P*Q)$; writing $S = P*Q$, the points $P,Q,S$ are collinear so $P+_EQ+_ES = P_0$, and $P_0, S, P_0*S$ are collinear so $S +_E (P_0*S) = P_0$ (using $P_0 +_E X = X$). Hence $P_0 * S = -_E S = P +_E Q$.
\end{proof}

This completes the proof of Theorem~\ref{thm:grouplaw}: associativity of the chord-and-tangent law is inherited from associativity of addition in $\Pic^0(C)$, which is free.

\begin{example}
	On $E \colon y^2 = x^3-x+1$ with $P_0$ the point at infinity, take $P = (-1,-1)$ and $Q = (1.6, 1.8698\ldots)$. The line through them has slope $m = 1.1038\ldots$, and substituting into the cubic gives $x^3 - m^2x^2 - (\ldots)x - (\ldots) = 0$ whose roots sum to $m^2$. Hence
	$$x_{P*Q} = m^2 - x_P - x_Q = 0.6183\ldots,$$
	and $P +_E Q = (0.6183\ldots, -0.7862\ldots)$. This is the first figure of Section~6.6. Similarly the tangent at $T = (1,1)$ has slope $(3\cdot 1 - 1)/2 = 1$, so $x_{2T} = 1 - 2 = -1$ and $2T = (-1,1)$.
\end{example}

\begin{remark}
	Over $\C$, a genus $1$ curve is topologically a torus $S^1\times S^1$, and in fact $C \cong \C/\Lambda$ for a lattice $\Lambda \leq \C$, as a complex manifold and as a group. The group law $+_E$ corresponds to addition in $\C/\Lambda$, and the Weierstrass form comes from the Weierstrass $\wp$-function, which satisfies
	$$(\wp')^2 = 4\wp^3 - g_2\wp - g_3.$$
	Over $\Q$, the group $E(\Q)$ of rational points is finitely generated (Mordell's theorem), and this is where the subject meets number theory.
\end{remark}

\subsection{Genus Two and Beyond}

\begin{proposition}
	Every curve of genus $2$ is hyperelliptic.
\end{proposition}

\begin{proof}
	Take $D = K_C$, of degree $2g-2 = 2$, with $\ell(D) = g = 2$. This is precisely the criterion for hyperellipticity noted after Example~\ref{ex:hyperelliptic}: $\varphi_{K}$ is a morphism $C \to \PP^1$, and by Theorem~\ref{thm:embedding} it is not an embedding (it maps a genus $2$ curve to a genus $0$ one), so it has degree $2$.
\end{proof}

\begin{example}
	Genus $2$ curves are exactly the smooth projective models of $y^2 = p(x)$ with $\deg p = 6$ (or $5$) and $p$ squarefree, by Example~\ref{ex:hyperelliptic} with $m = 3$. Up to isomorphism they are parametrised by the six branch points modulo $\PGL_2(\C)$, giving a $3$-dimensional family. In general, curves of genus $g \geq 2$ form a $(3g-3)$-dimensional family -- the \vocab{moduli space} $\mathcal M_g$ -- which one can already guess from this kind of count.
\end{example}

\begin{example}
	Let us determine all the curves of genus $3$. If $C$ is non-hyperelliptic, the canonical map embeds $C$ in $\PP^{g-1} = \PP^2$ as a curve of degree $\deg K = 2g-2 = 4$, i.e.\ a smooth plane quartic; and conversely every smooth plane quartic has genus $3$ by the genus--degree formula. If $C$ is hyperelliptic, it is the smooth model of $y^2 = p(x)$ with $\deg p = 8$. So genus $3$ curves are either plane quartics or hyperelliptic.
\end{example}

We close by recording the classification of maps in the simplest cases, which is a good test of everything we have set up.

\begin{proposition}
	Let $C$ be a curve of genus $g$.
	\begin{enumerate}[label=(\roman*)]
		\item If $g = 0$, the morphisms $C \to \PP^1$ of degree $d$ are exactly $\varphi_D$ for $D$ of degree $d$, and there is a $(2d+1)$-dimensional family of them.
		\item If $g \geq 1$, there is no morphism $C \to \PP^1$ of degree $1$.
		\item If $g\geq 2$, every morphism $C\to C$ of degree $1$ is an automorphism, and the group of these is finite (of order at most $84(g-1)$, by Hurwitz's automorphism theorem).
	\end{enumerate}
\end{proposition}

\begin{proof}[Proof of (ii)]
	A degree $1$ morphism induces an isomorphism of function fields, hence is an isomorphism of curves, so $C\cong\PP^1$ and $g = 0$.
\end{proof}

\begin{remark}
	Hurwitz's bound in (iii) is a beautiful application of Riemann--Hurwitz. If $G = \mathrm{Aut}(C)$ is finite of order $N$, the quotient $C/G$ is a curve $D$ of some genus $h$, and $C \to D$ has degree $N$; analysing the possible ramification data with Riemann--Hurwitz shows that
	$$2g-2 = N\left(2h - 2 + \sum_{i}\left(1 - \frac{1}{r_i}\right)\right)$$
	for the ramification orders $r_i$ over the branch points, and the smallest positive value the bracket can take is $1/42$, attained when $h = 0$ and $(r_1,r_2,r_3) = (2,3,7)$. Hence $N \leq 84(g-1)$.
\end{remark}

\subsection{A Summary}

It is worth stepping back. We began with the dictionary
$$\{\text{affine varieties}\} \longleftrightarrow \{\text{reduced finitely generated } \C\text{-algebras}\},$$
refined it to
$$\{\text{irreducible varieties up to birational equivalence}\} \longleftrightarrow \{\text{finitely generated field extensions of } \C\},$$
and then specialised to transcendence degree $1$, where the classification becomes tractable. For curves we found:
\begin{itemize}
	\item every function field of transcendence degree $1$ is the function field of a \emph{unique} smooth projective curve, since birational smooth projective curves are isomorphic (Corollary~\ref{cor:ratismorph});
	\item each such curve carries an intrinsic invariant, the genus $g = \ell(K_C)$, computable from any plane model by the genus--degree formula, and from any map to $\PP^1$ by Riemann--Hurwitz;
	\item the finite-dimensional spaces $L(D)$ are computed by Riemann--Roch, and they in turn produce all the maps from $C$ to projective space;
	\item genus $0$ gives exactly $\PP^1$; genus $1$ gives the elliptic curves, whose points form a group; and genus $\geq 2$ gives rigid objects with finite automorphism groups.
\end{itemize}
The natural next steps -- sheaves and their cohomology, which make Riemann--Roch a formal consequence of a general index theorem, and schemes, which restore the nilpotents we threw away in Section~2 -- are the subject of Part III.

%%%END%%%
\end{document}
